% !TEX root = ../../tree_notes.tex
% =====================================================================
%  Chapter 18 -- Fenwick Trees   WRITTEN
%  (folder ch14_fenwick_trees; label prefix c14; prints as Chapter 18;
%  folder number 14; build with:  python tools/build.py --only 14)
%
%  Owns the ARRAY-INDEX engine that Ch.5 and Ch.16 kept handing ranges to:
%    * 18.1 the BIT itself -- lowbit i&(-i), responsibility intervals
%      [i-lowbit(i)+1, i], the update/query walks, the O(n) linear build,
%      and why the array hides a binary tree (solved CSES 1648).
%    * 18.2 range variants -- difference BIT for range-add/point-query
%      (solved CSES 1651), the two-BIT identity for range-add/range-sum
%      (worked, since CSES has no such task), coordinate compression,
%      prefix-min with the append-only caveat.
%    * 18.3 frequency BIT and ORDER STATISTICS -- k-th by binary lifting
%      on the tree, inversion-style sweeps, decreasing-triplet counts
%      (solved CF 61E); k-th smallest in a static range via value-sweep.
%    * 18.4 Fenwick ON TREES -- Ch.5's tin/tout flat array driven by a BIT:
%      subtree add / subtree sum, the two-event enter +1 / leave -1 trick
%      for root-path reads (solved CF 383C Propagating tree).
%    * 18.5 offline sweeps -- sort one axis, BIT the other: 2D dominance
%      counting (worked) and parallel binary search whose predicate is a
%      BIT instead of Ch.17's rollback DSU (worked), discharging the
%      forward obligation Ch.17 pinned on this chapter.
%  Diagram-first: every responsibility-interval table, cascade walk,
%  difference cascade, lifting search and Euler overlay is hand-computed
%  on TWO running examples -- the 8-slot array a = (3,2,4,5,1,1,5,3)
%  (CSES 1648's own sample array) extended to 16 slots A, whose BIT is
%  T = (3,5,4,14,1,2,5,24,7,7,2,13,1,7,3,49), and the 10-node
%  tree T of Ch.12-17 (edges 1-2,1-3,2-4,2-5,3-6,3-7,5-8,7-9,9-10) with
%  the Ch.16 preorder tin = (1,2,6,3,4,7,8,5,9,10) (children ascending)
%  and INCLUSIVE tout = (10,5,10,3,5,7,10,5,10,10); every official sample
%  was live-verified against the problem before a trace was written.
%
%  Anti-repetition contract (teach once, cross-reference by label):
%    * Ch.5 (\chref{c24}) owns Euler-tour FLATTENING (tin/tout, subtree=
%      interval, \thmref{c24-range}); 18.4 reuses \tin/\tout/\size by label
%      and never re-derives them -- it owns the RANGE ENGINE (Ch.5's own
%      forward pointer \chref{c14} is discharged here).
%    * Ch.19/20 (\chref{c13}/\chref{c15}) own segment trees and lazy
%      propagation; the BIT-vs-segtree boundary is prose + forward-cite,
%      no segtree code appears in this chapter.
%    * Ch.1 (\chref{c01}) owns merge-sort inversion counting; 18.3 owns
%      the ONLINE frequency-BIT count and states the distinction.
%    * Ch.6 (\chref{c06}) owns CF 1354D Multiset, PBDS order-statistics
%      trees, and LC 307 (Range Sum Query - Mutable) with its URLs; 18.3
%      is the array-index RIVAL, cited by prose only, never re-linked.
%    * Ch.5's summit owns LC 315 Count of Smaller Numbers After Self
%      (size-augmented BST) and Ch.6's mixed list owns LC 493 Reverse
%      Pairs -- both URLs are taken; 18.3 solves CF 61E instead and
%      mentions 315/493 in prose as the same-sweep relatives.
%    * Ch.8 (\chref{c07}) heaps own the STREAM median two-heap trick;
%      18.3 handles k-th/offline only and back-cites.
%    * Ch.17 (\chref{c28}) owns PBS-with-rollback-DSU; 18.5 reuses the
%      lockstep SCHEDULING idea by reference (\secref{c28-offline}) and
%      swaps the predicate engine to a BIT.
%
%  Roster (every title/rating/sample live-verified before writing;
%  Codeforces-forward; URLs urlsweep-gated at write time):
%    solved: CSES 1648   Dynamic Range Sum Queries  (18.1, point add + range sum)
%            CSES 1651   Range Update Queries       (18.2, range add + point read)
%            CF  61E     Enemy is weak              (18.3, 1900, decreasing triplets)
%            CF  383C    Propagating tree           (18.4, 2000, Euler + parity BIT)
%    worked: two-BIT range-add/range-sum (18.2); kth-lifting + Hotel Queries
%            CSES 1143 (18.3); dominance counting + PBS-over-BIT (18.5)
%    summit: CF 1181D Irrigation (2200; offline + BIT + kth-lifting;
%            both official samples hand-replayed)
%    practice + mixed (10): all fresh, urlsweep-gated, 1600--2400 ladder
%    NOTE: the plan's original picks CSES 2166 (max-prefix, a segtree
%    problem), CF 915G (Easy Matrices -- wrong problem), CF 343D (URL used
%    by Ch.6 splay), LC 315/493/307 (URLs used by Ch.5/6) were all
%    replaced at the verification gate by the fresh equivalents above.
% =====================================================================
\bookchapter{Fenwick Trees}{c14}
\chapterinfo{CF 1600--2400, CSES, LC Hard}{Written}

% =====================================================================
%  Chapter 18 opener -- "Chapter at a Glance", three reading paths + the
%  ladder staircase, THE two running examples (the 16-slot array A with
%  the CSES 1648 sample in its first 8 slots, and the 10-node tree T of
%  Ch.12--17 with its Ch.16 preorder tins), the BIT-flavour cost chart,
%  the "who owns what" comparison (Ch.5 owns flattening; Ch.19/20 own
%  segment trees; Ch.1 owns merge counting; Ch.5/6 own LC 315/493/307 and
%  PBDS; Ch.17 owns PBS-with-rollback; THIS chapter owns the array-index
%  engine), recognition guide, constraint->tool table, notation.
%  Exports fig:c14-array and fig:c14-taxonomy.
%  Hand-computed values the whole chapter reuses (verified once here):
%    A=(3,2,4,5,1,1,5,3,7,0,2,4,1,6,3,2), sum 49;
%    T=(3,5,4,14,1,2,5,24,7,7,2,13,1,7,3,49);   (post-build block sums)
%    query(6)=T6+T4=16, query(7)=T7+T6+T4=21, query(13)=T13+T12+T8=38;
%    add 3 to A[3] -> T4=17,T8=27,T16=52 (query(7) then 24).
%  Every overlay below is hand-computed; every later section re-uses
%  exactly these numbers -- if a section disagrees, this opener is wrong
%  and the section is right; fix it here, never the other way.
% =====================================================================

\begin{chaptertheorybox}{Chapter at a Glance}
Chapter~5 ended every flattened tree with the sentence ``and now answer this
\emph{range} with a Fenwick tree, covered in \chref{c14}''.  Chapter~16 handed
its depth buckets to the same pointer.  This chapter is that structure: a
\term{Fenwick tree} (Binary Indexed Tree, BIT) stores, for every index \(i\),
one partial sum over the block \([i-\operatorname{lowbit}(i)+1,\,i]\) -- and
the single bit operation \(i\;\&\;(-i)\) turns both \emph{prefix queries} and
\emph{point updates} into \(\bigO{\log n}\) walks of at most \(\log n\) steps
in an ordinary array.  No nodes, no pointers, no rebalancing: the index
\emph{is} the tree.

\textbf{You will learn to}
\begin{itemize}[nosep]
  \item read responsibility intervals off the binary expansion, run the
        \code{update}/\code{query} walks, and build a BIT in \(\bigO{n}\);
        \hfill\lv{2}--\lv{3}
  \item answer point-update + range-sum end to end (CSES~1648, \(\secref{c14-core}\)); \hfill\lv{3}
  \item switch the BIT to \emph{differences} for range-add/point-read
        (CSES~1651) and pair two BITs for range-add/range-sum (\(\secref{c14-range}\)); \hfill\lv{4}
  \item use a BIT as a \emph{frequency} table: k-th smallest by binary
        lifting, inversion-style sweeps, decreasing triplets (CF~61E,
        \(\secref{c14-order}\)); \hfill\lv{4}--\lv{5}
  \item drive Ch.5's Euler array with a BIT -- subtree add/sum and the
        enter\(+1\)/leave\(-1\) root-path trick (CF~383C, \(\secref{c14-trees}\)); \hfill\lv{5}
  \item take queries \emph{offline}: sort one axis, sweep it, BIT the other
        (\(\secref{c14-offline}\)) -- dominance counting and parallel binary
        search whose predicate is a BIT instead of Ch.17's rollback DSU; and
        compose everything on the Summit CF~1181D. \hfill\lv{5}--\lv{6}
\end{itemize}

\textbf{Prerequisites.} Chapter~5 (\chref{c24}) for \(\tin/\tout\) and
``a subtree is one interval'' (\thmref{c24-range}) -- \secref{c14-trees}
consumes that reduction and never re-derives it.  Chapter~16 (\chref{c11})
for the running tree \(T\), its depths and its preorder stamps.  Chapter~17
(\chref{c28}) only for the \emph{scheduling} idea of parallel binary search,
reused by reference in \secref{c14-offline}.  Arrays, binary representation,
and \(O(\log n)\) as a concept are all that is assumed.
\end{chaptertheorybox}

% ---------------- Reading paths + ladder position ----------------
\begin{chaptertheorybox}{How to Read This Chapter (three paths)}
Level badges \(\lv1\)--\(\lv6\): \(\lv2\) is hand-feeding the lowbit walk on
a drawn array and reading intervals off the table; \(\lv3\) is a working
point-update/range-sum BIT on a real input (CSES~1648 is exactly this);
\(\lv4\) is choosing the right encoding -- differences, two BITs, or a
frequency table -- \emph{before} writing code; \(\lv5\) is composing two
chapters' machinery (Euler array \(+\) BIT; sweep \(+\) lifting) on problems
like CF~383C; \(\lv6\) is recognising that a ``static 2D query'' is really a
sorted offline sweep, or that a threshold question is PBS with a BIT
predicate -- the Summit's league.

\textbf{Ground-up path (\(\lv2\)--\(\lv4\), ``I have never written a BIT'').}
\secref{c14-core} in order: intervals, the two walks, the linear build,
CSES~1648 with every step traced on \(A\).  Then \secref{c14-range}: the
difference BIT is one change of perspective (\(\text{read}(i)=\) prefix of
differences) and the two-BIT identity is its algebraic cousin.

\textbf{Query-hammer path (\(\lv4\)--\(\lv5\), ``my TLE is a counting
problem'').}  \secref{c14-core} for the mechanics only, then
\secref{c14-order}: re-read the BIT as a table of \emph{frequencies}, where
\code{query}(k) counts elements \(\le k\).  K-th by lifting, inversion
sweeps, CF~61E.  If you only need one tool from the chapter, this section's
lifting search is it.

\textbf{Tree/offline path (\(\lv5\)--\(\lv6\), ``ranges on trees and
2D queries'').}  Ch.5's \secref{c24-flatten} (already read, by contract),
then \secref{c14-trees}: the flat array \emph{plus} this chapter's engine,
CF~383C fully traced on \(T\).  Finish with \secref{c14-offline} --
dominance sweeps and PBS-over-BIT -- and the Summit CF~1181D, which needs
all of it: offline sorting, a frequency BIT, and the lifting search
answering ``which city just crossed its quota?''.
\end{chaptertheorybox}

\begin{figure}[H]
\centering
\begin{tikzpicture}[diagram, x=9mm, y=7mm]
  \foreach \st/\lab in {1/lv1,2/lv2,3/lv3,4/lv4,5/lv5,6/lv6} {
    \fill[\lab!22, draw=\lab, line width=0.8pt] (\st-1,0) rectangle (\st,\st*0.52);
    \node[font=\scriptsize\bfseries, text=\lab] at (\st-0.5,\st*0.52-0.42) {L\st};}
  \node[annot, anchor=south west, font=\scriptsize]
    at (0.05,0.62) {18.1 lowbit + responsibility intervals (read the table)};
  \node[annot, anchor=south west, font=\scriptsize]
    at (1.05,1.14) {18.1 walks + O(n) build; solved CSES 1648};
  \node[annot, anchor=south west, font=\scriptsize]
    at (2.05,1.66) {18.2 difference BIT (CSES 1651) + two-BIT range-sum};
  \node[annot, anchor=south west, font=\scriptsize]
    at (3.05,2.18) {18.3 frequency BIT: kth-by-lifting, sweeps (CF 61E)};
  \node[annot, anchor=south west, font=\scriptsize, text=lv5]
    at (4.05,2.70) {18.4 BIT on Ch.5's Euler array: subtree/path (CF 383C)};
  \node[annot, anchor=south west, font=\scriptsize, text=lv6]
    at (5.05,3.22) {18.5 offline sweeps + PBS-over-BIT + Summit CF 1181D};
\end{tikzpicture}
\caption*{Chapter~18 on the six-step ladder: intervals and walks are
\(\lv2\)--\(\lv3\); choosing the encoding (differences, two BITs,
frequencies) is \(\lv4\); composing with Ch.5's flattening is \(\lv5\);
offline scheduling and the Summit CF~1181D are step~6.  Ground-up path
\(=\) steps 2--4; full contest path runs to~6.}
\end{figure}

% ---------------- The running examples ----------------
\subsection*{The Running Examples: One Array, One Tree}

\emph{Two} objects carry the whole chapter.  The first is a 16-slot array
whose opening half is CSES~1648's own sample, extended so that every
lowbit power up to \(16\) gets used:
\[
  A=(\,3,2,4,5,1,1,5,3,\;7,0,2,4,1,6,3,2\,),\qquad \textstyle\sum A = 49 ,
\]
with the prefix sums you will keep needing
(\(\mathit{pref}[k]=A_1+\cdots+A_k\)):
\[
  \mathit{pref}=(3,5,9,14,15,16,21,24,\;31,31,33,37,38,44,47,49).
\]
Slots \(9..16\) only ever exist to
show the \(\text{lowbit}=8\) and \(16\) intervals; every solved sample in the
chapter lives in slots \(1..8\).  The second is the tree \(T\) Chapters~12--17
have been sharing -- nodes \(1..10\), root~1, edges
\[
  1\!-\!2,\; 1\!-\!3,\; 2\!-\!4,\; 2\!-\!5,\; 3\!-\!6,\; 3\!-\!7,\; 5\!-\!8,\; 7\!-\!9,\; 9\!-\!10 .
\]
\[
  \depth=(0,1,1,2,2,2,2,3,3,4),\quad \size=(10,4,5,1,2,1,3,1,2,1),
\]
and the Ch.16 preorder (\secref{c11-paths}'s convention, \(\tout\) \emph{inclusive}):
\[
  \tin=(1,2,6,3,4,7,8,5,9,10),\qquad
  \tout=(10,5,10,3,5,7,10,5,10,10).
\]
So \(\subtree(2)=[2,5]\) holds nodes \(\{2,4,5,8\}\) and
\(\subtree(7)=[8,10]\) holds \(\{7,9,10\}\)
-- the intervals \secref{c14-trees} hands to the BIT.
\figref{c14-array} stamps \(A\) with its responsibility table and the two
walks every section reuses.

\begin{figure}[H]
\centering
\begin{tikzpicture}[diagram, scale=0.62, transform shape]
  % A row
  \foreach \i/\vv in {1/3,2/2,3/4,4/5,5/1,6/1,7/5,8/3,9/7,10/0,11/2,12/4,13/1,14/6,15/3,16/2} {
    \draw[primaryblue, line width=0.6pt, fill=primaryblue!6] (\i*0.9,0) rectangle (\i*0.9+0.9,0.9);
    \node[font=\scriptsize\ttfamily] at (\i*0.9+0.45,0.45) {\vv};
    \node[font=\tiny\ttfamily, text=darkgray] at (\i*0.9+0.45,1.15) {A[\i]};}
  % T row
  \foreach \i/\tv in {1/3,2/5,3/4,4/14,5/1,6/2,7/5,8/24,9/7,10/7,11/2,12/13,13/1,14/7,15/3,16/49} {
    \draw[difforange, line width=0.6pt, fill=difforange!8] (\i*0.9,-1.5) rectangle (\i*0.9+0.9,-0.6);
    \node[font=\scriptsize\ttfamily, text=orange!60!black] at (\i*0.9+0.45,-1.05) {\tv};
    \node[font=\tiny\ttfamily, text=darkgray] at (\i*0.9+0.45,-1.75) {T[\i]};}
  % responsibility interval bars [i-lowbit+1, i] (drawn implicitly by width below)
  \foreach \i/\lo in {1/1,2/1,3/3,4/1,5/5,6/5,7/7,8/1,9/9,10/9,11/11,12/9,13/13,14/13,15/15,16/1} {
    \fill[successgreen!30] (\lo*0.9+0.12,-2.55) rectangle (\i*0.9+0.78,-2.3);
    \node[font=\tiny, text=green!45!black] at (\i*0.9+0.45,-2.42) {[\lo,\i]};}
  \node[font=\scriptsize\ttfamily, text=primaryblue!60!black, anchor=east] at (0,0.45) {A:};
  \node[font=\scriptsize\ttfamily, text=orange!60!black, anchor=east] at (0,-1.05) {T:};
  \node[font=\scriptsize, text=green!45!black, anchor=east] at (0,-2.42) {covers:};
  % query walk 6 -> 4 -> 0
  \draw[warningred, line width=1.1pt, -latex] (6*0.9+0.45,-2.05) -- (4*0.9+0.45,-2.05);
  \draw[warningred, line width=1.1pt, -latex] (4*0.9+0.45,-2.05) -- (0*0.9+0.45,-2.05);
  \node[font=\scriptsize\ttfamily, text=warningred, anchor=north west] at (6.9,-2.1)
    {\(\text{query}(6)=T[6]+T[4]=2+14=16=\mathit{pref}[6]\)};
  % interval table
  \node[font=\scriptsize, text=darkgray, anchor=north west] at (0.3,-3.15) {%
    \(
    \begin{array}{c|cccccccccccccccc}
    i&1&2&3&4&5&6&7&8&9&10&11&12&13&14&15&16\\\hline
    \operatorname{lowbit}(i)&1&2&1&4&1&2&1&8&1&2&1&4&1&2&1&16\\
    \text{block of }T[i]&[1,1]&[1,2]&[3,3]&[1,4]&[5,5]&[5,6]&[7,7]&[1,8]&[9,9]&[9,10]&[11,11]&[9,12]&[13,13]&[13,14]&[15,15]&[1,16]
    \end{array}
    \)};
\end{tikzpicture}
\caption{The running array \(A\) (top), its BIT \(T\) (middle, \(T[i]=\) sum
of \(A\) over the green block below it), and the responsibility table
(bottom).  Build checks reused everywhere: \(T[4]=3{+}2{+}4{+}5{=}14\),
\(T[2]=3{+}2{=}5\),
\(T[8]=14{+}1{+}1{+}5{+}3{=}24\) (CSES~1648's first sample answer),
\(T[12]=7{+}0{+}2{+}4{=}13\), \(T[14]=A[13]{+}A[14]=1{+}6{=}7\), and \(T[16]=\sum A=49\).  The red
arrows show one \emph{query} walk: \(6\to4\to0\) drops the lowest set bit
each step and visits \(T[6],T[4]\) covering \([5,6]\uplus[1,4]=[1,6]\) --
disjoint blocks, no double counting.  The mirror walk (update \(A_3{+}3\))
climbs \(3\to4\to8\to16\), touching every interval that contains slot~3; both
walks are \secref{c14-core}'s opening trace.}
\label{fig:c14-array}
\end{figure}

Six angles, six hand-computed overlays -- each lives in the section named on
its row and never changes:

\begin{center}
\small
\renewcommand{\arraystretch}{1.25}
\begin{tabular}{@{}p{0.24\linewidth} p{0.50\linewidth} l@{}}
\toprule
\textbf{Angle} & \textbf{What comes out} & \textbf{Traced in} \\
\midrule
BIT over \(A\) & \(\text{query}(13){=}1{+}13{+}24{=}38{=}\mathit{pref}[13]\); add~3 to \(A_3\) cascades \(T[4],T[8],T[16]\) & \secref{c14-core} \\
difference BIT & CSES 1651's sample: \(+1\) on \([2,5]\) turns read\((4)\) from \(5\) to \(6\); read\((2)\to3\), read\((7)=5\) (outside) & \secref{c14-range} \\
two BITs & \(a{=}(3,2,4,5,1,1,5,3)\): \(\text{rangeSum}(1,4)=14\); after \(+1\) on \([2,5]\), \(\text{rangeSum}(1,4)=17\) & \secref{c14-range} \\
frequency lifting & hotels-as-counts \((3,2,4,1,5,5,2,6)\), prefix \((3,5,9,10,15,\dots)\): first prefix \(\ge6\) is slot~3; \(\ge11\) is slot~5 & \secref{c14-order} \\
sweep on \((10,8,3,1)\) & smaller-after counts \((3,2,1,0)\); inversions \(=6\); decreasing triplets \(=4\) & \secref{c14-order} \\
BIT on \(T\)'s Euler array & subtree(2) \(=[2,5]\) add \(+3\), whole tree \(+2\): node~4 (inside) reads \(5\), node~6 (outside) reads \(2\) & \secref{c14-trees} \\
\bottomrule
\end{tabular}
\end{center}

% ---------------- Which flavour, cost chart ----------------
\subsection*{Which BIT Flavour? (the cost chart)}

All five are the \emph{same} array-of-block-sums; they differ in what you
\emph{store} in slot \(i\), and that choice is forced by the query shape, not
by taste.

\begin{center}
\small
\renewcommand{\arraystretch}{1.25}
\begin{tabular}{@{}p{0.19\linewidth} p{0.15\linewidth} p{0.15\linewidth} p{0.39\linewidth}@{}}
\toprule
\textbf{Flavour} & \textbf{Update} & \textbf{Query} & \textbf{When it is the right tool} \\
\midrule
plain (18.1) & point \(\bigO{\log}\) & prefix/range sum \(\bigO{\log}\) & values change one index at a time; sums asked \\
difference (18.2) & \textbf{range} add \(\bigO{\log}\) & point read \(\bigO{\log}\) & whole intervals shift; you only ever read single positions \\
two BITs (18.2) & range add \(\bigO{\log}\) & \textbf{range} sum \(\bigO{\log}\) & both intervals move \emph{and} interval sums are asked \\
frequency (18.3) & insert/delete at rank & count \(\le k\), k-th by lifting & elements are \emph{piles of equal items}, not positions; order questions \\
on a tree (18.4) & any flavour & over \([\tin(v),\tout(v)]\) & Ch.5 already flattened the tree; now the ranges get an engine \\
\bottomrule
\end{tabular}
\end{center}

\begin{insightbox}[One array, and the trick is always ``which block owns index \(i\)'']
Every flavour in the table is the same two loops over
\(i\;\&\;(-i)\).  What changes is the \emph{interpretation} of slot \(i\): a
sum of \(A\), a sum of differences, one component of a paired encoding, a
count of ranks, a subtree's weight.  When a BIT problem feels hard, you are
almost always choosing a row of this table, not inventing a data structure.
\end{insightbox}

% ---------------- Taxonomy chart ----------------
\begin{figure}[H]
\centering
\begin{tikzpicture}[diagram, scale=0.95, transform shape,
    mnode/.style={circle, draw=#1, line width=0.9pt, fill=#1!8, minimum size=7mm,
                  font=\scriptsize\bfseries, text=#1!60!black}]
  \node[mnode=primaryblue] (a) at (0,0)     {lowbit};
  \node[annot, font=\scriptsize, anchor=west] at (1.25,0) {i\&(-i) = the block T[i] owns: [i-lowbit+1, i] (18.1)};
  \node[mnode=successgreen] (b) at (0,-1)  {walk};
  \node[annot, font=\scriptsize, anchor=west] at (1.25,-1) {query drops the last set bit; update adds it -- both <= log n steps (18.1)};
  \node[mnode=difforange]   (c) at (0,-2)  {difference};
  \node[annot, font=\scriptsize, anchor=west] at (1.25,-2) {store D[i]=A[i]-A[i-1]: range add becomes two point updates (18.2)};
  \node[mnode=warningred]   (d) at (0,-3)  {two BITs};
  \node[annot, font=\scriptsize, anchor=west] at (1.25,-3) {range-add/range-sum: prefix(x) = x*S1(x) - S2(x) over paired trees (18.2)};
  \node[mnode=softpurple]   (e) at (0,-4)  {lift};
  \node[annot, font=\scriptsize, anchor=west] at (1.25,-4) {find the first prefix >= k by walking powers of two down the tree (18.3)};
  \node[mnode=deepblue]     (f) at (0,-5)  {sweep};
  \node[annot, font=\scriptsize, anchor=west] at (1.25,-5) {sort queries on one axis, insert-and-count along the other; offline 2D and PBS (18.5)};
\end{tikzpicture}
\caption{The chapter's six-state taxonomy, read top to bottom as an
escalation.  Rows 1--2 are the machine; row~3 is a change of \emph{what you
store}; row~4 is a change of \emph{algebra} (two trees, one identity); row~5
runs the tree \emph{backwards} (search instead of sum); row~6 is a change of
\emph{schedule} -- queries arrive sorted, not online.  Every solved problem
sits on exactly one row: CSES~1648 on rows~1--2, CSES~1651 on row~3, the
two-BIT worked example on row~4, CF~61E and CSES~1143 on row~5, CF~383C on
row~3-plus-Ch.5, and the Summit CF~1181D spanning rows~5--6.}
\label{fig:c14-taxonomy}
\end{figure}

% ---------------- Who owns what ----------------
\subsection*{Who Owns What (so nothing is learned twice)}

\begin{center}
\small
\renewcommand{\arraystretch}{1.25}
\begin{tabular}{@{}p{0.30\linewidth} p{0.32\linewidth} p{0.32\linewidth}@{}}
\toprule
\textbf{Idea} & \textbf{Owner} & \textbf{This chapter's share} \\
\midrule
Euler tour, \(\tin/\tout\), ``subtree \(=\) interval'' & Ch.~5 (\chref{c24}) & cited by label; 18.4 supplies the range \emph{engine} Ch.~5 promised \\
segment trees, lazy propagation, arbitrary range ops & Ch.~19/20 (\chref{c13}, \chref{c15}) & boundary is prose only; BIT wins when the op is a sum/count and code size matters \\
merge-sort inversion counting & Ch.~1 (\chref{c01}) & same bound, offline divide\&conquer there; \emph{online} frequency sweep here (\secref{c14-order}) \\
LC 315 / LC 493 / LC 307, PBDS order-stat trees & Ch.~5 summit, Ch.~6 (\chref{c06}) & URLs live there; 18.3 is the array-index rival, mentioned in prose, never re-linked \\
stream median (two heaps) & Ch.~8 (\chref{c07}) & online medians stay there; 18.3 does k-th \emph{offline}/static \\
PBS scheduling, rollback DSU & Ch.~17 (\chref{c28}) & the lockstep idea is cited; 18.5 swaps the predicate to a BIT \\
\bottomrule
\end{tabular}
\end{center}

% ---------------- Recognition guide ----------------
\begin{recognitionbox}[Pattern card -- ``prefix structure under point changes, or counting across a sorted axis'']
\reckey{``update one position / sum a range'' arriving interleaved; ``add to a
whole interval, read one position''; ``how many elements \(\le k\) / the
k-th smallest'' with deletions; ``count pairs/triplets with index order and
value order''; ``for every subtree, add/read, and queries are many'';
``static 2D orthogonal counts, offline allowed''; ``the least value making a
sweep-count reach \(k\)'' (then lift the BIT inside a parallel binary
search).}
\recuse{
  \item point updates + prefix/range sums -- plain BIT, \secref{c14-core} (CSES~1648);
  \item whole intervals shift, single positions read -- difference BIT,
        \secref{c14-range} (CSES~1651); shifts \emph{and} interval sums -- two BITs;
  \item the answer is a \emph{rank} question (how many / k-th) over a
        compressed value range -- frequency BIT + lifting, \secref{c14-order} (CF~61E);
  \item the ranges are subtrees or root-paths of a tree -- Ch.5 flattening +
        any BIT flavour, \secref{c14-trees} (CF~383C);
  \item queries have a \emph{free} second dimension you can sort on, or a
        threshold you can binary-search offline -- \secref{c14-offline} (Summit).
}
\recnotuse{
  \item the range operation is not a sum/count (min with \emph{arbitrary}
        writes, gcd, max after decrements) -- that is a segment tree, Ch.~19
        (\chref{c13}); see \secref{c14-range}'s prefix-min caveat for the one
        order-preserving case;
  \item the tree ranges must be cut, linked or rerolled dynamically --
        Euler-tour \emph{trees}/splay live in Ch.~6 (\chref{c06});
  \item updates and queries interleave \emph{online} with no offline
        permission and you need full range-assign -- segment tree with lazy,
        Ch.~20 (\chref{c15});
  \item the whole array is static -- one prefix array is \(O(1)\) per query
        (Ch.~5, \propref{c24-prefix}); a BIT there is pure overhead.
}
One-line test: \emph{is the operation ``add at an index'' and the question
``sum/count a prefix'' -- possibly after sorting, possibly on a flattened
tree?  Then this chapter's array-index engine is the smallest structure that
answers it.}
\end{recognitionbox}

% ---------------- Constraint -> tool ----------------
\subsection*{Reading the Constraints (which flavour the limits want)}

\begin{center}
\small
\renewcommand{\arraystretch}{1.25}
\begin{tabular}{@{}p{0.30\linewidth} p{0.40\linewidth} p{0.24\linewidth}@{}}
\toprule
\textbf{Constraint signature} & \textbf{What it rules in/out} & \textbf{Tool} \\
\midrule
\(n,q\le2\cdot10^5\), point updates, range sums & prefix array dead (updates); BIT is the min-cost AC & \secref{c14-core} \\
range adds + point reads, \(u\) can be negative & point values drift far -- keep \code{long long} & \secref{c14-range} \\
values to \(10^9\), queries are rank/counting & no array of size \(10^9\): compress, then frequency BIT & \secref{c14-order} \\
\(n\le10^6\), ``count triplets'' with a middle index & two sweeps and a product; \(O(n\log n)\) with small constants only & \secref{c14-order} (CF 61E) \\
tree, subtree updates/reads, \(n,q\le2\cdot10^5\) & DFS order turns it into row 2 on slots \(1..n\) & \secref{c14-trees} \\
\(n,m,q\le5\cdot10^5\), ``which city gets the \(k\)-th shipment'' & simulate = TLE; sort + BIT + lifting answers each in \(\log n\) & \secref{c14-offline}, Summit \\
\bottomrule
\end{tabular}
\end{center}

% ---------------- Notation ----------------
\subsection*{Notation Used Only Here}

\(\operatorname{lowbit}(i)=i\,\&\,(-i)\) -- the value of the lowest set bit of
\(i\), i.e.\ the largest power of two dividing \(i\); \(T[i]\) is the BIT slot
covering the \term{responsibility interval} \([i{-}\operatorname{lowbit}(i){+}1,\,i]\);
\(\code{query}(i)=A_1+\cdots+A_i\) (a prefix sum, \(\mathit{pref}[i]\)) and
\(\code{update}(i,\delta)\) adds \(\delta\) to \(A_i\) everywhere it is
buried.  A subtree interval sum uses Ch.~5's reduction with this chapter's
\(1\)-based, \(\tout\) \textbf{inclusive} convention:
\(\sum_{x\in\subtree(v)}=\code{query}(\tout(v))-\code{query}(\tin(v)-1)\)
(Ch.~5's \(0\)-based half-open stamps compute the same number -- never mix
the two in one program; the pitfall box of \secref{c14-trees} shows the
symptom).  A \term{frequency BIT} indexes \emph{values} (after
\term{coordinate compression}: replace values by their rank in the sorted
set) and stores multiplicities; \term{lifting} denotes the search that walks
a candidate index down from the highest power of two, building the largest
prefix with sum \(<k\).  \term{Offline} means all queries are known upfront
and re-ordered before replay; \term{PBS} (parallel binary search, Ch.~17's
\secref{c28-offline}) runs one shared sweep per bit level for \emph{all}
queries at once.  Every section runs on \(A\) and \(T\) above: if a table
here and a trace there disagree, this opener is wrong and the section is
right -- fix it here, never the other way.

% =====================================================================
%  18.1  The Binary Indexed Tree                     sec:c14-core
%  Exports: def:c14-interval  thm:c14-query  prop:c14-linear
%           lst:c14-query (lowbit + prefix walk)  lst:c14-update  lst:c14-build
%           lst:c14-1648-soln  fig:c14-hiddentree  prob:c14-rangesum
%  All numbers agree with 00_opener (fig:c14-array):
%    A=(3,2,4,5,1,1,5,3,7,0,2,4,1,6,3,2), sum 49;
%    T=(3,2,4,14,1,2,5,24,7,0,2,13,1,7,3,49);
%    query(6)=2+14=16, query(7)=5+2+14=21, query(13)=1+13+24=38;
%    update(3,+3): T4 14->17, T8 24->27, T16 49->52, then query(7)=24.
%  CSES 1648 sample (live-verified): n=8 q=4, a=(3,2,4,5,1,1,5,3),
%    queries 2 1 4 / 2 5 6 / 1 3 1 / 2 1 4  ->  outputs 14, 2, 11.
% =====================================================================
\section{The Binary Indexed Tree}\label{sec:c14-core}

\begin{previouslybox}[Ch.~5 promised this engine; Ch.~1 met its static twin]
Chapter~1 learned that a \emph{static} array answers every range sum with one
prefix table in \(O(1)\) (Ch.~5, \propref{c24-prefix}).  Chapter~5 then
flattened trees into arrays and repeatedly said ``with updates, hand the
range to a Fenwick tree, \chref{c14}'' -- without ever building one.  The
promise is kept now: this section derives the structure from the binary
representation of indices, in three loops totalling about ten lines.
\end{previouslybox}

\begin{topicintro}
A \term{Fenwick tree} or \term{Binary Indexed Tree} (BIT, BITr, or
\emph{Fenwick array} after P.\ Fenwick, 1994) is a single array \code{t[1..n]}
sitting \emph{on top of} the data array \(A\).  Slot \(i\) stores the sum of
one carefully chosen block of \(A\)'s entries -- the block
\([i-\operatorname{lowbit}(i)+1,\,i]\), where
\(\operatorname{lowbit}(i)=i\,\&\,(-i)\) is the lowest set bit of \(i\).
Those \(n\) blocks are exactly structured so that every prefix
\([1..i]\) decomposes into \(\le\log_2 n\) disjoint blocks (walk
\(i\to i-\text{lowbit}(i)\)) and every position \(i\) is covered by
\(\le\log_2 n\) blocks (walk \(i\to i+\text{lowbit}(i)\)).  Result: prefix
sums and point updates in \(\bigO{\log n}\), with no pointers, no recursion,
and no balancing.

\textbf{Prerequisites.} Binary representation of integers and two's-complement
negation (\(-i\) flips bits and adds one, so \(i\&(-i)\) isolates the lowest
set bit).  Prefix sums from Ch.~1/Ch.~5 (\propref{c24-prefix}).  The running
array \(A\) and its table \(T\) from \figref{c14-array}.

\textbf{After this section you can}
\begin{itemize}[nosep]
  \item read the responsibility interval of every slot off the binary
        expansion, and verify a built BIT cell by cell; \hfill\lv{2}
  \item write \code{query}/\code{update} from memory and bound their walk
        length by \(\log_2 n\); \hfill\lv{2}--\lv{3}
  \item build a BIT from its data array in \(\bigO{n}\) (not
        \(\bigO{n\log n}\)); \hfill\lv{4}
  \item solve point-set + range-sum inputs end to end (CSES~1648). \hfill\lv{3}
\end{itemize}
\end{topicintro}

\begin{recognitionbox}[Pattern card -- ``prefix sums with values that change'']
\reckey{``update position \(k\) to \(u\) / add \(u\) to position \(k\)''
interleaved with ``sum of \(x_a+\cdots+x_b\)''; ``after these changes, how
many of X are currently \(\le\) \(\ldots\)''; any input whose brute force is
a prefix array rebuilt per update -- the phrase ``process \(q\) queries of
two types'' with one type mutating one position.}
\recuse{
  \item point mutations + prefix/range \emph{sums} (or counts): the plain
        BIT, this section -- about 10 lines;
  \item the same but the ranges are tree subtrees -- Ch.~5's flattening
        \(\;+\) this engine, \secref{c14-trees};
  \item static array, zero updates -- plain prefix array
        (\propref{c24-prefix}), \(O(1)\) per query; do not pay \(\log n\).
}
\recnotuse{
  \item the aggregate is min/max/gcd under \emph{arbitrary} writes -- a
        segment tree, Ch.~19 (\chref{c13}); the BIT trick needs an
        \emph{invertible} accumulator (sums, counts) or a monotone
        insertion order (see \secref{c14-range}'s prefix-min caveat);
  \item whole \emph{intervals} get added to -- still a BIT, but the
        difference flavour, \secref{c14-range};
  \item you must read an arbitrary \emph{element} after range adds and also
        sum ranges -- two BITs, \(\secref{c14-range}\).
}
One-line test: \emph{single positions move, prefix questions get asked?  The
smallest structure that does both is ten lines long and lives in one array.}
\end{recognitionbox}

% ---------------------------------------------------------------
\subsection{Intuition: the Blocks Binary Gives You}
% ---------------------------------------------------------------

Write the indices in binary.  Each index \(i\) ``owns'' the block whose
length is its lowest set bit:
\[
  \textstyle
  \operatorname{lowbit}(i)=i\,\&\,(-i),\qquad
  \text{block of slot }i=[\,i-\operatorname{lowbit}(i)+1,\; i\,].
\]
For \(n=16\) the whole schedule is mechanical (opener,
\figref{c14-array}, bottom row): six slots own length~1 (odd indices), four
own length~2 (\(\equiv 2 \pmod 4\)), two own length~4 (\(\equiv4\pmod8\)),
one owns length~8, one owns length~16.  Two facts make the structure work:

\begin{itemize}[nosep,leftmargin=1.4em]
  \item \textbf{Prefixes shatter into owned blocks.}  \(i\to
        i-\operatorname{lowbit}(i)\) peels the \emph{last} block off
        \([1..i]\): \(13\to12\to8\to0\) shatters \([1..13]\) into
        \([13,13]\uplus[9,12]\uplus[1,8]\).  Each step clears at least the
        lowest set bit, so at most \(\log_2 n\) blocks survive the walk.
  \item \textbf{Positions are covered by few blocks.}  \(A_i\) is stored in
        exactly the slots whose blocks contain \(i\): \(i\to
        i+\operatorname{lowbit}(i)\) climbs to each of them --
        \(3\to4\to8\to16\), again \(\le\log_2 n+1\) stops.
\end{itemize}

\begin{definition}{Responsibility interval}{c14-interval}
In a BIT over \([1..n]\), slot \(i\) is \emph{responsible} for the contiguous
range of data positions \([\,i-\operatorname{lowbit}(i)+1,\;i\,]\), and stores
\(t[i]=\sum_{p}\,A_p\) over exactly that range.  The \term{parent} of slot
\(i\) is the next slot responsible for a range strictly containing it,
\(i+\operatorname{lowbit}(i)\) (if it exists).
\end{definition}

\begin{theorem}{Both walks in \(\bigO{\log n}\)}{c14-query}
On a BIT of size \(n\), \code{query(i)} (sum of \([1..i]\)) visits at most
\(\lfloor\log_2 n\rfloor+1\) slots, and \code{update(i,delta)} touches at
most \(\lfloor\log_2 n\rfloor+1\) slots.  The blocks collected by
\code{query} are disjoint and cover \([1..i]\) exactly.
\end{theorem}
\begin{proofbox}
\textbf{Query.}  Each step replaces \(i\) by \(i-\operatorname{lowbit}(i)\),
which clears the lowest set bit and changes no higher bit, after
\(\operatorname{popcount}(i)\le\lfloor\log_2 n\rfloor+1\) steps
\(i=0\).\ \textbf{Coverage.}  The step at value \(i\) collects
\([i{-}\operatorname{lowbit}(i){+}1,\,i]\); the next value
\(i' = i-\operatorname{lowbit}(i)\) has
\(\operatorname{lowbit}(i')>\operatorname{lowbit}(i)\) (the bit just cleared
was \(i\)'s lowest, so \(i'\)'s lowest set bit is higher), hence
\(i'\)'s block ends at \(i'\ge i-\operatorname{lowbit}(i)\) and abuts the
previous block without overlap; by induction the collected blocks are
disjoint with union \([1..i]\).\ \textbf{Update.}  Each step
\(i\to i+\operatorname{lowbit}(i)\) clears the lowest set bit of \(i\) that
is \emph{followed by a zero}... formally: adding the lowest set bit to \(i\)
carries through the consecutive ones above it, strictly increasing
\(\operatorname{lowbit}\); since \(\operatorname{lowbit}\) is a power of two
\(\le n\), at most \(\lfloor\log_2 n\rfloor+1\) additions fit before the walk
passes \(n\).  The touched slots are precisely those whose blocks contain
position~\(i\) (a block \([a,b]\) contains \(i\le b\) iff starting from
\(i\)'s own slot the parent chain reaches it, because parents strictly
extend and the chain enumerates all extending dyadic blocks).
\end{proofbox}

\noindent\textbf{In plain words.}  Binary keeps a bookkeeping secret: the
number of set bits and the length of a carry chain are both \(\le\log n\).
The query walk spends set bits; the update walk spends carry steps.  Nothing
about the \emph{data} matters -- only indices -- which is why the structure
survives being just one flat array.

% ---------------------------------------------------------------
\subsection{The Two Walks, Then the Linear Build}
% ---------------------------------------------------------------

\begin{lstlisting}[caption={reading lowbit, and the prefix walk: sum of A[1..i]}, label={lst:c14-query}]
int lowbit(int i){ return i & -i; }        // lowest set bit of i
long long query(int i){                    // prefix sum A[1]+..+A[i]
    long long s = 0;
    for (; i > 0; i -= lowbit(i)) s += t[i];
    return s;
}
\end{lstlisting}

\begin{lstlisting}[caption={the point-update walk: add delta to A[i] everywhere it is buried}, label={lst:c14-update}]
void update(int i, long long delta){       // A[i] += delta
    for (; i <= n; i += lowbit(i)) t[i] += delta;
}
// range sum [a,b] = query(b) - query(a-1)   (invertible accumulator!)
\end{lstlisting}

Both loops are the walks of \thmref{c14-query} with no branch, no pointer and
no cache-hostile indirection.  On \(A\) of \figref{c14-array}:
\(\text{query}(13)\) collects \(t[13]{=}1\), \(t[12]{=}13\), \(t[8]{=}24\)
for \(38\); \(\text{update}(3,+3)\) rewrites \(t[4],t[8],t[16]\).

\begin{lstlisting}[caption={O(n) build: hand each partial block to the slot whose block extends it}, label={lst:c14-build}]
void build(){                       // assumes A[1..n]; fills t[1..n]
    for (int i = 1; i <= n; i++) t[i] = A[i];
    for (int i = 1; i <= n; i++){
        int j = i + lowbit(i);      // j's block ends just past i's
        if (j <= n) t[j] += t[i];   // pass i's full block up
    }
}
\end{lstlisting}

\begin{proposition}{The linear build}{c14-linear}
\(\coderef{c14-build}\) initialises \(\code{t}[i]=A[i]\) and pushes each slot
into its parent \(i{+}\operatorname{lowbit}(i)\) once, in increasing \(i\).
After processing \(i\), slot \(i\) already holds the full block sum
\(\defref{c14-interval}\), because the pushes into \(i\) come exactly from
the child slots \(i_1 < i_2 < \cdots\) whose blocks partition
\(i\)'s block minus its own single last entry; and every \(i_1,i_2,\dots<i\)
was processed earlier.  Total work: \(n\) copies \(+\;\le n\) pushes
\(=\bigO{n}\), versus \(\bigO{n\log n}\) for calling \code{update} per
position.
\end{proposition}

Trace on the first eight slots (opener table): start \(t=(3,2,4,5,1,1,5,3)\).
Pushes: \(i{=}1\!: t[2]{+}{=}3\to5\). \(i{=}2\!: t[4]{+}{=}5\to10\).
\(i{=}3\!: t[4]{+}{=}4\to14\). \(i{=}4\!: t[8]{+}{=}14\to17\).
\(i{=}5\!: t[6]{+}{=}1\to2\). \(i{=}6\!: t[8]{+}{=}2\to19\).
\(i{=}7\!: t[8]{+}{=}5\to24\).  Final row \((3,5,4,14,1,2,5,24)\) -- note
\(T[2]=A_1{+}A_2=5\), not \(2\): a slot's value after the build is its
\emph{whole block}, and the push phase is what turns single entries into
block sums.  Continuing to \(n=16\) gives the shared row
\[
  T=(\,3,5,4,14,1,2,5,24,\;7,7,2,13,1,7,3,49\,)
\]
(easy cells to recheck: \(T[10]=A_9{+}A_{10}=7{+}0=7\),
\(T[12]=7{+}0{+}2{+}4{=}13\), \(T[16]=\sum A=49\)).  The walks are unaffected:
\(\text{query}(6)=T[6]+T[4]=2+14=16\) \checkmark,
\(\text{query}(13)=T[13]+T[12]+T[8]=1+13+24=38\) \checkmark.

\begin{complexitybox}[Plain BIT]
\textbf{Build} \(\bigO{n}\) (\(\coderef{c14-build}\)).
\textbf{Query/update} \(\le\lfloor\log_2 n\rfloor+1\) array touches, e.g.\
\(18\) at \(n=2\cdot10^5\) -- a fraction of a microsecond per op.
\textbf{Memory} one \code{long long} array of \(n+1\).
\textbf{Constant factor}: the smallest of any \(\bigO{\log n}\) range
structure (segment trees: \(4\) times the memory, a quarter of the cache
hits per level; comparison is \secref{c14-range}'s closing table).
\end{complexitybox}

% ---------------------------------------------------------------
\subsection{Why It Is a Hidden Binary Tree}
% ---------------------------------------------------------------

The parent rule \(i\to i+\operatorname{lowbit}(i)\) makes slots into a
forest; for \(n=16\) the top of each tree and the children of slot~8 are:

\begin{figure}[H]
\centering
\begin{tikzpicture}[diagram, scale=0.72, transform shape,
    bnode/.style={circle, draw=primaryblue, line width=0.8pt, fill=primaryblue!8,
                  minimum size=6.5mm, font=\scriptsize\bfseries, text=primaryblue!60!black}]
  \node[bnode] (s8) at (0,0)   {8};
  \node[bnode] (s4) at (-2.2,-1.4) {4};
  \node[bnode] (s12) at ( 2.2,-1.4) {12};
  \node[bnode] (s2) at (-3.4,-2.8) {2};
  \node[bnode] (s6) at (-1.2,-2.8) {6};
  \node[bnode] (s10) at ( 1.2,-2.8) {10};
  \node[bnode] (s14) at ( 3.4,-2.8) {14};
  \node[bnode] (s1) at (-4.0,-4.1) {1};
  \node[bnode] (s3) at (-2.9,-4.1) {3};
  \node[bnode] (s5) at (-1.8,-4.1) {5};
  \node[bnode] (s7) at (-0.6,-4.1) {7};
  \node[bnode] (s9) at ( 0.6,-4.1) {9};
  \node[bnode] (s11) at ( 1.8,-4.1) {11};
  \node[bnode] (s13) at ( 2.9,-4.1) {13};
  \node[bnode] (s15) at ( 4.0,-4.1) {15};
  \node[bnode] (s16) at (0,2.2) {16};
  \draw[tedge] (s16)--(s8) (s8)--(s4) (s8)--(s12)
    (s4)--(s2) (s4)--(s6) (s12)--(s10) (s12)--(s14)
    (s2)--(s1) (s2)--(s3) (s6)--(s5) (s6)--(s7)
    (s10)--(s9) (s10)--(s11) (s14)--(s13) (s14)--(s15);
  \node[annot, font=\scriptsize, text=darkgray, anchor=north west] at (4.9,2.2)
    {\(\begin{array}{@{}l@{}}
      \text{parent of }i = i+\text{lowbit}(i)\\
      \text{children of }8=\{4,12\};\ \text{of }4=\{2,6\}\\
      \text{query path: child}\to\text{root sums}\le\log n\ \text{siblings}\\
      \text{update path: root}\to\text{leaf refreshes the path}
    \end{array}\)};
\end{tikzpicture}
\caption{The BIT over \([1..16]\) drawn as the forest its parent rule
defines.  Slot~16 roots the whole array (block \([1,16]\)); the children of
any slot \(s\) are \(s-\operatorname{lowbit}(s)/2\) and its recursive
split, mirroring a segment tree \emph{whose split points are the binary
carries instead of the midpoints}.  The \code{query} walk is a root path in
this forest read bottom-up with sibling sums skipped; \code{update} is the
same path top-down.  This is why textbook proofs of \(\bigO{\log n}\) need no
amortisation: the forest has height \(\lceil\log_2 n\rceil\).}
\label{fig:c14-hiddentree}
\end{figure}

\begin{comparisonbox}[Prefix array vs BIT vs segment tree (sum workload)]
\begin{itemize}[nosep,leftmargin=1.4em]
  \item \textbf{Prefix array}: build \(\bigO{n}\), query \(O(1)\), point
        update \(\bigO{n}\) -- dead the moment values move.
  \item \textbf{BIT}: build \(\bigO{n}\), query \(\bigO{\log n}\), update
        \(\bigO{\log n}\), \(\sim\)10 lines, one array of \(n{+}1\) -- the
        default for sums/counts under point changes.
  \item \textbf{Segment tree} (Ch.~19, \chref{c13}): build \(\bigO{n}\),
        query/update \(\bigO{\log n}\), \(\sim\)40 lines, array of
        \(4n\) -- required the moment the accumulator stops being
        invertible (max, min, gcd) or the updates stop being additive.
\end{itemize}
The honest gap between BIT and segment tree is \emph{expressiveness}, not
speed; the BIT wins on code size, constant factor and memory, and this book
keeps the segment tree's full power for \chref{c13}/\chref{c15}.
\end{comparisonbox}

% ---------------------------------------------------------------
\subsection{Problem Lab: CSES 1648 Dynamic Range Sum Queries}
% ---------------------------------------------------------------

\solvedproblem{c14-rangesum}{Dynamic Range Sum Queries}{CSES 1648}
  {https://cses.fi/problemset/task/1648}{CSES L3}
  {plain BIT, point set via delta, range sum by prefix difference}

\begin{problemstatement}[CSES 1648 Dynamic Range Sum Queries]
Given an array of \(n\) integers, process \(q\) queries of two types:
\texttt{1 k u} --- set the value at position \(k\) to \(u\) (not add);
\texttt{2 a b} --- print the sum of values in range \([a,b]\).

\textbf{Input.} \(n,q\le 2\cdot10^5\); values \(1\le x_i,u\le10^9\);
\(1\le a\le b\le n\).

\textbf{Output.} The answer to each type-2 query, one per line.

\textbf{Example.}
\begin{verbatim}
Input:
8 4
3 2 4 5 1 1 5 3
2 1 4
2 5 6
1 3 1
2 1 4

Output:
14
2
11
\end{verbatim}
\end{problemstatement}

\begin{attemptgate}[Attempt it yourself first -- \lv{3}, give yourself 20 minutes]
The brute force is exactly the two textbook approaches colliding: a prefix
array answers \texttt{2} in \(O(1)\) but needs \(\bigO{n}\) per
\texttt{1}; re-summing the window per \texttt{2} is \(\bigO{n}\) too, and
\(n\cdot q=4\cdot10^{10}\) says \tle.  The task is to make \emph{both} sides
pay \(\log n\) and nothing more.
\peek{which quantity changes when one position moves, and by how much?}
\nudge{set \(x_3=1\) from \(4\) is really ``add \(-3\) to position 3'' --
       an \emph{increment} is what a BIT stores, so keep the raw array to
       compute the delta.}
\fullhint{Maintain a plain BIT over \(x\).  Type~1: \(\delta = u - x[k]\);
\(x[k]\!\leftarrow\!u\); \code{update(k,}\(\delta\)\code{)}.  Type~2:
\code{query(b)-query(a-1)}.  Everything is \(\bigO{\log n}\); sums need
\code{long long} (\(2\cdot10^5\cdot10^9=2\cdot10^{14}\)).  The full build is
below.}
\end{attemptgate}

\solstep{Understanding}
Two operations, one shared object: the multiset of \emph{prefix sums}.  A
type-2 answer is a difference of two prefixes (the accumulator is invertible
-- this is the license to use a BIT at all); a type-1 set perturbs every
prefix \(\ge k\) by the same \(\delta=u-x[k]\), which is exactly what the
update walk of \thmref{c14-query} propagates through the \(O(\log n)\) blocks
containing position~\(k\).  Nothing about the problem is tree-shaped yet --
this is the pure array-index case; \secref{c14-trees} adds the tree later.

\solstep{Key Observations}
\begin{observation}
\texttt{set} \(\ne\) \texttt{add}: the BIT supports adding increments to
slots; to honour a set you must know the \emph{current} value, so a shadow
array \code{x[]} stays.  Forgetting this is this problem's signature WA.
\end{observation}
\begin{observation}
A range \([a,b]\) is answered as \(\mathit{pref}(b)-\mathit{pref}(a{-}1)\)
-- the walk never computes a range directly, and \(a=1\) needs
\(\mathit{pref}(0)\overset{\text{def}}{=}0\), which \coderef{c14-query}'s
loop condition gives for free.
\end{observation}
\begin{observation}
Bounds: max magnitude \(2\cdot10^{14}\) fits \code{long long} with margin;
\(4\cdot10^5\) walks of \(\le18\) touches each \(=\lesssim7\cdot10^6\) array
ops -- the 1.00\,s limit is a non-issue once the \(\bigO{n}\) build (or
\(n\) updates) is in place.
\end{observation}

\solstep{Think It Through With Me}
\begin{thinkbox}
\deadend{Rebuild the prefix array after every type-1 query}
{This is the \(\bigO{n}\)-per-update brute force in disguise: \(2\cdot10^5\)
updates of \(O(n)\) is \(4\cdot10^{10}\) -- \tle.  The whole point of the
BIT is that a change at one position touches only \(\log n\) stored blocks.}
\deadend{\code{update(k, u)} instead of \code{update(k, u - x[k])}}
{Silently wrong for every set that is not an ``add \(u\)''; passes samples
whose first update happens to be a large increase.  The BIT is an
\emph{incremental} structure; absolute writes must be converted to deltas
against the shadow array.}
\aha{Set = add of (new $-$ old); range = difference of prefixes; and the lowbit
walks make both directions logarithmic at once.}
\end{thinkbox}

\solstep{The Approach Ladder}
\begin{center}\renewcommand{\arraystretch}{1.25}
\begin{tabular}{@{}c p{0.235\linewidth} p{0.30\linewidth} l c@{}}
\toprule
\textbf{\#} & \textbf{Approach} & \textbf{Idea} & \textbf{Time} & \textbf{Verdict}\\
\midrule
1 & \lv{2}\ re-sum window & loop \(a..b\) per type-2 & \(\bigO{q\cdot n}\) & \tle \\
2 & \lv{2}\ rebuild prefixes & fresh \(\mathit{pref}\) after each set & \(\bigO{q\cdot n}\) & \tle \\
3 & \lv{3}\ \textbf{plain BIT} & delta-update + prefix difference (\(\coderef{c14-query}\)--\(\coderef{c14-update}\)) & \(\bigO{(n+q)\log n}\) & \ac \\
4 & \lv{4}\ segment tree & works, but 4x memory + full code for an invertible op & \(\bigO{(n+q)\log n}\) & overkill \\
\bottomrule
\end{tabular}\end{center}
Rung~4 is the Ch.~19 tool brought to a knife fight; rung~3 is why this
chapter exists before it.

\solstep{Solution}
Store the array \(x[1..n]\) (1-based -- the lowbit machinery \emph{requires}
index \(0\) to be unused).  Build the BIT once with
\(\coderef{c14-build}\).  For \texttt{1 k u}: \(\delta = u-x[k]\),
\code{update(k,}\(\delta\)\code{)}, \(x[k]=u\).  For \texttt{2 a b}: print
\(\code{query(b)}-\code{query(a-1)}\).

\solstep{Proof of Correctness}
\begin{proofbox}
\textbf{Invariant.} For every \(i\), \(t[i]=\sum_{p=i-\operatorname{lowbit}(i)+1}^{i} x[p]\).
\emph{Init:} \propref{c14-linear} establishes it. \emph{Step:} a type-1 at
position~\(k\) changes exactly the blocks containing~\(k\); by
\thmref{c14-query} the update walk visits exactly those slots and adds
\(\delta\) to each, restoring the invariant elsewhere trivially (other blocks
omit \(k\)).
\textbf{Answer:} \thmref{c14-query}'s disjoint-cover property gives
\(\code{query}(i)=\sum_{p\le i}x[p]\) for all \(i\); subtracting at
\(a{-}1\) (with \(\code{query}(0)=0\), the empty walk) yields the range sum.
\textbf{Termination/bounds:} both walks shrink/grow \(i\) strictly past \(n\)
within \(\lfloor\log_2 n\rfloor+1\) steps; magnitudes \(\le 2\cdot10^{14}\)
stay inside \code{long long}.
\end{proofbox}

\solstep{Implementation}
\begin{lstlisting}[caption={CSES 1648 full solution (plain BIT, set-as-delta)}, label={lst:c14-1648-soln}]
#include <bits/stdc++.h>
using namespace std;
const int MX = 200005;
int n, q, x[MX];
long long t[MX];

int lowbit(int i){ return i & -i; }
long long query(int i){
    long long s = 0;
    for (; i > 0; i -= lowbit(i)) s += t[i];
    return s;
}
void update(int i, long long d){
    for (; i <= n; i += lowbit(i)) t[i] += d;
}
void build(){                          // O(n), prop:c14-linear
    for (int i = 1; i <= n; i++) t[i] = x[i];
    for (int i = 1; i <= n; i++){
        int j = i + lowbit(i);
        if (j <= n) t[j] += t[i];
    }
}
int main(){
    scanf("%d%d", &n, &q);
    for (int i = 1; i <= n; i++) scanf("%d", &x[i]);
    build();
    while (q--){
        int ty, a, b; scanf("%d%d%d", &ty, &a, &b);
        if (ty == 1){
            long long d = (long long)b - x[a];   // set k->u as delta
            x[a] = b; update(a, d);
        } else printf("%lld\n", query(b) - query(a-1));
    }
}
\end{lstlisting}

\solstep{Trace on the Official Sample}
\begin{tracebox}[CSES 1648: a=(3,2,4,5,1,1,5,3), then set x[3]=1]
After the \(\bigO{n}\) build, \(t=(3,5,4,14,1,2,5,24)\) (the first eight
cells of the corrected opener row).
\[
\renewcommand{\arraystretch}{1.15}
\begin{array}{c|l|l|l}
\text{\#} & \text{query} & \text{walk} & \text{result}\\ \hline
1 & \texttt{2 1 4} & \text{query}(4){=}t[4]{=}14;\ \text{query}(0){=}0 & 14\\
2 & \texttt{2 5 6} & \text{query}(6){=}t[6]{+}t[4]{=}2{+}14{=}16;\ \text{query}(4){=}14 & 16-14=2\\
3 & \texttt{1 3 1} & \delta=1{-}4{=}-3:\ t[4]{=}11,\ t[8]{=}21 & \text{(update)}\\
4 & \texttt{2 1 4} & \text{query}(4){=}t[4]{=}11 & 11
\end{array}
\]
Printed: \texttt{14 2 11} -- the official output.  Post-update sanity:
\(\text{query}(7)=t[7]+t[6]+t[4]=5+2+11=18=(3{+}2{+}1{+}5{+}1{+}1{+}5)\).
\end{tracebox}

\solstep{Complexity}
\begin{complexitybox}
\(\bigO{n}\) build \(+\;\bigO{\log n}\) per query:
\(\lesssim 2\cdot10^5 + 4\cdot10^5\cdot18\approx7.6\cdot10^6\) primitive
ops.  Memory \(8\cdot(2\cdot10^5)\) bytes for \code{t}, plus the shadow
\code{int} array.  \textbf{AC} with ~0.1\,s typical.
\end{complexitybox}

\solstep{Verification and Edge Cases}
\begin{center}\renewcommand{\arraystretch}{1.2}
\begin{tabular}{@{}p{0.20\linewidth} p{0.30\linewidth} p{0.38\linewidth}@{}}
\toprule
\textbf{Case} & \textbf{Why tricky} & \textbf{Code behaviour}\\
\midrule
\(a=1\) range & \(\mathit{pref}(0)\) must be \(0\) & \code{query(0)} loop body never runs \\
set to same value & \(\delta=0\) & harmless no-op walk (could skip: \(x[a]==b\)) \\
\(n=1\) & single slot, parent chain empty & \code{build} pushes nothing; walks stop at \(i>n\) \\
large sums & \(2\cdot10^{14}\) overflows \code{int} & \code{t}, \code{query} are \code{long long} \\
last index update & \(k=n\): walk visits \(n\) then exits & loop bound \(i\le n\) correct \\
\bottomrule
\end{tabular}\end{center}

\begin{hardeningbox}
\begin{itemize}[nosep,leftmargin=1.4em]
  \item \textbf{1-index or die:} slot \(0\) has no block (\(\text{lowbit}(0)=0\)
        makes both walks spin forever).  Reading 0-based input?  Convert at
        the door: \code{x[i]} for \(i=1..n\), and shift query endpoints.
  \item \textbf{Two's complement note:} \code{i \& -i} relies on signed
        integers being two's complement (true on all contest judges; C++20
        made it mandatory for the builtin types).
  \item \textbf{Fast IO:} \(q=2\cdot10^5\) lines -- \code{scanf} as above or
        \code{cin} with \code{sync\_with\_stdio(false)} is enough; the
        structure is never the bottleneck.
  \item \textbf{Array sizing:} \code{t[MX]} with \(MX=n_{\max}+5\); an
        off-by-one makes \code{update} at \(i=n\) write \code{t[n]} legally
        and nothing else -- but a too-small \(MX\) is silent UB.
\end{itemize}
\end{hardeningbox}

% ---------------------------------------------------------------
\subsection{Common Pitfalls}
% ---------------------------------------------------------------

\begin{warningbox}[Trap 1: mixing the two walks' directions]
\code{query} must \emph{subtract} the lowbit (\(6{\to}4{\to}0\));
\code{update} must \emph{add} it (\(3{\to}4{\to}8{\to}16\)).  Swapping the
signs produces an infinite loop (update walking down into \(0\)) or a wrong
prefix (query walking up).  The mnemonic: \emph{questions descend the forest,
answers climb it} (\figref{c14-hiddentree}).
\end{warningbox}

\begin{warningbox}[Trap 2: \code{update(i, value)} instead of \code{update(i, delta)}]
The BIT accumulates increments; a type-1 ``set'' is only an
\(\text{add}(u - x_i)\).  Keep the shadow array.  (Ch.~6's LC~307 -- the
mutable range-sum problem you may already know -- has the same trap, and its
treap/prefix solutions there all convert to deltas first.)
\end{warningbox}

\begin{warningbox}[Trap 3: \code{int} accumulator]
\(\sum x_i\) reaches \(2\cdot10^{14}\) on this very problem's constraints;
the first \(10^9\)-sized test that overflows is also the one that flips a
range sum negative.  \code{long long} for \code{t[]}, \code{query}'s return,
and the delta \emph{before} the cast.
\end{warningbox}

\begin{warningbox}[Trap 4: \(n\) updates from zeros as the build]
Calling \code{update}(i, x[i]) for every position costs \(\bigO{n\log n}\)
and is \emph{correct} -- but when \(n=2\cdot10^5\) and the time limit is
1.00\,s with several test files, the \(\bigO{n}\) build of
\(\coderef{c14-build}\) is free insurance.  Never rebuild by updates inside
a loop over test cases.
\end{warningbox}

% ---------------------------------------------------------------
\subsection{Practice Problems}
% ---------------------------------------------------------------

\begin{stuckbox}
\begin{itemize}[nosep,leftmargin=1.4em]
  \item Did you convert every ``set'' into a delta against a shadow array?
  \item Are you 1-indexed everywhere, with \code{query(0)}\(=0\) as the
        \(a=1\) guard?
  \item Is every sum type \code{long long} (array, return, printf)?
  \item Build by \(\coderef{c14-build}\) or by \(n\) updates -- which did
        you actually write?
\end{itemize}
\end{stuckbox}

\begin{practiceset}
  \practice{Forest Queries}{CSES 1652}
    {https://cses.fi/problemset/task/1652}
    {CSES L3}
    {Static \(n\cdot n\) grid, rectangle queries ``how many trees in
     \([x_1..x_2][y_1..y_2]\)'' -- the \(O(1)\) 2D \emph{prefix} cousin of
     this section: inclusion--exclusion over four corner values of
     \(\mathit{pref}[i][j]\).  Solve it to feel exactly what the BIT later
     buys you: this works only because nothing ever changes.}
  \practice{Multidimensional Queries}{CF 1093G}
    {https://codeforces.com/problemset/problem/1093/G}
    {CF 2300 L6}
    {Points in a \(k\le4\)-dimensional space, ``set a point's coordinates''
     and ``sum of points in an axis-aligned box''.  The box becomes
     \(2^k\) signed prefix queries of a \(k\)-dimensional BIT -- nested
     lowbit loops, \(O(\log^k n)\) per op -- and the \(\bigO{n}\) build of
     \(\coderef{c14-build}\) generalises axis by axis.  Check your loop
     nesting against \(\coderef{c14-1648-soln}\) first.}
\end{practiceset}

\paragraph{Back-cites and forwards.}
Chapter~5 (\propref{c24-prefix}) owns the static prefix array this replaces
under updates.  Forward: \secref{c14-range} changes \emph{what the slots
store} (differences, pairs); \secref{c14-order} changes \emph{what the
indices mean} (ranks); \secref{c14-trees} changes \emph{what the ranges are}
(subtrees); \secref{c14-offline} changes \emph{when the queries run}
(sorted).  Segment trees, the heavier rival, are \chref{c13}/\chref{c15}.

% ---------------------------------------------------------------
\begin{recallbox}[Section 18.1 in 10 seconds]
Slot \(i\) owns block \([i{-}\operatorname{lowbit}(i){+}1,i]\);
\(\operatorname{lowbit}(i)=i\mathbin{\&}(-i)\).  Query walks \(i\to i{-}\text{lowbit}\)
(peeling blocks, \(\le\log n\)); update walks \(i\to i{+}\text{lowbit}\)
(refreshing containers); range \(=\) prefix difference (invertibility
licence!).  Build in \(\bigO{n}\) by pushing \(t[i]\to t[i{+}\text{lowbit}]\).
Set-queries need deltas plus a shadow array; sums need \code{long long}.
Worked: CSES~1648 (\(\probref{c14-rangesum}\), \(\coderef{c14-1648-soln}\)).
\selftest{
  \item List the blocks the walk \(13\to\cdots\to0\) collects on \(n=16\).
  \item Which slots does \code{update(10,}\(\delta\)\code{)} touch?
  \item Why is \code{lowbit(0)}=0 a segfault magnet, and what invariant of
        the code prevents visiting \(0\) in \code{update}?
  \item Give the one-line reason a segment tree is \emph{not} needed here.
  \item Convert ``\texttt{1 k u}'' to the BIT vocabulary in one sentence.
}
\end{recallbox}

% =====================================================================
%  18.2  Range Variants: Differences, Pairs, Compression   sec:c14-range
%  Exports: def:c14-diff  thm:c14-twobit  prop:c14-compress
%           prop:c14-prefixmin  lst:c14-diff-build  lst:c14-diff-range
%           lst:c14-twobit  lst:c14-compress  lst:c14-min  lst:c14-1651-soln
%           prob:c14-rangeupd
%  Shared numbers (all hand-reverified, ASCII only):
%    D = (3,-1,2,1,-4,0,4,-2); diff-BIT Td = BIT over D = (3,2,2,5,-4,-4,4,3);
%    add +1 to [2,5]: Td -> (3,3,2,6,-4,-5,4,3); reads: p4=6, p2=3, p7=5, p8=3.
%    Two-BIT (T1 over D, T2 over j*D_j), prefix(x)=(x+1)Q1(x)-Q2(x):
%    original T1=(3,2,2,5,-4,-4,4,3), T2=(3,1,6,11,-20,-20,28,3);
%    prefix(4)=5*5-11=14, prefix(6)=7*1+9=16, prefix(3)=4*4-7=9;
%    after +1 on [2,5]: T1=(3,3,2,6,-4,-5,4,3), T2=(3,3,6,13,-20,-26,28,-1);
%    sum[1,4]=5*6-13=17, sum[1,6]=7*1+13=20, sum[4,6]=20-11=9.
%  CSES 1651 sample (live-verified): 8 3 / 3 2 4 5 1 1 5 3 /
%    2 4 / 1 2 5 1 / 2 4  ->  outputs 5 then 6 (x[4]=5; after +1 on [2,5], 6).
%  Two-BIT convention used EVERYWHERE here: T1 over D, T2 over j*D_j, and
%    prefix(x) = (x+1)Q1(x) - Q2(x)   (thm:c14-twobit, lst:c14-twobit).
%    Verified on a=(3,2,4,5,1,1,5,3): prefix(4)=5*5-11=14=sum a[1..4];
%    after +1 on [2,5]: T1=(3,3,2,6,-4,-5,4,3), T2=(3,3,6,13,-20,-26,28,-1),
%    sum[1,4]=5*6-13=17; sum[4,6]=20-11=9 (see the side-by-side tracebox).
% =====================================================================
\section{Range Variants: Adding to Intervals, Summing Intervals}\label{sec:c14-range}

\begin{previouslybox}[18.1 stores sums of \emph{data}; now store sums of \emph{changes}]
The plain BIT answers a prefix of \(A\) and mutates one entry of \(A\).
This section swaps which side gets the expensive operation: encode the
\emph{difference array} \(D_i=A_i-A_{i-1}\) in the BIT, and a whole-interval
addition collapses to two point updates -- at the price that a read becomes
a prefix walk.  Push the algebra one step further (a \emph{paired} BIT) and
even ``range add \emph{and} range sum'' become logarithmic.  Two supporting
tools travel with them: coordinate compression (when indices are not small
integers) and the prefix-min variant (when the accumulator is \(\min\), not
\(\sum\)).
\end{previouslybox}

\begin{topicintro}
\textbf{Prerequisites.} \secref{c14-core} in full -- walks, blocks,
\(\coderef{c14-build}\).  The telescoping identity
\(A_i=\sum_{j\le i}(A_j-A_{j-1})\) (Ch.~5 used it for edge-span counts).

\textbf{After this section you can}
\begin{itemize}[nosep]
  \item build the difference BIT and answer range-add/point-read inputs
        (CSES~1651); \hfill\lv{3}--\lv{4}
  \item derive and implement the two-BIT range-add/range-sum identity from
        the contribution count, not from memory; \hfill\lv{4}--\lv{5}
  \item coordinate-compress arbitrary values into BIT indices without
        breaking equal-range semantics; \hfill\lv{4}
  \item state exactly when a prefix-min BIT is legal (append-only decreases)
        and when it silently lies. \hfill\lv{4}
\end{itemize}
\end{topicintro}

\begin{recognitionbox}[Pattern card -- ``intervals move, positions/segments are asked'']
\reckey{``increase each value in \([a,b]\) by \(u\)'' + ``value at position
\(k\)'' (CSES~1651's exact words); ``add to a range, then ask a range sum''
(two BITs); values up to \(10^9\) but only \(\le2\cdot10^5\) of them appear
(compression); ``\(\mathit{dp}[i]=\min_{j<i,\;\text{cond}}\mathop(\ldots)\)
with \(j\) queried by a prefix of \(\min\) values'' (prefix-min BIT).}
\recuse{
  \item range add, point read -- difference BIT, this section (row~3 of the
        opener's cost chart);
  \item range add \emph{and} range sum -- two paired BITs
        (\thmref{c14-twobit}); the segment tree with lazy (\chref{c15}) also
        does it, at 4x the code;
  \item index domain huge, order the only thing that matters -- compress
        first (\propref{c14-compress});
  \item accumulator \(\min\) with monotone (decreasing-only) writes --
        prefix-min BIT (\propref{c14-prefixmin}).
}
\recnotuse{
  \item point \emph{sets} and range sums -- the plain BIT of
        \secref{c14-core} (differences would force you to re-point two
        slots per set and buy nothing);
  \item arbitrary overwrites of a \(\min\) table -- the monotonicity
        hypothesis of \propref{c14-prefixmin} fails; go to \chref{c13};
  \item range \emph{assign} (set all of \([a,b]\) to \(v\)) -- neither
        difference nor pair survives non-invertible assignment; lazy Ch.~20.
}
One-line test: \emph{if every mutation touches a whole interval, something
telescoping is being stored -- put \emph{that} in the BIT.}
\end{recognitionbox}

% ---------------------------------------------------------------
\subsection{The Difference BIT: Range Add, Point Read}
% ---------------------------------------------------------------

Let \(D_1=A_1\), \(D_i=A_i-A_{i-1}\).  Then \(A_i=\sum_{j\le i}D_j\): a
\emph{point read of \(A\) is a prefix query of \(D\)} -- exactly the
operation a BIT is born for.  The magic is the other direction:

\begin{definition}{Difference encoding of interval adds}{c14-diff}
Adding \(\delta\) to every \(A_p\) with \(p\in[a,b]\) changes only two
differences: \(D_a\mathrel{+}=\delta\) (the step up into the interval) and
\(D_{b+1}\mathrel{-}=\delta\) (the step back out).  All other \(D_j\) are
untouched because inside the interval both neighbours moved by \(\delta\).
\end{definition}

So the whole workload ``range add / point read'' becomes: BIT over \(D\);
\(\texttt{add}(a,\delta),\ \texttt{add}(b{+}1,-\delta)\); read
\(\code{query}(k)\).  \(\bigO{\log n}\) per op, three lines of code.

\begin{lstlisting}[caption={the difference BIT: build from A by differencing (O(n)), range add as two point updates}, label={lst:c14-diff-build}]
// t[] is a BIT over D; every "A[i]" below is 1-based
void build_diff(){                       // O(n): t[i] = D[i] first...
    for (int i = 1; i <= n; i++)
        t[i] = A[i] - (i > 1 ? A[i-1] : 0);
    for (int i = 1; i <= n; i++){        // ...then the same push as
        int j = i + lowbit(i);           // coderef:c14-build
        if (j <= n) t[j] += t[i];
    }
}
void range_add(int a, int b, long long d){
    update(a, d);                        // t is the point-update walk of
    if (b + 1 <= n) update(b+1, -d);     // coderef:c14-update on D
}
long long point_read(int k){ return query(k); }   // prefix of D = A[k]
\end{lstlisting}

\begin{tracebox}[Difference cascade on the CSES 1651 sample array]
\(a=(3,2,4,5,1,1,5,3)\) gives the raw difference array
\(D=(3,-1,2,1,-4,0,4,-2)\).  The BIT is built on top of \(D\) by
\(\coderef{c14-diff-build}\): start \(t=D\) and push \(t[i]\to t[i{+}\text{lowbit}(i)\)
(\(\coderef{c14-build}\)), giving the \emph{block sums}
\[
  Td=(\,3,2,2,5,-4,-4,4,3\,),
\]
cell by cell: \(Td[1]=3\), \(Td[2]=D_1{+}D_2=2\), \(Td[3]=D_3=2\),
\(Td[4]=\sum_{1..4}D=5\), \(Td[5]=D_5=-4\), \(Td[6]=D_5{+}D_6=-4\),
\(Td[7]=D_7=4\), \(Td[8]=\sum_{1..8}D=3\).  This is \(Td\), \emph{not} the
raw \(D\) above -- conflating them is Pitfall~1 below.

Now \(\texttt{range\_add}(2,5,+1)\) fires \(\mathit{update}(2,+1)\) (walk
\(2{\to}4{\to}8\): \(Td[2]\!{\to}\!3,\ Td[4]\!{\to}\!6,\ Td[8]\!{\to}\!4\)) and
\(\mathit{update}(6,-1)\) (walk \(6{\to}8\): \(Td[6]\!{\to}\!-5,\ Td[8]\!{\to}\!3\)),
so the live tree is \(Td=(3,3,2,6,-4,-5,4,3)\).  Point reads are prefix walks:
\(\text{read}(4)=Td[4]=6=a_4{+}1\) \checkmark,\quad
\(\text{read}(2)=Td[2]=3=a_2{+}1\) \checkmark,\quad
\(\text{read}(7)=Td[7]{+}Td[6]{+}Td[4]=4-5+6=5=a_7\) \checkmark\ (position~7
lies outside \([2,5]\)); \(\text{read}(8)=Td[8]=3=a_8\) is unchanged because
the two difference writes contributed \(+1-1=0\) to \(Td[8]\) -- the encoding
cancels exactly as \defref{c14-diff} promises.
\end{tracebox}

% ---------------------------------------------------------------
\subsection{Two BITs: Range Add \emph{and} Range Sum}
% ---------------------------------------------------------------

Difference reads a point; but a \emph{sum} of \(A\)'s prefix in terms of
\(D\) weights each difference by how many entries it still affects:
\[
  \sum_{i=1}^{x} A_i
  = \sum_{i=1}^{x}\sum_{j\le i} D_j
  = \sum_{j\le x} D_j\,(x-j+1)
  = (x+1)\underbrace{\sum_{j\le x} D_j}_{S_1(x)}
    \;-\; \underbrace{\sum_{j\le x} j\,D_j}_{S_2(x)} .
\]
Two plain BITs -- one over \(D_j\), one over \(j\,D_j\) -- answer range
add/range sum in \(4\) updates and \(4\) walks per op.

\begin{theorem}{The two-BIT identity}{c14-twobit}
Maintain BITs \(T_1\) over \(D\) and \(T_2\) over \(j\cdot D_j\).  Then
\(\displaystyle\mathit{prefix}(x)=(x+1)\,\code{query}_1(x)-\code{query}_2(x)\)
and a range add \(\delta\) on \([a,b]\) is
\[
\text{at }a:\ \mathit{add_1}(a,\delta),\ \mathit{add_2}(a,a\delta);
\qquad
\text{at }b{+}1:\ \mathit{add_1}(b{+}1,-\delta),\ \mathit{add_2}(b{+}1,-(b{+}1)\delta).
\]
Every operation is \(\bigO{\log n}\).
\end{theorem}
\begin{proofbox}
The prefix formula is the telescoping rearrangement above -- each \(D_j\)
lifts \(A_j,\dots,A_x\), i.e.\ \(x-j+1\) entries, so its contribution to
\(\sum_{i\le x}A_i\) is \((x+1)D_j-jD_j\).  A range add changes exactly
\(D_a{+}=\delta\), \(D_{b+1}{-}=\delta\) (\defref{c14-diff}); each change
must be entered in \emph{both} trees with the respective weight
(\(\delta\) vs.\ \(j\delta\)), giving the four updates.  \code{query}
correctness is \thmref{c14-query} applied to each tree.
\end{proofbox}

\begin{lstlisting}[caption={range add + range sum with two BITs (worked example, not a linked problem: CSES has no such task)}, label={lst:c14-twobit}]
// T1 over D, T2 over j*D[j]; both use coderef:c14-update/coderef:c14-query
void ra(int i, long long d){                 // one difference point change
    for (int k = i; k <= n; k += lowbit(k)){
        t1[k] += d;                          // d contribution
        t2[k] += d * i;                      // (j * d) contribution, j = i
    }
}
void range_add(int a, int b, long long d){
    ra(a, d);  ra(b + 1, -d);
}
long long bitq(long long* t, int i){
    long long s = 0;
    for (; i > 0; i -= lowbit(i)) s += t[i];
    return s;
}
long long prefix(int x){ return (x + 1) * bitq(t1, x) - bitq(t2, x); }
long long range_sum(int a, int b){ return prefix(b) - prefix(a - 1); }
\end{lstlisting}

\begin{tracebox}[Both BITs side by side: start a=(3,2,4,5,1,1,5,3), add +1 to [2,5], then sum [1,4] and [4,6]]
Before any update, \(D=(3,-1,2,1,-4,0,4,-2)\).  \(T_1\) is the BIT of block
sums over \(D\) -- exactly the \(Td\) row of the box above,
\(T_1=(3,2,2,5,-4,-4,4,3)\).  \(T_2\) is the BIT over \(jD_j=(3,-2,6,4,-20,0,28,-16)\),
which \(\coderef{c14-build}\) turns into \(T_2=(3,1,6,11,-20,-20,28,3)\).
Sanity check on the untouched array: \(\mathit{sum}[1,4]=\mathit{prefix}(4)=5{\cdot}Q_1(4)-Q_2(4)=5{\cdot}5-11=14=3{+}2{+}4{+}5\) \checkmark.

Apply \(\texttt{range\_add}(2,5,+1)\): four point writes -- at \(j=2\) put
\((\delta,\,j\delta)=(+1,+2)\) into \((T_1,T_2)\), at \(j=6\) put \((-\delta,\,-j\delta)=(-1,-6)\).
Their update walks give \(T_1=(3,3,2,6,-4,-5,4,3)\) and \(T_2=(3,3,6,13,-20,-26,28,-1)\).
Now the range sums read off the live pair:
\[
\mathit{sum}[1,4]=\mathit{prefix}(4)=5{\cdot}T_1[4]-T_2[4]=5{\cdot}6-13=17=(3{+}3{+}5{+}6)\ \checkmark ,
\]
\[
\mathit{sum}[4,6]=\mathit{prefix}(6)-\mathit{prefix}(3)=(7{\cdot}1+13)-(4{\cdot}5-9)=20-11=9=(6{+}2{+}1)\ \checkmark .
\]
Both trees moved by the \emph{same} two difference events; the second simply
weights each by its index \(j\), which is what turns a point read into a range
sum.  This is \(\thmref{c14-twobit}\) running, no memorised constant required.
\end{tracebox}

% ---------------------------------------------------------------
\subsection{Coordinate Compression}
% ---------------------------------------------------------------

A BIT is an array indexed by \(1..n\), so it silently assumes the things you
store live at small \emph{integer} positions.  When the real indices are huge
(values up to \(10^9\)) or non-integral, only their \emph{relative order} can
matter to a prefix count -- and order is exactly what survives being replaced
by rank.  Compression is the one-time relabelling that makes that legal.

\begin{proposition}{Compression preserves the game}{c14-compress}
Let \(\mathit{vals}\) be the sorted, deduplicated set of all values that will
\emph{ever} be indexed (read, updated, or compared against).  Replacing each
value by its index in \(\mathit{vals}\) turns the value domain into
\([1..m]\), \(m\le\) (number of events), and preserves every comparison
\(v_1\le v_2\), hence every prefix count/sum the BIT answers.  Cost:
\(\bigO{M\log M}\) once.
\end{proposition}

\begin{lstlisting}[caption={standard compression: collect, sort, unique, lower\_bound}, label={lst:c14-compress}]
// Compress BEFORE any update/query: every value that will ever be an index
// must already sit in vals when the sort happens.  See thm:c14-twobit for a
// value-weighted tree when magnitudes (not just order) are needed.
vector<int> vals;                          // 1. COLLECT every value that
for (...) vals.push_back(x);               //    will ever be an index
sort(vals.begin(), vals.end());            // 2. SORT + 3. UNIQUE
vals.erase(unique(vals.begin(), vals.end()), vals.end());
int rank_(int x){                          // 4. RANK by lower_bound
    return (int)(lower_bound(vals.begin(), vals.end(), x) - vals.begin()) + 1;
}
// rank_ returns a 1-based index into vals -- the BIT's [1..m] domain.
// from here on, the BIT is sized m = vals.size(), never 1e9
\end{lstlisting}

% ---------------------------------------------------------------
\subsection{Prefix Min, and the Only Way It Is Legal}
% ---------------------------------------------------------------

\begin{proposition}{The append-only prefix-min BIT}{c14-prefixmin}
If the accumulator is \(\min\) and slot values only ever \emph{decrease} (the
standard DP pattern \(\mathit{dp}[i]=\min_{j\le i}\) over processed prefixes,
each \(j\) written once), then storing block minima and updating
\(t[k]=\min(t[k],v)\) along the \emph{update} walk answers
\(\min_{j\le x}\) along a \emph{superset-cover} walk: query visits the same
\(i\to i{-}\operatorname{lowbit}(i)\) blocks, which cover \([1..x]\) but now
may also be joined by over-covering ancestors -- for \(\min\) (idempotent)
extra coverage is harmless.  Arbitrary point increases are \textbf{not}
supported: the stale minimum remains buried in some slot with no rule for
which one to fix.
\end{proposition}

\begin{lstlisting}[caption={prefix-min BIT under monotone writes (query min over [1..x])}, label={lst:c14-min}]
void push(int i, long long v){             // position i may only decrease
    for (; i <= n; i += lowbit(i)) t[i] = min(t[i], v);
}
long long qmin(int x){                     // min over [1..x] (over-covers)
    long long m = INF;
    for (; x > 0; x -= lowbit(x)) m = min(m, t[x]);
    return m;
}
\end{lstlisting}

\begin{warningbox}[The prefix-min BIT answers min only while writes stay monotone]
``Cover a superset, harmless for \(\min\)'' is exactly as true for \(\max\);
what is \emph{not} true is mixing directions -- a \(\min\)-BIT whose writes
sometimes raise a value has no local repair rule (the raised entry may hide
in any ancestor).  When monotonicity breaks, you need recomputation over
blocks: that is a segment tree, \chref{c13}, and the honest cost is stating
so here instead of debugging a silent WA for a night.
\end{warningbox}

% ---------------------------------------------------------------
\subsection{BIT vs Segment Tree: the Boundary in One Table}
% ---------------------------------------------------------------

\begin{comparisonbox}[Where each range engine lives (and why this book teaches both)]
\begin{itemize}[nosep,leftmargin=1.4em]
  \item \textbf{BIT}: sums/counts (invertible) or monotone \(\min/\max\)
        writes; point or difference updates; \(\sim\)8 lines per flavour;
        trivially 2D-able (\(\secref{c14-offline}\), practice CF~1093G);
        cache-friendly half the memory of a segtree.
  \item \textbf{Segment tree} (Ch.~19, \chref{c13}): \emph{any} associative
        aggregate (max/min/gcd under arbitrary writes), non-invertible
        queries, custom combine functions; \(\sim\)40 lines; lazy variant
        (Ch.~20, \chref{c15}) for range assign/add with range asks.
  \item \textbf{Rule of thumb for contests}: if you can phrase every query
        as ``prefix of an invertible accumulator'' or a monotone DP write,
        the BIT is the lower-effort AC; the moment the statement asks a
        \emph{suffix min under increases}, \emph{range mode}, or
        \emph{first position where a non-additive predicate flips}, reach
        for \chref{c13}.
\end{itemize}
This is prose, per the chapter contract -- segtree code appears only there.
\end{comparisonbox}

% ---------------------------------------------------------------
\subsection{Problem Lab: CSES 1651 Range Update Queries}
% ---------------------------------------------------------------

\solvedproblem{c14-rangeupd}{Range Update Queries}{CSES 1651}
  {https://cses.fi/problemset/task/1651}{CSES L4}
  {difference BIT, range add as two point updates, point read as prefix}

\begin{problemstatement}[CSES 1651 Range Update Queries]
An array of \(n\) values; process \(q\) queries of two types:
\texttt{1 a b u} --- increase each value in \([a,b]\) by \(u\);
\texttt{2 k} --- print the value at position \(k\).

\textbf{Input.} \(n,q\le 2\cdot10^5\); \(|x_i|\le10^9\); \(|u|\le10^9\)
(negative updates allowed); \(1\le a\le b\le n\), \(1\le k\le n\).

\textbf{Output.} Each type-2 query on its own line.

\textbf{Example.}
\begin{verbatim}
Input:
8 3
3 2 4 5 1 1 5 3
2 4
1 2 5 1
2 4

Output:
5
6
\end{verbatim}
\end{problemstatement}

\begin{attemptgate}[Attempt it yourself first -- \lv{4}, give yourself 25 minutes]
The type-1 brute force loops \(a..b\): a single \([1,n]\) update is
\(\bigO{n}\), and \(q\) of them is \(\tle\).  The hint is in what type-1
\emph{does not} change: differences between adjacent values.
\peek{after adding \(+1\) to \([2,5]\), which adjacent pairs actually see a
change in \(A_i-A_{i-1}\)?}
\nudge{exactly two boundaries: at \(2\) (enters) and \(6\) (leaves) -- store
the difference array in the BIT of \secref{c14-core} and you already know
every remaining line of code.}
\fullhint{BIT over \(D\).  \texttt{1 a b u}: \code{update(a,u)},
\code{update(b+1,}-u\code{)} (guard \(b+1\le n\)).  \texttt{2 k}:
\code{query(k)} \(=\sum_{j\le k}D_j=A_k\).  Build: \(t[i]\!\leftarrow\!D_i\)
then the push loop of \(\coderef{c14-diff-build}\).  \code{long long}
everywhere: \(|u|\le10^9\), \(2\cdot10^5\) overlapping adds reach
\(2\cdot10^{14}\).}
\end{attemptgate}

\solstep{Understanding}
The problem never asks a range sum -- it asks \emph{single positions} after
\emph{whole-interval} moves.  That is the difference encoding's exact
signature (\defref{c14-diff}): the mutation is local in \(D\) (two slots),
and the read is global-in-prefix in \(D\) (one walk).  A plain BIT over
\(A\) would have it backwards: reads \(O(\log)\), writes \(O(n)\).

\solstep{Key Observations}
\begin{observation}
An interval add is a \emph{rank-one} change in the difference view: \(+u\) at
the left edge, \(-u\) at one past the right edge.  Everything inside shifts
together, so no difference between consecutive inside-positions changes.
\end{observation}
\begin{observation}
The initial build must difference \emph{first}, then BIT-build (\(\coderef{c14-diff-build}\));
feeding \(A\) into \(\coderef{c14-build}\) and then calling updates
``range-style'' answers \(A_k\) with a sum of raw values, not the value --
the Pitfall~1 trace above shows how plausible this looks.
\end{observation}
\begin{observation}
\(b=n\) is safe: \(\mathit{update}(n{+}1,\cdot)\) is out of the array's
domain \([1,n]\) and no read walk ever passes \(n\); guard with
\(b{+}1\le n\) (or size \(t\) to \(n{+}1\) and skip the guard).
\end{observation}

\solstep{Think It Through With Me}
\begin{thinkbox}
\deadend{Lazy segment tree (Ch.~20's tool)}
{Correct and \(\bigO{\log n}\) per op, but ~45 lines, \(4n\) memory, and a
propagation discipline to debug -- for ``range add, point read'' the lazy
machinery answers a question nobody asked (range aggregates).  Reserve it
for \chref{c15} proper.}
\deadend{Difference \emph{array} (no BIT)}
{The static array version supports each update in \(O(1)\) but each read in
\(\bigO{n}\) (re-prefix the whole array).  Interleaving is what kills it --
\(q\) reads after \(q\) updates is the same \(\bigO{nq}\).  The BIT is what
makes the prefix-of-differences incremental.}
\aha{Store in the tree what makes the \emph{update} cheap, if the \emph{query}
is an operation the tree already answers fast -- here differences turn
intervals into points, and prefix sums are the BIT's native tongue.}
\end{thinkbox}

\solstep{The Approach Ladder}
\begin{center}\renewcommand{\arraystretch}{1.25}
\begin{tabular}{@{}c p{0.235\linewidth} p{0.30\linewidth} l c@{}}
\toprule
\textbf{\#} & \textbf{Approach} & \textbf{Idea} & \textbf{Time} & \textbf{Verdict}\\
\midrule
1 & \lv{2}\ loop the interval & \(A_p+=u\) for \(p\in[a,b]\) & \(\bigO{nq}\) & \tle \\
2 & \lv{3}\ diff array only & two \(O(1)\) writes, re-prefix per read & \(\bigO{nq}\) & \tle \\
3 & \lv{4}\ \textbf{difference BIT} & two \(\bigO{\log}\) writes, \(\bigO{\log}\) read & \(\bigO{(n+q)\log n}\) & \ac \\
4 & \lv{4}\ lazy segtree & full power, mostly unused here & \(\bigO{(n+q)\log n}\) & overkill \\
\bottomrule
\end{tabular}\end{center}

\solstep{Solution}
Build \(D\) from the input, then the \(\bigO{n}\) BIT build
(\(\coderef{c14-diff-build}\)).  Type~1: two point updates at \(a\) and
\(b{+}1\).  Type~2: \code{query(k)}.  \code{long long} for the tree and the
answers.

\solstep{Proof of Correctness}
\begin{proofbox}
\textbf{Invariant.} After any sequence of processed queries, \(t\) is a BIT
over the current difference array \(D\) of the logical array \(A\).
\emph{Init:} \(\coderef{c14-diff-build}\) establishes it (\propref{c14-linear}
applied to \(D\)).  \emph{Type-1 step:} by \defref{c14-diff} the interval add
changes exactly \(D_a,D_{b+1}\) by \(+u,-u\); the two \code{update} walks
maintain the BIT for those point changes (\thmref{c14-query}).
\emph{Type-2 step:} \(A_k=\sum_{j\le k}D_j\) telescopes; \code{query(k)}
returns that prefix by \thmref{c14-query}'s disjoint-cover property.
\textbf{Bounds:} values \(|A_k|\le 10^9+2\cdot10^5\cdot10^9<3\cdot10^{14}\)
-- \code{long int long}; walks \(\le18\). \textbf{Termination:} each walk
strictly increases/decreases its index past the guards.
\end{proofbox}

\solstep{Implementation}
\begin{lstlisting}[caption={CSES 1651 full solution (difference BIT)}, label={lst:c14-1651-soln}]
#include <bits/stdc++.h>
using namespace std;
const int MX = 200005;
int n, q;
long long t[MX];

int lowbit(int i){ return i & -i; }
void update(int i, long long d){
    for (; i <= n; i += lowbit(i)) t[i] += d;
}
long long query(int i){
    long long s = 0;
    for (; i > 0; i -= lowbit(i)) s += t[i];
    return s;
}
int main(){
    scanf("%d%d", &n, &q);
    long long prev = 0, cur;
    for (int i = 1; i <= n; i++){          // build D in place
        scanf("%lld", &cur);
        t[i] = cur - prev; prev = cur;     // t = raw D so far
    }
    for (int i = 1; i <= n; i++){          // O(n) BIT push (build_diff)
        int j = i + lowbit(i);
        if (j <= n) t[j] += t[i];
    }
    while (q--){
        int ty, a, b; long long u;
        scanf("%d", &ty);
        if (ty == 1){
            scanf("%d%d%lld", &a, &b, &u);
            update(a, u);
            if (b + 1 <= n) update(b + 1, -u);
        } else { scanf("%d", &a); printf("%lld\n", query(a)); }
    }
}
\end{lstlisting}

\solstep{Trace on the Official Sample}
\begin{tracebox}[CSES 1651: read x[4], add +1 to [2,5], read x[4] again]
\(D=(3,-1,2,1,-4,0,4,-2)\); after the push phase the live tree is
\(Td=(3,2,2,5,-4,-4,4,3)\) (the row verified in the first trace box of this
section).
\[
\renewcommand{\arraystretch}{1.15}
\begin{array}{c|l|l|l}
\text{\#} & \text{query} & \text{walk} & \text{result}\\ \hline
1 & \texttt{2 4} & \text{query}(4)=Td[4]=5 & 5\\
2 & \texttt{1 2 5 1} & \text{update}(2,+1):\ Td[2]{\to}3,\ Td[4]{\to}6,\ Td[8]{\to}4 & \text{(add)}\\
  & \text{(cont.)} & \text{update}(6,-1):\ Td[6]{\to}-5,\ Td[8]{\to}3 & \\
3 & \texttt{2 4} & \text{query}(4)=Td[4]=6 & 6
\end{array}
\]
Printed: \texttt{5, 6} -- matches the live statement (each type-2 answer on
its own line).  Cross-check on \(\text{query}(8)=Td[8]=3=a_8\) after the
update: \(a_8=3\) was outside \([2,5]\), and \(Td[8]\) received \(+1-1=0\)
net -- the difference encoding cancelling itself exactly as
\defref{c14-diff} promises.
\end{tracebox}

\solstep{Complexity}
\begin{complexitybox}
\(\bigO{n}\) build, \(\bigO{\log n}\) per update/read; total
\(\lesssim 2\cdot10^5 + 3\cdot10^5\cdot18\approx5.6\cdot10^6\) ops.
Memory one \code{long long} array.  \textbf{AC}.
\end{complexitybox}

\solstep{Verification and Edge Cases}
\begin{center}\renewcommand{\arraystretch}{1.2}
\begin{tabular}{@{}p{0.20\linewidth} p{0.30\linewidth} p{0.38\linewidth}@{}}
\toprule
\textbf{Case} & \textbf{Why tricky} & \textbf{Code behaviour}\\
\midrule
\(b=n\) & \(D_{n+1}\) has no slot & guard \code{b+1<=n}; no read ever reaches \(n{+}1\) \\
negative \(u\) & tree entries go negative & \code{long long}, no logic change \\
overlapping adds & value drifts to \(2\cdot10^{14}\) & magnitudes bounded above; fits \\
first element & \(D_1=A_1\) (no \(A_0\)) & \code{prev=0} initialisation handles it \\
read after many adds & needs \emph{current} value & prefix walk reads live tree \\
\bottomrule
\end{tabular}\end{center}

\begin{hardeningbox}
\begin{itemize}[nosep,leftmargin=1.4em]
  \item \textbf{Do not mix flavours:} a tree built over \(D\) answers point
        reads; feeding it ``set \(A_k=u\)'' queries (as in CSES~1648) is
        \emph{wrong} -- a set at \(k\) is two delta-updates in the difference
        view (\(+\delta\) at \(k\), \(-\delta\) at \(k{+}1\)) \emph{and}
        still needs the shadow array.  Mixing problem types is this
        section's signature WA.
  \item \textbf{Guard or size-up:} pick one convention for \(b=n\) and keep
        it in the template (\code{if (b+1<=n)} above; or \(MX=n+2\) and no
        guard).
  \item \textbf{\(\pm u\) stress:} generate random \([a,b]\) adds and compare
        \code{query(k)} against a brute-force array for \(n\le50\), \(q\le100\)
        -- five lines of Python; catches an inverted sign at \(b{+}1\)
        instantly.
\end{itemize}
\end{hardeningbox}

% ---------------------------------------------------------------
\subsection{Common Pitfalls}
% ---------------------------------------------------------------

\begin{warningbox}[Trap 1: storing the raw difference array \emph{as if} it were the BIT]
\(D\) and the BIT over \(D\) are different objects: the second stores
\emph{block sums} of the first.  Skipping the push phase of
\(\coderef{c14-diff-build}\) (leaving \(t=D\)) makes \code{query} walk
individual \(D\) entries and return \(A_k\) only for the indices the walk
happens to touch -- e.g.\ on the sample, \(\text{query}(4)\) would answer
\(1\) (the raw \(D_4\)) instead of \(5\).  Wrong answers that ``look like
off-by-one'' are usually this.
\end{warningbox}

\begin{warningbox}[Trap 2: forgetting \(-u\) at \(b+1\)]
A one-sided difference update turns every later read into an accumulated
flood of \(+u\).  Visible only when a query position exceeds \(b\) -- many
hand samples never query there.  Say the two words out loud when writing:
\emph{enter, leave}.
\end{warningbox}

\begin{warningbox}[Trap 3: memorising the two-BIT formula with the wrong weight]
The pair is \(\big(D_j,\;(j{-}1)D_j\big)\) with
\(\mathit{prefix}(x)=x\,Q_1(x)-Q_2(x)\), \emph{or} \(\big(D_j,\;jD_j\big)\)
with \((x{+}1)Q_1-Q_2\) -- one convention, consistently.  The bug appears
when code and derivation come from different memories.  Re-derive in one
line on the page before submitting (the derivation is in
\thmref{c14-twobit}'s proof); do not recall.
\end{warningbox}

\begin{warningbox}[Trap 4: compression on a \emph{sum} that needs real values]
Compression preserves \emph{order} (counts, k-th, \(\le k\) predicates).  It
destroys \emph{magnitudes}: ``sum of values \(\le k\)'' cannot be answered
by counting compressed ranks.  Keep sums on a value-carrying structure or
store \(v\)-weighted trees -- \(\coderef{c14-twobit}\)'s second tree is the
template for weighting.
\end{warningbox}

% ---------------------------------------------------------------
\subsection{Practice Problems}
% ---------------------------------------------------------------

\begin{stuckbox}
\begin{itemize}[nosep,leftmargin=1.4em]
  \item Which side is cheap in your encoding -- update or read?  Choose the
        stored object so the \emph{frequent} op is the cheap one.
  \item Did you build the difference tree with the push phase, not raw \(D\)?
  \item Compression: are you using only \emph{order} properties?
  \item Two-BIT: which weight convention is your code actually using?
\end{itemize}
\end{stuckbox}

\begin{practiceset}
  \practice{Little Girl and Maximum Sum}{CF 276C}
    {https://codeforces.com/problemset/problem/276/C}
    {CF 1500 L3}
    {\(q\) fixed ranges, reorder the array so that the \emph{sum over the
     ranges} is maximal.  Count how many ranges cover each position --
     that is the difference-array idea of \defref{c14-diff} run once
     offline (two \(O(1)\) writes per range, one prefix sweep) -- then sort
     and match greedily.  No BIT survives the queries here (the array never
     changes), which makes the difference \emph{technique} visible without
     the tree.}
  \practice{Iahub and Xors}{CF 341D}
    {https://codeforces.com/problemset/problem/341/D}
    {CF 1900 L5}
    {A \(n\times n\) matrix, two ops: xor \(x\) into every cell of a
     subrectangle, and xor of a subrectangle.  \(\oplus\) is its own inverse
     (invertible accumulator -- the BIT licence holds), so the rectangle is
     a 2D difference encoding with four corner events per update, and a
     range-xor needs the same paired-tree trick as \thmref{c14-twobit} in
     two axes.  Budget a page of index algebra before coding; verify on a
     \(2\times2\) hand case first.}
\end{practiceset}

\paragraph{Back-cites and forwards.}
\secref{c14-core} owns the walks this section reuses without restating.
Chapter~5 (\propref{c24-prefix}) owns telescoping in its static form.
Forward: the frequency reading of compressed indices is
\secref{c14-order}; difference-on-Euler arrays is \secref{c14-trees}
(CF~383C); 2D sweep encodings return in \secref{c14-offline}.  Arbitrary
aggregates under arbitrary writes are the segment tree's -- \chref{c13},
\chref{c15}; the boundary table is above, the code is there.

% ---------------------------------------------------------------
\begin{recallbox}[Section 18.2 in 10 seconds]
Store \(D=A_{\mathit{diff}}\): interval add \(=\) two point updates
(\defref{c14-diff}); point read \(=\) prefix walk (CSES~1651,
\(\coderef{c14-1651-soln}\)).  Need range sums too?  Pair a second tree of
\((j{-}1)D_j\): \(\mathit{prefix}(x)=xQ_1-Q_2\) (\thmref{c14-twobit}) --
pick ONE weight convention.  Huge index domains: collect--sort--unique--
lower\_bound (\propref{c14-compress}); order survives, magnitudes don't.
\(\min\)-BIT legal only while writes are monotone (\propref{c14-prefixmin});
otherwise \chref{c13}.  Never confuse raw \(D\) with the BIT over \(D\).
\selftest{
  \item Which two slots does ``\(+5\) on \([3,7]\)'' touch in the difference
        tree, and what are the increments?
  \item Re-derive \(\mathit{prefix}(x)=xQ_1-Q_2\) in two lines from
        \(\sum_{j\le x}D_j(x-j+1)\).
  \item Give the workload where a lazy segment tree is the \emph{smaller}
        program than the two-BIT pair.
  \item Why does compression answer ``how many \(\le k\)'' but not ``sum of
        values \(\le k\)''?
  \item State the monotonicity hypothesis of the prefix-min BIT and the
        failure mode when it breaks.
}
\end{recallbox}

% =====================================================================
%  18.3  Order Statistics and Sweeps with a Frequency BIT   sec:c14-order
%  Exports: def:c14-freq  thm:c14-lift  prop:c14-inversion
%           lst:c14-freq  lst:c14-kth  lst:c14-inv  lst:c14-61-soln
%           fig:c14-lift  prob:c14-triplets
%  Shared numbers (agree with 00_opener, six-angles rows 5 and 6):
%    frequency (weight) array w=(3,2,4,1,5,5,2,6), total 28,
%      prefix=(3,5,9,10,15,20,22,28), BIT  tf=(3,5,4,10,1,6,5,28);
%      kth(6)=3, kth(11)=5  (opener says "first prefix >=6 -> slot 3; >=11 -> slot 5").
%    inversion sweep on a=(10,8,3,1): ranks (4,3,2,1),
%      smaller-after=(3,2,1,0), inversions=6, decreasing triplets=4.
%    CF 61E (Enemy is weak, 1900) samples, live-verified:
%      3 / 3 2 1 -> 1 ;  3 / 2 3 1 -> 0 ;  4 / 10 8 3 1 -> 4 ;  4 / 1 5 4 3 -> 1.
%    leftGreater/rightLess on (10,8,3,1): leftG=(0,1,2,3), rightL=(3,2,1,0),
%      sum of products = 0*3+1*2+2*1+3*0 = 4.
% =====================================================================
\section{Order Statistics and Sweeps with a Frequency BIT}\label{sec:c14-order}

\begin{previouslybox}[Ch.~1 counted these offline; Ch.~6 owns the tree versions]
Chapter~1 (\chref{c01}) counted inversions by splitting the array and
merging -- divide and conquer, all the data on the table at once.  The
answer is the same number; here we get it \emph{online}, one element at a
time, by reading the BIT not as ``a sum over positions'' but as ``a count
over values''.  Order-statistic \emph{trees} (PBDS, the multiset of Ch.~6,
\chref{c06}) answer the same questions with pointers and rotations; the
frequency BIT is their flat-array twin -- it needs the values to be ranks,
which \secref{c14-range}'s compression already taught us to produce.
\end{previouslybox}

\begin{topicintro}
Reinterpret slot \(i\) of a BIT: instead of ``the block of positions
\([i{-}\text{lowbit}(i){+}1,\,i]\)'', read \(i\) as a \emph{value} (a rank)
and \(t[i]\) as ``how many stored elements fall in the value-range of that
block''.  Then \(\code{query}(k)\) counts elements \(\le k\) and
\(\code{add}(\mathit{rank},+1)\) inserts one element -- a full multiset with
\(\bigO{\log}\) insert, delete, and rank-count.  The single new trick is the
\emph{reverse} query: given a count \(k\), recover the value that is the
\(k\)-th smallest, by walking the tree top-down a bit at a time.

\textbf{Prerequisites.} \secref{c14-core} (the two walks) and
\secref{c14-range} (coordinate compression: values \(\to\) ranks).  The
running frequency array \((3,2,4,1,5,5,2,6)\) and the sweep array
\((10,8,3,1)\) from \figref{c14-array}'s six-angle table.

\textbf{After this section you can}
\begin{itemize}[nosep]
  \item drive a BIT as a multiset: insert, erase, ``how many \(\le k\)''; \hfill\lv{3}
  \item find the \(k\)-th smallest in \(\bigO{\log n}\) by binary lifting; \hfill\lv{4}
  \item count inversions and \emph{decreasing triplets} with one or two
        directional sweeps (CF~61E); \hfill\lv{4}--\lv{5}
  \item state, precisely, why this beats Ch.~1's merge count when updates
        interleave and matches it when they do not. \hfill\lv{4}
\end{itemize}
\end{topicintro}

\begin{recognitionbox}[Pattern card --``counts across a sorted value axis'']
\reckey{``how many elements are \(\le k\) / \(< k\)'' with insertions and
deletions; ``the \(k\)-th smallest'' of a changing multiset; ``for each
element, count the ones after/before it that are smaller/greater''
(inversions, LC~315's shape); ``count triples \(i<j<k\) with
\(a_i>a_j>a_k\)'' (CF~61E); ``smallest value whose running count reaches
\(k\)''; any problem whose brute force is a scan that \emph{compares to one
element} and \emph{counts over a prefix}.}
\recuse{
  \item rank/count over a value range, online -- frequency BIT (this
        section): \(\code{add}\) at a rank, \(\code{query}\) counts \(\le k\);
  \item \(k\)-th smallest / select -- binary lifting (\(\thmref{c14-lift}\));
  \item pairwise or triple counts with an index constraint -- sweep one
        direction, keep a frequency BIT of what you've seen
        (\(\propref{c14-inversion}\), CF~61E).
}
\recnotuse{
  \item the values arrive in one pass and you only ever need the
        \emph{median} of a \emph{stream} -- the two heaps of Ch.~8
        (\chref{c07}) are \(O(\log)\) with less setup; here select is
        offline/static or by rank;
  \item you must also read arbitrary \emph{sums of values} (not counts) of a
        multiset under order statistics -- store a value-carrying tree (the
        paired trick of \(\secref{c14-range}\));
  \item the range aggregate is over \emph{positions} with an arbitrary
        combine (max/min/gcd) -- segment tree, Ch.~19 (\chref{c13}); a
        frequency BIT counts \emph{by value}, not by position.
}
One-line test: \emph{is the question a \emph{rank} among \emph{values} (how
many below / which is the \(k\)-th)?  Compress the values, and let the index
of the BIT \emph{be} the value.}
\end{recognitionbox}

% ---------------------------------------------------------------
\subsection{The Frequency BIT: an Index that Means a Value}
% ---------------------------------------------------------------

\begin{definition}{Frequency BIT}{c14-freq}
A \term{frequency BIT} (a.k.a.\ \term{binary indexed multiset}) is a BIT
\(\mathit{tf}[1..m]\) whose index \(r\) stands for the \(r\)-th
\emph{rank} in a sorted set of values.  \(\code{add}(r,d)\) changes the
multiplicity of rank \(r\) by \(d\); \(\code{count\_le}(r)=\sum_{p\le r}
\mathit{tf}[p]\) is the number of stored elements \(\le\) the \(r\)-th
value.  The accumulator is a \emph{count}, so it is invertible and the plain
prefix-difference of \(\secref{c14-core}\) applies: elements strictly below
rank \(r\) are \(\code{count\_le}(r-1)\).
\end{definition}

\begin{lstlisting}[caption={the frequency BIT: add a copy of a rank, count how many sit at or below a rank}, label={lst:c14-freq}]
int lowbit(int i){ return i & -i; }
void add(int r, int d){                 // one more item of rank r
    for (; r <= m; r += lowbit(r)) tf[r] += d;
}
int count_le(int r){                    // items with rank <= r
    int s = 0;
    for (; r > 0; r -= lowbit(r)) s += tf[r];
    return s;
}
// strictly below r:  count_le(r-1)      items in (l..r]: count_le(r)-count_le(l-1)
\end{lstlisting}

The weight array \(w=(3,2,4,1,5,5,2,6)\) of \figref{c14-array} is read here
as ``rank~1 appears three times, rank~2 twice, \ldots'' -- \(\sum w=28\)
items total -- and its frequency BIT is
\(\mathit{tf}=(3,5,4,10,1,6,5,28)\).

\begin{figure}[H]
\centering
\begin{tikzpicture}[diagram, scale=0.82, transform shape]
  \foreach \i/\vv in {1/3,2/5,3/4,4/10,5/1,6/6,7/5,8/28} {
    \draw[primaryblue, line width=0.6pt, fill=primaryblue!6] (\i*0.9,0) rectangle (\i*0.9+0.9,0.9);
    \node[font=\scriptsize\ttfamily] at (\i*0.9+0.45,0.45) {\vv};
    \node[font=\tiny\ttfamily, text=darkgray] at (\i*0.9+0.45,1.12) {tf[\i]};}
  \node[font=\scriptsize, text=primaryblue!60!black, anchor=east] at (0,0.45)
    {\(\mathit{tf}:\)};
  \foreach \i/\pv in {1/3,2/5,3/9,4/10,5/15,6/20,7/22,8/28} {
    \node[font=\tiny\ttfamily, text=green!45!black] at (\i*0.9+0.45,-0.55) {\pv};}
  \node[font=\tiny, text=green!45!black, anchor=east] at (0,-0.55) {prefix:};
  \draw[warningred, line width=1.0pt, -latex] (0.15,-1.05) -- node[below,font=\scriptsize,text=warningred] {step 4: swallow tf[4]=10 < 11} (4*0.9+0.75,-1.05);
  \draw[warningred, line width=1.0pt, -latex] (4*0.9+0.85,-1.05) -- node[below,font=\scriptsize,text=warningred] {then 5} (5*0.9+0.75,-1.05);
\end{tikzpicture}
\caption{The frequency BIT \(\mathit{tf}\) over \(w=(3,2,4,1,5,5,2,6)\), its
cumulative counts on the middle row, and one \emph{lifting} descent for
\(k{=}11\): the search swallows the length-4 block (\(\mathit{tf}[4]{=}10<11\)),
then cannot take the next two, and stops with answer \(\mathit{idx}{+}1=5\).
The red walk is \(\thmref{c14-lift}\); the numeric trace is the box below.}
\label{fig:c14-lift}
\end{figure}

% ---------------------------------------------------------------
\subsection{Binary Lifting: the k-th Smallest in \texorpdfstring{\(\bigO{\log n}\)}{O(log n)}}
% ---------------------------------------------------------------

\(\code{count\_le}\) turns a value into a count; we want the inverse -- turn
a count into a value.  Because a BIT is the forest of \figref{c14-hiddentree},
we can decide the answer one bit at a time: walk a candidate index
\(\mathit{idx}\) from the most significant bit down, and \emph{keep} a
candidate step whenever the block we would swallow is still strictly short of
the target.  This is the same ``drop one bit per step'' budget as
\thmref{c14-query}, run in reverse -- we build the largest prefix whose sum
stays \(<k\), and the answer is the next index.

\begin{theorem}{Lifting to the first prefix reaching \(k\)}{c14-lift}
Let \(\mathit{tf}\) be a frequency BIT over \(m\) ranks and let \(k\ge1\) be
at most the total count.  The following returns the smallest rank
\(r\) with \(\code{count\_le}(r)\ge k\): start \(\mathit{idx}=0\), take
\(\mathit{step}=2^{\lfloor\log_2 m\rfloor}\), and while \(\mathit{step}>0\)
set \(\mathit{idx}\mathrel{+}=\mathit{step},\;k\mathrel{-}=\mathit{tf}[\mathit{idx}]\)
whenever \(\mathit{idx}{+}\mathit{step}\le m\) and \(\mathit{tf}[\mathit{idx}{+}\mathit{step}]<k\);
else just halve \(\mathit{step}\).  Return \(\mathit{idx}+1\).  Runs in
\(\bigO{\log m}\).
\end{theorem}
\begin{proofbox}
The invariant is ``\(\mathit{idx}\) is the largest index with
\(\sum_{p\le \mathit{idx}} w_p \le k_{\text{orig}} - k_{\text{now}}\) that is
a prefix of the already-fixed high bits, and \(k_{\text{now}}\) is the count
still unaccounted for beyond \(\mathit{idx}\)''.  At step size
\(\mathit{step}=2^{b}\), the block \(\mathit{tf}[\mathit{idx}{+}\mathit{step}]\)
covers exactly the \(\mathit{step}\) ranks immediately after \(\mathit{idx}\)
(because \(\mathit{idx}\)'s fixed high bits make
\(\operatorname{lowbit}(\mathit{idx}{+}\mathit{step})=\mathit{step}\)).  If
that block sum is \(<k\) we may safely consume all of it and set the bit to
\(1\); otherwise the target rank lies inside the block, we leave the bit
\(0\), and descend.  Every bit is decided once, so \(\le\lceil\log_2
m\rceil\) steps, and by construction \(\mathit{idx}\) ends as the largest
rank whose cumulative count is \(<k_{\text{orig}}\), making \(\mathit{idx}+1\)
the smallest rank whose cumulative count is \(\ge k_{\text{orig}}\).
\end{proofbox}

\noindent\textbf{In plain words.}  Reading a frequency BIT as a cumulative
distribution, lifting is the binary search that a segment-tree descent also
performs -- but it never touches a child pointer, only the flat
\(\mathit{tf}[]\).  With \(\mathit{tf}\) storing \emph{ranks of equal items},
\(\text{kth}(k)\) is literally the \(k\)-th smallest; with \(\mathit{tf}\)
storing \emph{weights}, it is ``the first slot whose running total reaches
\(k\)'', the shape CSES~1143's booking query would need if it were a
\emph{sum} threshold.

\begin{lstlisting}[caption={binary lifting on the BIT: smallest index whose prefix sum reaches k}, label={lst:c14-kth}]
// smallest r with count_le(r) >= k ; needs 1 <= k <= total, m >= 1
int kth(int k){
    int idx = 0;
    int step = 1;
    while (step <= m) step <<= 1;          // power of two strictly above m
    for (step >>= 1; step > 0; step >>= 1){
        int next = idx + step;             // candidate block just past idx
        if (next <= m && tf[next] < k){    // whole block still short of k
            idx  = next;                   //   -> swallow it
            k   -= tf[next];               //   -> owe k fewer items
        }
    }
    return idx + 1;                        // first prefix reaching the target
}
\end{lstlisting}

\begin{tracebox}[Lifting on w=(3,2,4,1,5,5,2,6): kth(6) and kth(11)]
The tree is \(\mathit{tf}=(3,5,4,10,1,6,5,28)\), \(m=8\), so
\(\mathit{step}\) starts at \(8\).

\emph{\(k=6\):} step~8: \(\mathit{tf}[8]=28\not<6\), skip. step~4:
\(\mathit{tf}[4]=10\not<6\), skip. step~2: next\(=2\),
\(\mathit{tf}[2]=5<6\) \(\Rightarrow\) \(\mathit{idx}=2,\;k=6-5=1\). step~1:
next\(=3\), \(\mathit{tf}[3]=4\not<1\), skip. Return \(\mathit{idx}+1=3\)
\checkmark\ (\(\text{prefix}=(3,5,\mathbf{9})\) first reaches \(\ge6\) at
slot~3).

\emph{\(k=11\):} step~8: \(28\not<11\). step~4: \(\mathit{tf}[4]=10<11\)
\(\Rightarrow\) \(\mathit{idx}=4,\;k=1\). step~2: next\(=6\),
\(\mathit{tf}[6]=6\not<1\). step~1: next\(=5\), \(\mathit{tf}[5]=1\not<1\)
(strict). Return \(4+1=5\) \checkmark\ (\(\text{prefix}=(\ldots,10,\mathbf{15})\):
slot~4 gives \(10<11\), slot~5 gives \(15\ge11\)).  Both match
\figref{c14-lift} and the opener's lifting row exactly.
\end{tracebox}

% ---------------------------------------------------------------
\subsection{Inversions as a One-Sided Sweep}
% ---------------------------------------------------------------

\begin{proposition}{Inversions by a right-to-left frequency sweep}{c14-inversion}
Sweep the array from the right, inserting each element's rank into a fresh
frequency BIT.  When the pointer is at position \(i\), the BIT already holds
exactly the elements \(a_{i+1},\dots,a_n\), so
\(\code{count\_le}(\mathit{rank}_i-1)\) is the number of them \emph{strictly
smaller} than \(a_i\) -- the ``smaller-after self'' count of position \(i\).
Summing these over \(i\) is the number of inversions (pairs \(i<j\),
\(a_i>a_j\)).  Total cost \(\bigO{n\log n}\) after compression; each query is
a single walk.
\end{proposition}

\begin{lstlisting}[caption={inversion / smaller-after count: one right-to-left sweep}, label={lst:c14-inv}]
// ranks r[1..n] in 1..m from coderef:c14-compress; tf[] starts all zero
long long inv = 0;
for (int i = n; i >= 1; i--){
    inv += count_le(r[i] - 1);          // items to the RIGHT that are smaller
    add(r[i], 1);                       // this one joins the seen multiset
}
// inv == sum of the smaller-after counts; the per-i value IS LC 315's answer
\end{lstlisting}

\begin{comparisonbox}[Same count, two mechanics (Ch.~1 vs this section)]
\begin{itemize}[nosep,leftmargin=1.4em]
  \item \textbf{Ch.~1 merge sort (\chref{c01})}: divide, conquer, and charge
        inversions at merge time.  \(\bigO{n\log n}\), \emph{offline} -- the
        whole array must exist, and it cannot answer an \emph{interleaved}
        update without re-sorting.
  \item \textbf{Frequency BIT (here)}: one walk per element, \(\bigO{n\log
        n}\), \emph{online} -- and it generalises to ``smaller after'',
        ``greater before'', and triplet products for free, because the seen
        multiset is a data structure, not a merge step.
\end{itemize}
If the task is \emph{only} the inversion count of a fixed array, Ch.~1 is the
shorter program; the moment the question varies per element (LC~315's output
list) or the array mutates, the sweep wins.  This is the stated division of
labour, not a re-derivation of merge sort.
\end{comparisonbox}

\begin{tracebox}[Smaller-after and inversions on a=(10,8,3,1)]
Compress ascending: \(1{\to}1,\,3{\to}2,\,8{\to}3,\,10{\to}4\), so
\(r=(4,3,2,1)\).  Right-to-left sweep (the BIT holds ranks seen so far):
\[
\renewcommand{\arraystretch}{1.15}
\begin{array}{c|c|c|c|c}
i & a_i & r_i & \text{count\_le}(r_i-1)\ \text{(smaller after)} & \text{seen set}\\ \hline
4 & 1  & 1 & 0 & \{1\}\\
3 & 3  & 2 & 1 & \{1,2\}\\
2 & 8  & 3 & 2 & \{1,2,3\}\\
1 & 10 & 4 & 3 & \{1,2,3,4\}
\end{array}
\]
Smaller-after \(=(3,2,1,0)\) (reading positions \(1..4\)), sum \(=6\) -- the
inversion count of a strictly decreasing length-4 array, \(\binom{4}{2}=6\)
\checkmark.  These same partial counts are the input to the triplet sweep
below, and match the opener's sweep row exactly.
\end{tracebox}

% ---------------------------------------------------------------
\subsection{Problem Lab: CF 61E Enemy is weak}
% ---------------------------------------------------------------

\solvedproblem{c14-triplets}{Enemy is weak}{CF 61E}
  {https://codeforces.com/problemset/problem/61/E}{CF 1900 L5}
  {frequency BIT over compressed ranks, two directional sweeps, per-middle product}

\begin{problemstatement}[CF 61E Enemy is weak]
An army is a line of \(n\) men with pairwise distinct powers
\(a_1,\dots,a_n\).  Its \emph{weakness} is the number of triplets
\((i,j,k)\) with \(i<j<k\) and \(a_i>a_j>a_k\).  Count them.

\textbf{Input.} \(3\le n\le 10^6\); \(n\) distinct integers
\(1\le a_i\le 10^9\).

\textbf{Output.} A single integer -- the weakness (it may exceed \(2^{31}\)).

\textbf{Example.}
\begin{verbatim}
Input:
4
10 8 3 1

Output:
4
\end{verbatim}
Other official samples: \texttt{3 / 3 2 1 -> 1},
\texttt{3 / 2 3 1 -> 0}, \texttt{4 / 1 5 4 3 -> 1}.
\end{problemstatement}

\begin{attemptgate}[Attempt it yourself first -- \lv{5}, give yourself 30 minutes]
Three nested loops are \(\bigO{n^3}\) -- at \(n=10^6\), hopeless.  A triplet
\((i,j,k)\) is really a \emph{middle} \(j\) glued to an independent choice
on its left and on its right: for a fixed \(j\), the count is
\((\#\,i<j,\;a_i>a_j)\cdot(\#\,k>j,\;a_k<a_j)\).  Can you make each factor a
prefix of a frequency BIT?
\peek{separate the two conditions around the middle index \(j\).}
\nudge{two independent sweeps: left-to-right fills ``greater to the left of
\(j\)'', right-to-left fills ``smaller to the right of \(j\)''; then sum the
products.  Each sweep is exactly \(\coderef{c14-inv}\).}
\fullhint{Compress to ranks.  Forward: \(\mathit{lg}[j]=(j-1)-
\code{count\_le}(r_j)\) then \(\code{add}(r_j,1)\).  Clear \(\mathit{tf}\),
backward: \(\mathit{rl}[j]=\code{count\_le}(r_j-1)\) then
\(\code{add}(r_j,1)\).  Answer \(\sum_j \mathit{lg}[j]\cdot\mathit{rl}[j]\)
in \code{long long}.  Two \(\bigO{n\log n}\) passes, memory \(\bigO{n}\).}
\end{attemptgate}

\solstep{Understanding}
The index order \(i<j<k\) and the value order \(a_i>a_j>a_k\) are two
different axes: position and magnitude.  Pin the middle \(j\) and the problem
\emph{factors}: everything to the left only has to beat \(a_j\), everything
to the right only has to lose to \(a_j\).  Each side is a running ``how many
values exceed / fall below a threshold'' question -- precisely what a
frequency BIT over ranks answers in one walk.  This is
\propref{c14-inversion} run twice, in both directions, then multiplied.

\solstep{Key Observations}
\begin{observation}
The distinctness of the \(a_i\) means every rank is unique; the ``strictly
greater / strictly smaller'' conditions become clean rank comparisons with no
tie handling: greater-left \(=(j-1)-\text{count\_le}(r_j)\), smaller-right
\(=\text{count\_le}(r_j-1)\).
\end{observation}
\begin{observation}
A triplet is fully determined by its middle plus one left and one right
choice, so the total is \(\sum_j \mathit{lg}[j]\cdot\mathit{rl}[j]\) -- no
over- or under-counting, because each valid \((i,j,k)\) has exactly one
middle \(j\).
\end{observation}
\begin{observation}
\(\binom{10^6}{3}\approx1.7\cdot10^{17}\) fits \code{long long} (\(<9.2\cdot
10^{18}\)) but not \code{int}; and \(n=10^6\) rules out any
\(\bigO{n\sqrt n}\) or cache-hostile tree -- this is a flat-array sweep, and
fast I/O is not optional.
\end{observation}

\solstep{Think It Through With Me}
\begin{thinkbox}
\deadend{Fix \(j\) and scan both sides naively}
{That is \(\bigO{n^2}\) per choice of \(j\)... really \(\bigO{n^2}\) total
with the right prefix idea, but \(\bigO{n^3}\) if done plainly; at
\(n=10^6\) even \(\bigO{n^2}\) is \(10^{12}\).  Only a per-side
\(\bigO{\log}\) query survives.}
\deadend{Count triplets in a single left-to-right pass}
{One pass cannot see the future: \(\mathit{rl}[j]\) needs the elements
\(\ge j\).  Trying to fold the right count into the left sweep makes the
product depend on data not yet read -- this is exactly why the sweep runs
twice, in opposite directions.}
\aha{The middle index splits a triplet into two independent rank-counts; a
frequency BIT gives each count in \(\bigO{\log n}\), so the weakness is a sum
of products of two one-dimensional sweeps.}
\end{thinkbox}

\solstep{The Approach Ladder}
\begin{center}\renewcommand{\arraystretch}{1.25}
\begin{tabular}{@{}c p{0.235\linewidth} p{0.30\linewidth} l c@{}}
\toprule
\textbf{\#} & \textbf{Approach} & \textbf{Idea} & \textbf{Time} & \textbf{Verdict}\\
\midrule
1 & \lv{2}\ three loops & test every \((i,j,k)\) & \(\bigO{n^3}\) & \tle \\
2 & \lv{3}\ fix middle, scan & for each \(j\), count sides linearly & \(\bigO{n^2}\) & \tle \\
3 & \lv{5}\ \textbf{two BIT sweeps} & lg, rl by rank-counts; sum products (\(\coderef{c14-inv}\)) & \(\bigO{n\log n}\) & \ac \\
4 & \lv{4}\ merge-sort variant & possible but awkward for a per-middle product & \(\bigO{n\log n}\) & clunkier \\
\bottomrule
\end{tabular}\end{center}
Rung~3 is the point: the middle-index factorisation is what makes a
\emph{triplet} problem reduce to \emph{pair}-counting, which the frequency
BIT already solves.

\solstep{Solution}
Compress the distinct powers to ranks \(1..n\).  Left-to-right sweep to fill
\(\mathit{lg}[j]\) (greater to the left); reset the frequency BIT; right-to-
left sweep to fill \(\mathit{rl}[j]\) (smaller to the right).  Print
\(\sum_j \mathit{lg}[j]\cdot\mathit{rl}[j]\) as \code{long long}.

\solstep{Proof of Correctness}
\begin{proofbox}
\emph{Compression:} \propref{c14-compress} lets us replace the powers by
their ranks with no change to any comparison (distinct \(\Rightarrow\) rank
order equals value order).  \emph{Left count:} before \(\code{add}(r_j,1)\),
the forward BIT holds exactly ranks of \(a_1..a_{j-1}\); \(\code{count\_le}(r_j)\)
counts those \(\le a_j\), so \((j-1)-\code{count\_le}(r_j)\) counts those
\(\;>\;a_j\)=\(\mathit{lg}[j]\).  \emph{Right count:} symmetrically the
backward BIT holds \(a_{j+1}..a_n\) and \(\code{count\_le}(r_j-1)\) counts
those \(<a_j\)=\(\mathit{rl}[j]\).  \emph{Product:} a triplet with middle
\(j\) picks one greater-left and one smaller-right independently, so it is
counted once among \(\mathit{lg}[j]\cdot\mathit{rl}[j]\); summing over all
\(j\) counts every valid \((i,j,k)\) exactly once (each has a unique middle).
\emph{Bounds:} \(\bigO{n\log n}\); magnitude \(\le\binom{n}{3}<1.7\cdot10^{17}\)
fits \code{long long}.
\end{proofbox}

\solstep{Implementation}
\begin{lstlisting}[caption={CF 61E full solution: two frequency-BIT sweeps around the middle index}, label={lst:c14-61-soln}]
#include <bits/stdc++.h>
using namespace std;
const int MX = 1000005;
int n, a[MX], rnk[MX], tf[MX], m;
int lg[MX], rl[MX];

int lowbit(int i){ return i & -i; }
void add(int r, int d){ for (; r <= m; r += lowbit(r)) tf[r] += d; }
int count_le(int r){
    int s = 0;
    for (; r > 0; r -= lowbit(r)) s += tf[r];
    return s;
}
int main(){
    scanf("%d", &n);
    vector<int> v(n);
    for (int i = 1; i <= n; i++){ scanf("%d", &a[i]); v[i-1] = a[i]; }
    sort(v.begin(), v.end());
    v.erase(unique(v.begin(), v.end()), v.end());
    m = (int)v.size();
    for (int i = 1; i <= n; i++)
        rnk[i] = (int)(lower_bound(v.begin(), v.end(), a[i]) - v.begin()) + 1;

    for (int i = 1; i <= n; i++){                 // greater to the LEFT
        lg[i] = (i - 1) - count_le(rnk[i]);
        add(rnk[i], 1);
    }
    for (int i = 1; i <= m; i++) tf[i] = 0;       // fresh tree
    for (int i = n; i >= 1; i--){                 // smaller to the RIGHT
        rl[i] = count_le(rnk[i] - 1);
        add(rnk[i], 1);
    }
    long long ans = 0;
    for (int i = 1; i <= n; i++) ans += (long long)lg[i] * rl[i];
    printf("%lld\n", ans);
}
\end{lstlisting}

\solstep{Trace on the Official Sample}
\begin{tracebox}[CF 61E: a=(10,8,3,1), ranks (4,3,2,1)]
Forward (greater-left): \(j{=}1\): \(0-\text{count\_le}(4)=0\).
\(j{=}2\): seen \(\{4\}\), \(1-0=1\).  \(j{=}3\): seen \(\{3,4\}\),
\(2-0=2\).  \(j{=}4\): seen \(\{2,3,4\}\), \(3-0=3\).  So
\(\mathit{lg}=(0,1,2,3)\).\\
Backward (smaller-right): \(j{=}4\): \(0\).  \(j{=}3\): seen \(\{1\}\),
\(\text{count\_le}(1)=1\).  \(j{=}2\): seen \(\{1,2\}\), \(\text{count\_le}(2)=2\).
\(j{=}1\): seen \(\{1,2,3\}\), \(\text{count\_le}(3)=3\).  So
\(\mathit{rl}=(3,2,1,0)\).\\
Products: \(0{\cdot}3+1{\cdot}2+2{\cdot}1+3{\cdot}0=0+2+2+0=4\).  Printed
\texttt{4} -- the official output.  Sanity on \texttt{3 2 1}
(ranks \(3,2,1\)): \(\mathit{lg}=(0,1,2)\), \(\mathit{rl}=(2,1,0)\),
products \(0+1+0=1\) \checkmark; \texttt{2 3 1}: \(\mathit{lg}=(0,0,2)\),
\(\mathit{rl}=(1,1,0)\), \(0+0+0=0\) \checkmark.
\end{tracebox}

\solstep{Complexity}
\begin{complexitybox}
Two sweeps of \(n\) \(\bigO{\log n}\) queries/updates plus an
\(\bigO{n\log n}\) sort for compression: total \(\bigO{n\log n}\), at
\(n=10^6\) about \(2\cdot10^7\) tree touches -- fine under the 5\,s limit.
Memory: three \code{int} arrays and one frequency tree of size \(n\), plus the
compression copy, \(\approx 6\) arrays \(\sim 24\) MB of \(4\) bytes -- inside
256 MB.  \textbf{AC}.
\end{complexitybox}

\solstep{Verification and Edge Cases}
\begin{center}\renewcommand{\arraystretch}{1.2}
\begin{tabular}{@{}p{0.20\linewidth} p{0.30\linewidth} p{0.38\linewidth}@{}}
\toprule
\textbf{Case} & \textbf{Why tricky} & \textbf{Code behaviour}\\
\midrule
strictly increasing & no triplet is decreasing & \(\mathit{lg}\equiv0\) for every \(j\) except the far right; products all \(0\) \\
strictly decreasing & every \(\binom{n}{3}\) counts & \(\mathit{lg}=(0,1,2,\dots)\), \(\mathit{rl}=(\dots,2,1,0)\), sum \(=\binom{n}{3}\) \\
\(n=3\) minimum & single candidate middle & works: only \(j=2\) can have both sides \\
answer size & \(\binom{10^6}{3}\) overflows \code{int} & \code{long long} accumulator and cast \\
\(r_j=1\) & \(\code{count\_le}(0)\) must be \(0\) & walk loop body never runs \\
\bottomrule
\end{tabular}\end{center}

\begin{hardeningbox}
\begin{itemize}[nosep,leftmargin=1.4em]
  \item \textbf{Fast I/O:} \(10^6\) integers -- \code{scanf}/\code{printf} as
        above (or \code{cin} with \code{sync\_with\_stdio(false)} and
        \code{tie(nullptr)}); \code{cout} per line would TLE on I/O alone.
  \item \textbf{Reset the tree between sweeps:} forgetting to zero
        \(\mathit{tf}\) after the forward pass corrupts \(\mathit{rl}\); this
        is the problem's signature WA.  One \(\bigO{m}\) clear is enough.
  \item \textbf{Distinctness is a licence, not a luck:} the strict
        \(\code{count\_le}(r-1)\) for ``smaller'' is correct \emph{because}
        no two powers tie.  If a variant allows equal values, fix ties by a
        stable secondary key before ranking.
  \item \textbf{\code{int} vs \code{long long}:} \(\mathit{lg},\mathit{rl}\)
        fit \code{int} (\(\le n\)); only their product summed over \(n\) needs
        \code{long long} -- cast at multiply time, \code{(long long)lg[i]*rl[i]}.
\end{itemize}
\end{hardeningbox}

% ---------------------------------------------------------------
\subsection{Common Pitfalls}
% ---------------------------------------------------------------

\begin{warningbox}[Trap 1: forgetting that lifting needs a power-of-two cap]
\(\coderef{c14-kth}\) must start \(\mathit{step}\) at the highest power of
two \(\le m\), not at \(m\).  Starting mid-block means the very first
candidate \(\mathit{tf}[\mathit{idx}+\mathit{step}]\) no longer covers a
contiguous prefix and the search returns a wrong rank that still ``looks
plausible''.  Compute \(\mathit{step}\) with the \code{while (step <= m) step
<<= 1;} then \code{step >>= 1;} idiom.
\end{warningbox}

\begin{warningbox}[Trap 2: counting \emph{positions} when you mean \emph{ranks}]
A frequency BIT's index is a value's rank, not its location.  For a ``how
many elements \(\le k\)'' query the answer is \(\code{count\_le}(\mathit{rank}(k))\),
and \(\mathit{rank}(k)\) must come from the same sorted set you compressed
against -- and \(k\) must have been \emph{collected} first
(\propref{c14-compress}).  Querying a rank that was never inserted is the
compression-off-by-one we warned about in \(\secref{c14-range}\).
\end{warningbox}

\begin{warningbox}[Trap 3: reusing the forward tree for the backward sweep]
CF~61E's two counts come from \emph{opposite} directions.  Leaving the
forward BIT populated makes \(\mathit{rl}[j]\) ``smaller to the right''
silently include left-side elements -- answers too big.  Clear \(\mathit{tf}\)
(or use a second tree).  The same trap hides in ``pairs with property''
problems that need greater-before \emph{and} less-before
(CF~1579E2): they share the tree only if you reset it.
\end{warningbox}

\begin{warningbox}[Trap 4: \code{int} product on a triplet count]
The weakness reaches \(\binom{n}{3}\).  Even one \(\mathit{lg}[j]\cdot
\mathit{rl}[j]\) product can overflow \code{int} (up to \(\sim n^2/4\));
summed, it is astronomically past \(2^{31}\).  Multiply into \code{long long}
and accumulate in \code{long long} -- CF~61E's own note that \code{\%lld}
is fine (unlike most modern judges) is a hint that this is a 64-bit answer.
\end{warningbox}

% ---------------------------------------------------------------
\subsection{Practice Problems}
% ---------------------------------------------------------------

\begin{stuckbox}
\begin{itemize}[nosep,leftmargin=1.4em]
  \item Did you compress \emph{all} values (and every query bound) before the
        sort, so no rank is ever missing?
  \item Is the lifting \(\mathit{step}\) a power of two \(\le m\)?
  \item For a two-sided count, did you reset the tree between directions?
  \item Is the accumulation \code{long long} at the multiply, not just the sum?
\end{itemize}
\end{stuckbox}

\begin{practiceset}
  \practice{Array Optimization by Stack}{CF 1579E2}
    {https://codeforces.com/problemset/problem/1579/E2}
    {CF 1800 L4}
    {Cost of a build step is the number of earlier elements \emph{greater
     than} plus those \emph{less than} the pushed value.  For each position
     \(j\) that is \((j{-}1)-\code{count\_le}(r_j)\) (greater before) plus
     \(\code{count\_le}(r_j-1)\) (less before) -- the very two prefix counts
     of \(\coderef{c14-61-soln}\), now \emph{summed} not multiplied.  One
     left-to-right frequency sweep; distinct ranks need the tie-handling note
     of the hardening box.}
  \practice{Invariant}{CF 540E}
    {https://codeforces.com/problemset/problem/540/E}
    {CF 1900 L5}
    {Two descriptions of the same permutation: count the pairs whose
     relative order the two disagree on.  Map one sequence's values to the
     other's positions, reducing it to an \emph{inversion count} of a single
     permutation -- exactly \(\coderef{c14-inv}\)'s right-to-left sweep, then
     subtract from \(\binom{n}{2}\).  The lesson: a ``difference of two
     orders'' is a frequency-BIT inversion count after a relabelling.}
\end{practiceset}

\paragraph{Back-cites and forwards.}
\secref{c14-range} (\(\coderef{c14-compress}\)) supplies the ranks this
section indexes by; \propref{c14-compress}'s ``order survives, magnitudes do
not'' is why a count-based tree is the right companion.  Chapter~1
(\chref{c01}) owns the offline merge count; Ch.~6 (\chref{c06}) owns PBDS and
the LC~315/LC~493 URLs (the same sweep, a balanced tree); Ch.~8 (\chref{c07})
owns the online stream median -- lifting here is the \emph{static/offline}
select.  Forward: the same ``sort one axis, count the other'' shape returns,
fully offline, in \secref{c14-offline}, and the lifting search is the exact
predicate the Summit CF~1181D runs per binary-search step.

% ---------------------------------------------------------------
\begin{recallbox}[Section 18.3 in 10 seconds]
Read the index as a \emph{value}: frequency BIT \(\;\code{add}(r,\pm1)\)
inserts, \(\code{count\_le}(r)\) counts \(\le\) rank \(r\), and
\(\code{count\_le}(r-1)\) is strictly-below (\(\defref{c14-freq}\)).
Select backwards with \emph{lifting} -- build the largest prefix still \(<k\)
bit by bit, answer \(\mathit{idx}+1\) (\(\thmref{c14-lift}\),
\(\coderef{c14-kth}\)).  Inversions \(=\) one right-to-left sweep
(\(\coderef{c14-inv}\)); decreasing triplets \(=\) two sweeps whose counts
multiply per middle (CF~61E, \(\probref{c14-triplets}\),
\(\coderef{c14-61-soln}\)).  Distinct ranks, reset between directions, 64-bit
answer.
\selftest{
  \item State the BIT's two interpretations (position-sum vs value-count) and
        one problem each.
  \item Why must the lifting \(\mathit{step}\) be a power of two, and what
        breaks otherwise?
  \item Express CF~61E's \(\mathit{lg}[j]\) and \(\mathit{rl}[j]\) using only
        \(\code{count\_le}\).
  \item Give the one reason the merge count of Ch.~1 cannot answer an
        interleaved ``how many now \(\le k\)'' query.
  \item Why is \(\code{count\_le}(r-1)\) ``strictly smaller'' only when the
        values are distinct?
}
\end{recallbox}

% =====================================================================
%  18.4  Fenwick on Trees: driving Ch.5's Euler array   sec:c14-trees
%  Exports: def:c14-eulerbit  prop:c14-eulerrange  prop:c14-path
%           thm:c14-parity    lst:c14-flatten  lst:c14-subtree
%           lst:c14-path  lst:c14-383-soln  fig:c14-euler  prob:c14-prop
%  Reuses (never re-derives): Ch.5 (\chref{c24}) tin/tout + subtree=interval
%    (\thmref{c24-range}); \secref{c14-core} plain BIT; \secref{c14-range}
%    difference BIT (range add / point read).
%  Running tree T (Ch.12--17), root 1, children ascending, edges
%    1-2,1-3,2-4,2-5,3-6,3-7,5-8,7-9,9-10.  Ch.16 preorder (INCLUSIVE tout):
%      tin  (by node 1..10) = (1,2,6,3,4,7,8,5,9,10)
%      tout (by node 1..10) = (10,5,10,3,5,7,10,5,10,10)
%      size=(10,4,5,1,2,1,3,1,2,1)  depth=(0,1,1,2,2,2,2,3,3,4)
%      Euler order (position -> node) = (1,2,4,5,8,3,6,7,9,10)
%      sgn (depth parity, by node)    = (+1,-1,-1,+1,+1,+1,+1,-1,-1,+1)
%  Opener six-angle row (BIT on T) reproduced in the subtree-add trace:
%    subtree(2)=[2,5] add +3, then whole tree +2: node4 reads 5, node6 reads 2.
%  CF 383C Propagating tree (2000) -- live-verified via the CF API + statement:
%    n,m<=200000, a_i<=1000, val<=1000, root 1.
%    Query "1 x val": add val to x, then -val to children, +val to grand-
%    children, ... alternating by depth through the whole subtree.
%    Query "2 x": print current value at x.
%    Sample (verified): 5 5 / 1 2 1 1 2 / edges 1-2,1-3,2-4,2-5 /
%      1 2 3 / 1 1 2 / 2 1 / 2 2 / 2 4  => output 3 / 3 / 0.
%    Solution: single difference BIT over Euler positions; store delta=val*sgn[x]
%      on range [tin[x],tout[x]]; read a[x] + sgn[x]*point(tin[x]).
% =====================================================================
\section{Fenwick on Trees: Driving the Euler Array}\label{sec:c14-trees}

\begin{previouslybox}[Ch.~5 flattened the tree; it just had no engine yet]
Chapter~5 (\chref{c24}) proved the reduction we now exploit without re-deriving
it: a depth-first \term{Euler tour} stamps each node \(v\) an entry time
\(\tin(v)\) and an exit time \(\tout(v)\) so that \emph{the subtree of \(v\) is
exactly the index interval} \([\tin(v),\tout(v)]\) (\thmref{c24-range}).  What
Ch.~5 could do was read static sums over that interval with a prefix array
(\(\propref{c24-prefix}\)); what it \emph{could not} do was keep answering
after the values changed, because a prefix array dies under updates.  This
section bolts \secref{c14-core}'s BIT onto Ch.~5's flat array, discharging the
forward pointer Ch.~5 left here: subtree ranges now get an engine.
\end{previouslybox}

\begin{topicintro}
Once the tree is flattened, every ``add to a subtree / read a node / sum a
subtree'' question stops being a \emph{tree} problem and becomes a \emph{range}
problem over the slots \(1..n\) -- and ranges over slots with an invertible
accumulator are exactly what a BIT is built for.  Concretely we place each
node's value at slot \(\tin(v)\), then reach for whichever \secref{c14-range}
flavour the query shape forces: a \emph{plain} BIT for point-update + subtree
sum, a \emph{difference} BIT for subtree-add + node-read, and the
\emph{enter\(+1\)/leave\(-1\)} re-interpretation of the difference BIT for
root-path reads.  The single new idea -- used to crack the solved problem -- is
that a sign that alternates with \emph{depth} becomes a \(\pm1\) multiplier
attached to each slot, so a depth-propagating update is still \emph{one} range
addition.

\textbf{Prerequisites.} Ch.~5 (\chref{c24}) for \(\tin/\tout\);
\secref{c14-core} for the two walks; \secref{c14-range} for the difference BIT
(range add / point read).  The running example is the tree \(T\) with the
stamps recorded in \figref{c14-array}'s opener block.

\textbf{After this section you can}
\begin{itemize}[nosep]
  \item turn any subtree update/query into a range operation on the Euler
        array and pick the right BIT flavour; \hfill\lv{3}
  \item answer root-path reads with the two-event \(\mathit{enter}{+}1 /
        \mathit{leave}{-}1\) trick, no LCA machinery; \hfill\lv{4}
  \item absorb a depth-parity (or any per-node \(\pm1\)) into the stored
        value so a propagating update stays a single range add; \hfill\lv{4}--\lv{5}
  \item solve CF~383C with one difference BIT over the tour. \hfill\lv{5}
\end{itemize}
\end{topicintro}

\begin{recognitionbox}[Pattern card -- ``many subtree / root-path operations on one static tree'']
\reckey{``add \(x\) to every node in the subtree of \(v\)'' interleaved with
``read node \(v\)'`; ``invert the colour of a whole subtree, query one
vertex''; ``sum of the values on the path from the root to \(v\)''; ``the
value at \(v\) after updates that propagate to descendants with alternating
sign''; ``how many marked nodes are ancestors of \(v\)''.  The tree itself
never changes -- only per-node numbers do, and queries cluster on
subtrees/paths.}
\recuse{
  \item subtree add + point read (or point add + subtree sum) -- flatten to
        \([\tin,\tout]\), difference BIT (this section, CF~383C, CSES~1137);
  \item root-path read with point ``activation'' -- enter\(+1\)/leave\(-1\)
        on the same difference BIT (\(\propref{c14-path}\));
  \item updates whose effect depends on depth parity or a per-node sign --
        fold the sign into the stored delta (\(\thmref{c14-parity}\)).
}
\recnotuse{
  \item the tree is \emph{cut, linked or rerooted} between queries --
        Euler-tour trees / splay, Ch.~6 (\chref{c06}); a static tour array
        cannot survive a structural change;
  \item the range aggregate is not invertible (max/min/gcd over the subtree)
        -- segment tree with the flat array, Ch.~19 (\chref{c13});
  \item paths are \emph{arbitrary} \(u\)-to-\(v\) with non-commutative data
        -- heavy-light decomposition on top of a segment tree, not a BIT;
  \item the tree is a single path -- it is just an array; use
        \secref{c14-range} directly.
}
One-line test: \emph{is the tree static and does every operation address a
subtree or a root-path?  Flatten with Ch.~5's \(\tin/\tout\), then let a BIT
own the range.}
\end{recognitionbox}

% ---------------------------------------------------------------
\subsection{The Flattened Array and Its Two Range Shapes}
% ---------------------------------------------------------------

\begin{definition}{Euler-indexed BIT}{c14-eulerbit}
Fix a DFS order of the tree and let \(\mathit{euler}[p]\) be the node stamped
\(\tin(\mathit{euler}[p])=p\).  An \term{Euler-indexed BIT} is a BIT
\(B[1..n]\) whose slot \(p\) stores the aggregate of the \emph{node}
\(\mathit{euler}[p]\).  A node operation on \(v\) becomes a slot operation at
\(\tin(v)\); a subtree operation on \(v\) becomes a range operation on
\([\tin(v),\tout(v)]\) (both \(1\)-based, \(\tout\) \emph{inclusive}, per this
chapter's convention).  Because \(\tin\) is a permutation, ``node \(v\)'' and
``slot \(\tin(v)\)'' name the same object -- we translate silently between
them.
\end{definition}

\begin{figure}[H]
\centering
\begin{tikzpicture}[diagram, scale=0.62, transform shape,
    tnode/.style={circle, draw=primaryblue, line width=0.8pt, fill=primaryblue!8,
                  minimum size=6.5mm, font=\scriptsize\bfseries, text=primaryblue!55!black}]
  % tree edges (root 1 at top), children ascending left-to-right
  \node[tnode] (n1) at (0,3)   {1};
  \node[tnode] (n2) at (-3,2)  {2};
  \node[tnode] (n3) at (3,2)   {3};
  \node[tnode] (n4) at (-4,1)  {4};
  \node[tnode] (n5) at (-1.5,1){5};
  \node[tnode] (n6) at (2,1)   {6};
  \node[tnode] (n7) at (4.2,1) {7};
  \node[tnode] (n8) at (-1.5,0){8};
  \node[tnode] (n9) at (4.2,0) {9};
  \node[tnode] (n10) at (4.2,-1){10};
  \draw (n1)--(n2) (n1)--(n3) (n2)--(n4) (n2)--(n5) (n3)--(n6) (n3)--(n7)
        (n5)--(n8) (n7)--(n9) (n9)--(n10);
  % tin/tout annotations
  \node[annot, font=\tiny, anchor=south] at (n1.north) {\(1[1,\!10]\)};
  \node[annot, font=\tiny, anchor=east]  at (n2.west)  {\(2[2,\!5]\)};
  \node[annot, font=\tiny, anchor=east]  at (n3.west)  {\(3[6,\!10]\)};
  \node[annot, font=\tiny, anchor=south] at (n4.north) {\(4[3]\)};
  \node[annot, font=\tiny, anchor=north] at (n5.south) {\(5[4,\!5]\)};
  \node[annot, font=\tiny, anchor=south] at (n6.north) {\(6[7]\)};
  \node[annot, font=\tiny, anchor=south] at (n7.north) {\(7[8,\!10]\)};
  \node[annot, font=\tiny, anchor=east]  at (n8.west)  {\(8[5]\)};
  \node[annot, font=\tiny, anchor=east]  at (n9.west)  {\(9[9,\!10]\)};
  \node[annot, font=\tiny, anchor=west]  at (n10.east) {\(10[10]\)};
  % Euler slot strip below the tree
  \foreach \p/\nd in {1/1,2/2,3/4,4/5,5/8,6/3,7/6,8/7,9/9,10/10} {
    \draw[difforange, line width=0.6pt, fill=difforange!8] (\p*0.62,-2.6) rectangle (\p*0.62+0.62,-2.0);
    \node[font=\tiny\ttfamily, text=orange!55!black] at (\p*0.62+0.31,-2.3) {\nd};
    \node[font=\tiny\ttfamily, text=darkgray] at (\p*0.62+0.31,-1.82) {\p};}
  \node[font=\scriptsize, text=orange!55!black, anchor=east] at (0.4,-2.3) {euler:};
  % subtree(2)=[2,5] and subtree(3)=[6,10] interval bars
  \fill[successgreen!35] (2*0.62+0.05,-3.0) rectangle (5*0.62+0.57,-2.75);
  \node[font=\tiny, text=green!45!black] at (3.5*0.62+0.3,-3.15) {subtree(2)=[2,5]};
  \fill[softpurple!28] (6*0.62+0.05,-3.0) rectangle (10*0.62+0.57,-2.75);
  \node[font=\tiny, text=purple!55!black] at (8*0.62+0.3,-3.15) {subtree(3)=[6,10]};
\end{tikzpicture}
\caption{Tree \(T\) with each node annotated \emph{v}[\(\tin(v),\tout(v)\)]
(Ch.~5's tour, \(\tout\) inclusive).  The orange strip is the Euler array:
slot~\(p\) holds node \(\mathit{euler}[p]\), read as \((1,2,4,5,8,3,6,7,9,10)\).
The two shaded bars are the intervals a BIT is handed to answer ``sum/add
subtree(2)'' and ``sum/add subtree(3)'' -- the whole tree collapses to
contiguous slots.  Node~3 sits at slot~6, so it is \emph{outside} the
subtree(2) bar even though it is drawn beside node~2 in the picture.}
\label{fig:c14-euler}
\end{figure}

Two operation shapes cover almost every tree-BIT problem.  Both are just a
flavour from \secref{c14-core}/\secref{c14-range} wearing the \(\tin\)
translation:

\begin{proposition}{Subtree sum with point updates; subtree add with point reads}{c14-eulerrange}
Let \(B\) be a BIT over the Euler array.
\begin{itemize}[nosep,leftmargin=1.4em]
  \item \emph{Point update of a node, subtree sum query} (plain BIT,
        \secref{c14-core}): to add \(\delta\) to node \(v\), run
        \(\code{update}(\tin(v),\delta)\); the sum of the whole subtree of
        \(u\) is \(\code{query}(\tout(u))-\code{query}(\tin(u)-1)\).
        This is CSES~1137 (``Subtree Queries''), solved in Ch.~5 with a static
        prefix array; here updates are free.
  \item \emph{Subtree add, single-node read} (difference BIT,
        \secref{c14-range}): to add \(\delta\) to every node of
        \(\subtree(v)\), run \(\code{range\_add}(\tin(v),\tout(v),\delta)\);
        the current value at node \(u\) is the point read
        \(\code{point}(\tin(u))\).  Cost \(\bigO{\log n}\) each.
\end{itemize}
\end{proposition}

\begin{lstlisting}[caption={Euler-indexed BIT: subtree sum with point updates, and subtree add with point reads}, label={lst:c14-subtree}]
// Ch.5 produced tin[1..n], tout[1..n]; we never recompute them here.
int lowbit(int i){ return i & -i; }

// (A) PLAIN BIT: add delta to node v, read the sum over subtree(u).
void bit_add(long long B[], int i, long long d){ for(; i<=n; i+=lowbit(i)) B[i]+=d; }
long long bit_pre(long long B[], int i){
    long long s=0; for(; i>0; i-=lowbit(i)) s+=B[i]; return s; }
void update_node(int v, long long d){ bit_add(B, tin[v], d); }          // 18.1
long long subtree_sum(int u){                                          // 18.1
    return bit_pre(B, tout[u]) - bit_pre(B, tin[u]-1); }

// (B) DIFFERENCE BIT: add d to every node of subtree(v), read one node.
void diff_add(long long D[], int i, long long d){ for(; i<=n; i+=lowbit(i)) D[i]+=d; }
void subtree_add(int v, long long d){        // range [tin, tout] = subtree(v)
    diff_add(D, tin[v], d); diff_add(D, tout[v]+1, -d); }   // 18.2 (i>n is a no-op)
long long node_value(int u){                 // point read at slot tin[u]
    long long s=0; for(int i=tin[u]; i>0; i-=lowbit(i)) s+=D[i]; return s; }
\end{lstlisting}

\begin{tracebox}[Subtree add on \(T\): \(\subtree(2){+}3\), then whole tree \(+2\)]
All nodes start at \(0\).  \(\subtree(2)=[\tin_2,\tout_2]=[2,5]\); the whole
tree is \(\subtree(1)=[1,10]\).  Two difference-BIT range adds:
\(\text{diff\_add}(2,+3),\text{diff\_add}(6,-3)\) then
\(\text{diff\_add}(1,+2),\text{diff\_add}(11,-2)\) (the last is \(i>n\), a
no-op).  Reading \(\code{node\_value}\) as a prefix sum over the difference
array at each node's \(\tin\):
\[
\renewcommand{\arraystretch}{1.1}
\begin{array}{c|c|c|c|c}
\text{node }u & \tin(u) & \tin\in[2,5]? & \tin\in[1,10]? & \text{value}\\ \hline
4 & 3 & \checkmark\ (+3) & \checkmark\ (+2) & \mathbf{5}\\
6 & 7 & \times & \checkmark\ (+2) & \mathbf{2}\\
2 & 2 & \checkmark\ (+3) & \checkmark\ (+2) & 5\\
3 & 6 & \times & \checkmark\ (+2) & 2
\end{array}
\]
Node~4 (a child of~2) carries both adds \(\to5\); node~6 (under~3) misses the
subtree(2) add \(\to2\).  This is exactly the opener's \textsc{BIT on \(T\)}
row -- and \figref{c14-euler}'s interval bars make it obvious without any
walking.
\end{tracebox}

\begin{complexitybox}[Why a BIT and not a segment tree here]
Both shapes of \(\propref{c14-eulerrange}\) are \(\bigO{\log n}\) with a
\emph{single} array and four-line code.  A segment tree does the same in the
same asymptotics but carries lazy tags and a \(4n\) array; the BIT is smaller,
faster, and cache-friendly.  The BIT is the right tool \emph{because} the
aggregate is a sum (invertible).  The moment the subtree query wants a
non-invertible combine -- ``maximum value in the subtree'' or ``first node
passing a threshold with arbitrary decrements'' -- the flat interval no longer
supports the difference trick and you owe the range to Ch.~19's segment tree
(\chref{c13}).  That boundary is stated here, not coded here.
\end{complexitybox}

% ---------------------------------------------------------------
\subsection{Root-Path Reads: the Enter\(	ext{+}1\) / Leave\(	ext{-}1\) Trick}
% ---------------------------------------------------------------

The difference BIT already stores a range add as a \(+\delta\) entering the
interval and a \(-\delta\) leaving it.  Read that backwards: if we \emph{tag}
a node \(v\) with a value \(x\) by writing \(+x\) at \(\tin(v)\) and \(-x\) at
\(\tout(v){+}1\), then a point read at \(\tin(u)\) sums \(x\) over exactly the
tags whose interval contains \(\tin(u)\) -- and \(\tin(v)\le\tin(u)\le\tout(v)\)
holds \emph{iff} \(v\) is an ancestor of \(u\) (Ch.~5, \(\thmref{c24-range}\)).
So the point read is the \emph{root-path} sum.

\begin{proposition}{Point-read at \(\tin(u)\) equals the root-to-\(u\) path sum}{c14-path}
Maintain a difference BIT \(D\).  To ``activate'' node \(v\) with weight \(x\)
for all its descendants, do \(D[\tin(v)]{+}{=}x,\;D[\tout(v){+}1]{-}{=}x\).
Then \(\sum_{p\le\tin(u)} D[p]\) (the point read at node \(u\)) equals
\(\sum_{v \text{ ancestor of } u} x_v\): the running prefix at slot
\(\tin(u)\) is inside \(v\)'s live window precisely when \(v\) is on the path
from the root to \(u\).  Each activation and each path read is \(\bigO{\log n}\).
\end{proposition}

\begin{lstlisting}[caption={two-event activation: subtree-add semantics become root-path reads}, label={lst:c14-path}]
// +x on ENTERING v's window, -x on LEAVING it -- one range add over [tin, tout].
void activate(int v, long long x){
    diff_add(D, tin[v],     x);        // switch x ON for v's whole subtree
    diff_add(D, tout[v] + 1, -x);      // switch it OFF after the subtree
}
// point read at tin[u] = sum of x over every ancestor v of u (root path of u):
long long root_path_sum(int u){
    long long s=0; for(int i=tin[u]; i>0; i-=lowbit(i)) s+=D[i]; return s; }
// So "add x to subtree(v); read value at u" and
//    "activate v with x; read the root-to-u sum" are the SAME two numbers
//    (x, -x) at (tin(v), tout(v)+1) -- just read with opposite intent.
\end{lstlisting}

\begin{insightbox}[One pair of writes, two stories]
The lines \(\mathit{diff\_add}(D,\tin(v),x)\) and
\(\mathit{diff\_add}(D,\tout(v){+}1,-x)\) are the same whether you call them
``add \(x\) to subtree \(v\)'' or ``activate \(v\) for its descendants''.  The
difference is only \emph{what you query}: at the same \(\tin(u)\), a subtree
add sees \(x\) when \(u\) is a \emph{descendant} of the source; a root-path
read sees \(x\) when the source is an \emph{ancestor} of \(u\).  Ancestor and
descendant are one relation read from opposite ends -- so one difference BIT
serves both, and this is the whole trick behind CF~383C below.
\end{insightbox}

% ---------------------------------------------------------------
\subsection{Folding a Depth Sign into the Stored Value}
% ---------------------------------------------------------------

\begin{theorem}{Alternating-by-depth becomes a \(\pm1\) multiplier}{c14-parity}
Root the tree; let \(\operatorname{sgn}(v)=+1\) if \(\depth(v)\) is even and
\(-1\) otherwise.  An update ``add \(x\) at node \(r\), and let it propagate
with a sign flip at every edge'' gives node \(u\in\subtree(r)\) the
contribution
\[
  x\cdot(-1)^{\depth(u)-\depth(r)} \;=\; x\cdot\operatorname{sgn}(u)\cdot\operatorname{sgn}(r).
\]
Hence if we store, for each update, the single number \(\Delta=x\cdot\operatorname{sgn}(r)\)
as a \emph{plain} subtree-range add on \([\tin(r),\tout(r)]\), then reading a
node is \(\operatorname{sgn}(u)\cdot(\text{point read at }\tin(u))\), which
reproduces the alternating contribution exactly.  The sign pattern costs
nothing but one \(\pm1\) multiply on each end.
\end{theorem}
\begin{proofbox}
Since \((-1)^{a-b}=(-1)^{a+b}=(-1)^a(-1)^b\) for integers, the propagated
sign of update \(r\) at node \(u\) is \(\operatorname{sgn}(u)\operatorname{sgn}(r)\).
The range add stores \(\Delta=x\operatorname{sgn}(r)\) so that \emph{every}
slot in \([\tin(r),\tout(r)]\) (i.e.\ every descendant of \(r\)) accumulates
\(\Delta\); a point read at \(u\) is the sum \(\sum_{r\,\text{anc. of}\,u}
x_r\operatorname{sgn}(r)\) over exactly the updates whose subtree contains
\(u\).  Multiplying by \(\operatorname{sgn}(u)\) gives
\(\sum_r x_r\operatorname{sgn}(u)\operatorname{sgn}(r)=\sum_r x_r(-1)^{\depth(u)-\depth(r)}\),
the definition of the propagating process.  Nodes outside \(\subtree(r)\) never
see \(\Delta\) because their \(\tin\) is outside its range.
\end{proofbox}

\noindent\textbf{In plain words.}  Do not simulate the up-down-up propagation.
Precompute each node's \(\operatorname{sgn}\) once in the same DFS that stamps
\(\tin/\tout\), then treat ``add to a propagating subtree'' as \emph{one}
difference-BIT range add of a signed amount.  \(\thmref{c14-parity}\) is what
turns CF~383C from a scary statement into \(\coderef{c14-subtree}\).

\begin{tracebox}[Parity store on \(T\): ``add 4 at node 3'' propagates]
\(\depth(3)=1\) so \(\operatorname{sgn}(3)=-1\); \(\Delta=4\cdot(-1)=-4\) added
to range \([\tin_3,\tout_3]=[6,10]\).  Read \(u\) as
\(\operatorname{sgn}(u)\cdot(\text{point at }\tin(u))\), point over the range
\(=-4\) for all of \(\subtree(3)=\{3,6,7,9,10\}\):
\[
\renewcommand{\arraystretch}{1.1}
\begin{array}{c|c|c|c|c}
u & \depth & \operatorname{sgn}(u) & \text{point} & \text{contrib }=\operatorname{sgn}(u)\cdot(-4)\\ \hline
3 & 1 & -1 & -4 & +4\ (\text{self})\\
6 & 2 & +1 & -4 & -4\ (\text{child})\\
7 & 2 & +1 & -4 & -4\\
9 & 3 & -1 & -4 & +4\ (\text{grandchild})\\
10& 4 & +1 & -4 & -4
\end{array}
\]
Self \(+4\), children \(-4\), grandchildren \(+4\), great-grandchild \(-4\) --
the alternating propagation of the statement, from a single range add and a
read-time sign (\(\thmref{c14-parity}\)).
\end{tracebox}

% ---------------------------------------------------------------
\subsection{Problem Lab: CF 383C Propagating Tree}
% ---------------------------------------------------------------

\solvedproblem{c14-prop}{Propagating tree}{CF 383C}
  {https://codeforces.com/problemset/problem/383/C}{CF 2000 L5}
  {Euler tour + one difference BIT, depth-parity sign folded into the stored delta}

\begin{problemstatement}[CF 383C Propagating tree]
A tree of \(n\) nodes numbered \(1..n\), root~\(1\); node \(i\) has value
\(a_i\).  Its \emph{propagating} rule: when a value \(\mathit{val}\) is added
to node \(i\), \(-\mathit{val}\) is added to every child of \(i\); and when
\(-\mathit{val}\) reaches a child, \(+\mathit{val}\) reaches that child's
children, and so on down the whole subtree -- the sign flips every level.
Two query types: \(\texttt{1 }x\,\mathit{val}\) performs the propagating
addition starting at \(x\); \(\texttt{2 }x\) prints the current value at \(x\).

\textbf{Input.} \(n,m\le 2\cdot10^5\); then \(a_1..a_n\) (\(1\le a_i\le1000\));
then \(n-1\) edges; then \(m\) queries with \(1\le x\le n\), \(1\le\mathit{val}\le1000\).

\textbf{Output.} For each type~2 query, one line: the value at \(x\).

\textbf{Example.}
\begin{verbatim}
Input:
5 5
1 2 1 1 2
1 2
1 3
2 4
2 5
1 2 3
1 1 2
2 1
2 2
2 4

Output:
3
3
0
\end{verbatim}
\end{problemstatement}

\begin{attemptgate}[Attempt it yourself first -- \lv{5}, give yourself 30 minutes]
``Add to a subtree with a sign that flips by level'' sounds like you must
walk the subtree.  But the tree is static, so the subtree is a flat interval
(\(\thmref{c24-range}\)) and the sign is a fixed \(\pm1\) per node
(\(\thmref{c14-parity}\)).  Can you make the update \emph{one} range addition?
\peek{store \(\mathit{val}\cdot\!\operatorname{sgn}(x)\) on the interval and read with a sign.}
\nudge{a difference BIT over Euler slots handles ``add to subtree, read one
node'' (\(\coderef{c14-subtree}\) shape B).  All that is left is the parity
factor of \(\thmref{c14-parity}\).}
\fullhint{Precompute \(\tin,\tout,\operatorname{sgn}\) in one DFS.  Query 1
\((x,\mathit{val})\): \(\code{subtree\_add}(x,\;\mathit{val}\cdot\operatorname{sgn}(x))\).
Query 2 \((x)\): print \(a_x+\operatorname{sgn}(x)\cdot\code{node\_value}(x)\).
Both \(\bigO{\log n}\); use \code{long long} for the accumulated values
(\(m\cdot1000\) drifts past \code{int}).}
\end{attemptgate}

\solstep{Understanding}
Two facts collapse the problem.  (1) The tree never changes, so ``the subtree
of \(x\)'' is the fixed index interval \([\tin(x),\tout(x)]\) -- Ch.~5's
\(\thmref{c24-range}\), which we use as a black box.  (2) The propagating sign
is not a process to simulate: it is \(+1/-1\) determined by the parity of the
\emph{depth difference}, which \(\thmref{c14-parity}\) factors into
\(\operatorname{sgn}(u)\operatorname{sgn}(x)\).  Once factored, the update is a
range addition of a signed constant and the read multiplies by a per-node
constant -- \(\propref{c14-eulerrange}\) shape B with one extra \(\pm1\).

\solstep{Key Observations}
\begin{observation}
A type~1 query touches only \(\subtree(x)\); in Euler order that is a
\emph{contiguous} range, so the difference BIT's two writes
(\(+\Delta\) at \(\tin(x)\), \(-\Delta\) at \(\tout(x)+1\)) do the whole
update in \(\bigO{\log n}\).
\end{observation}
\begin{observation}
The node value never needs to be materialised: the answer to a type~2 query is
\(a_x\) (frozen) plus the propagated total, which is
\(\operatorname{sgn}(x)\) times a single point read.
\end{observation}
\begin{observation}
\(m\le2\cdot10^5\) additions of \(\mathit{val}\le1000\) can move a value by
\(\pm2\cdot10^8\), and with alternating signs the accumulator can be negative
-- \code{long long} for the BIT, and signed arithmetic throughout.
\end{observation}

\solstep{Think It Through With Me}
\begin{thinkbox}
\deadend{Recolour the subtree level by level on each query}
{A subtree can be \(\Theta(n)\) and there are \(\Theta(n)\) queries --
\(\bigO{n^2}\) worst case at \(n=2\cdot10^5\), far past the 2\,s budget.  The
propagation is a red herring: it is a \emph{formula}, not a \emph{process}.}
\deadend{Keep two segment trees, one per depth parity}
{Works, but you pay a \(4n\) tree twice and lazy tags for a task that is just
``range add, point read''.  That is the difference BIT's signature job
(\secref{c14-range}); a segment tree is overkill and easy to get wrong.}
\aha{One difference BIT plus a \(\pm1\) sign array does everything: store
\(\mathit{val}\cdot\operatorname{sgn}(x)\) on the subtree interval, and read
\(a_x+\operatorname{sgn}(x)\cdot\text{point}(\tin(x))\).}
\end{thinkbox}

\solstep{The Approach Ladder}
\begin{center}\renewcommand{\arraystretch}{1.25}
\begin{tabular}{@{}c p{0.245\linewidth} p{0.29\linewidth} l c@{}}
\toprule
\textbf{\#} & \textbf{Approach} & \textbf{Idea} & \textbf{Time} & \textbf{Verdict}\\
\midrule
1 & \lv{2}\ BFS the subtree & apply the flip edge by edge & \(\bigO{nm}\) & \tle \\
2 & \lv{3}\ two segment trees & even/odd depth, lazy range add & \(\bigO{n\log n}\) & heavy \\
3 & \lv{5}\ \textbf{difference BIT + parity} & one tree, store \(\mathit{val}\cdot\operatorname{sgn}\) (\(\thmref{c14-parity}\)) & \(\bigO{n\log n}\) & \ac \\
\bottomrule
\end{tabular}\end{center}
Rung~3 is \(\coderef{c14-subtree}\)'s shape B plus \(\thmref{c14-parity}\) --
no new data structure, one extra multiply.

\solstep{Solution}
One DFS fills \(\tin,\tout\) (Ch.~5) and, in the same walk, \(\operatorname{sgn}(v)=(-1)^{\depth(v)}\).
A single difference BIT \(D\) over slots \(1..n\).  Query 1\((x,\mathit{val})\):
\(\code{subtree\_add}(x,\mathit{val}\cdot\operatorname{sgn}(x))\).  Query
2\((x)\): print \(a_x+\operatorname{sgn}(x)\cdot\code{node\_value}(x)\).

\solstep{Proof of Correctness}
\begin{proofbox}
\emph{Range coverage:} \(\code{subtree\_add}\) writes \(+\Delta\) at
\(\tin(x)\) and \(-\Delta\) at \(\tout(x)+1\); its point read at \(u\) equals
\(\Delta\) iff \(\tin(x)\le\tin(u)\le\tout(x)\), i.e.\ iff \(u\in\subtree(x)\)
(\(\thmref{c24-range}\)), and \(0\) otherwise -- matching ``propagation only
reaches descendants''.  \emph{Sign:} the stored delta is
\(\Delta=\mathit{val}\operatorname{sgn}(x)\); the read multiplies by
\(\operatorname{sgn}(u)\), contributing
\(\mathit{val}\operatorname{sgn}(u)\operatorname{sgn}(x)=\mathit{val}(-1)^{\depth(u)-\depth(x)}\),
exactly the level-by-level flip after \(\depth(u)-\depth(x)\) edges
(\(\thmref{c14-parity}\)).  \emph{Superposition:} each query~1 is linear and
independent, so a node's value is \(a_x\) plus the sum of its contributions,
which the prefix sum accumulates.  \emph{Bounds:} \(\bigO{\log n}\) per query,
\(\bigO{n\log n}\) total with the sort-free DFS; \code{long long} is sound
since each value stays within \(\pm(1000\cdot2\cdot10^5+1000)\).
\end{proofbox}

\solstep{Implementation}
\begin{lstlisting}[caption={CF 383C: Euler tour + one difference BIT with a depth-parity sign}, label={lst:c14-383-soln}]
#include <bits/stdc++.h>
using namespace std;
const int MX = 200005;
int n, m, tin[MX], tout[MX], sgn[MX], tmr_;
long long a[MX], D[MX];
vector<int> adj[MX];

void dfs(int u, int p, int dep){                 // Ch.5 tour + parity, one pass
    tin[u] = ++tmr_;  sgn[u] = (dep & 1) ? -1 : 1;
    for (int v : adj[u]) if (v != p) dfs(v, u, dep + 1);
    tout[u] = tmr_;                              // inclusive
}
int lowbit(int i){ return i & -i; }
void bitadd(long long t[], int i, long long d){ for(; i<=n; i+=lowbit(i)) t[i]+=d; }
void range_add(int l, int r, long long d){ bitadd(D,l,d); bitadd(D,r+1,-d); } // 18.2
long long point(int i){ long long s=0; for(; i>0; i-=lowbit(i)) s+=D[i]; return s; }

int main(){
    scanf("%d%d", &n, &m);
    for (int i=1;i<=n;i++) scanf("%lld", &a[i]);
    for (int i=1;i<n;i++){ int u,v; scanf("%d%d",&u,&v);
        adj[u].push_back(v); adj[v].push_back(u); }
    dfs(1, 0, 0);
    while (m--){
        int type, x; scanf("%d%d", &type, &x);
        if (type==1){ long long val; scanf("%lld", &val);
            range_add(tin[x], tout[x], val * sgn[x]);    // signed subtree add
        } else {
            printf("%lld\n", a[x] + sgn[x] * point(tin[x]));
        }
    }
}
\end{lstlisting}

\solstep{Trace on the Official Sample}
\begin{tracebox}[CF 383C: tree \(1\!\to\!2,1\!\to\!3,2\!\to\!4,2\!\to\!5\)]
Root~1, so \(\depth=(0,1,1,2,2)\) and \(\operatorname{sgn}=(+1,-1,-1,+1,+1)\)
for nodes \(1..5\).  \(\tin=(1,2,5,3,4)\), \(\tout=(5,4,5,3,4)\) (subtree\(2=[2,4]\)).
\(a=(1,2,1,1,2)\).

\emph{Q1 \texttt{1 2 3}:} \(\Delta=3\cdot\operatorname{sgn}(2)=-3\) added to
\([2,4]\): \(D\mathrel{+}=(-3\) at~2, \(+3\) at~5\()\).
\emph{Q2 \texttt{1 1 2}:} \(\Delta=2\cdot(+1)=+2\) added to \([1,5]\):
\(D\mathrel{+}=(+2\) at~1, \(-2\) at~6 no-op\()\).
Now the difference array (slots \(1..5\)) is \((+2,-3,0,0,+3)\); prefix sums
\(=(2,-1,-1,-1,2)\).

\emph{Q3 \texttt{2 1}:} point\((\tin_1{=}1)=2\); \(a_1+\operatorname{sgn}(1)\cdot2=1+2=\mathbf{3}\).\\
\emph{Q4 \texttt{2 2}:} point\((\tin_2{=}2)=-1\); \(a_2+\operatorname{sgn}(2)(-1)=2+1=\mathbf{3}\).\\
\emph{Q5 \texttt{2 4}:} point\((\tin_4{=}3)=-1\); \(a_4+\operatorname{sgn}(4)(-1)=1-1=\mathbf{0}\).\\
Printed \(3,3,0\) -- the official output \checkmark.  (Node~5 reads \(a_5+1\cdot(-1)=1\),
node~3 reads \(a_3+(-1)\cdot2=-1\), matching the statement's worked note.)
\end{tracebox}

\solstep{Complexity}
\begin{complexitybox}
One DFS (\(\bigO{n}\)) to fill \(\tin/\tout/\operatorname{sgn}\), then \(m\)
queries at \(\bigO{\log n}\) each: total \(\bigO{(n+m)\log n}\) with
\(n,m\le2\cdot10^5\) -- roughly \(4\cdot10^6\) tree touches, easily inside 2\,s.
Memory: two \(O(n)\) \code{long long} arrays, two \(O(n)\) \code{int} arrays,
and the adjacency list \(2(n-1)\) -- about \(20\) MB, well under 256 MB.
\textbf{AC}.
\end{complexitybox}

\solstep{Verification and Edge Cases}
\begin{center}\renewcommand{\arraystretch}{1.2}
\begin{tabular}{@{}p{0.20\linewidth} p{0.30\linewidth} p{0.38\linewidth}@{}}
\toprule
\textbf{Case} & \textbf{Why tricky} & \textbf{Code behaviour}\\
\midrule
update at a leaf & subtree is one slot & \(\tin=\tout\), range add hits one point \\
\(\tout(x)=n\) & write at \(\tout+1=n+1\) & \code{bitadd} guard \(i\le n\) drops it safely \\
root update & whole array flips sign & \(\operatorname{sgn}\) already encodes it, no special case \\
value goes negative & alternating adds & \code{long long} and signed read \(a_x+\operatorname{sgn}\cdot\text{pt}\) \\
deep chain \(n=2\cdot10^5\) & recursion depth & raise the stack or DFS iteratively \\
\bottomrule
\end{tabular}\end{center}

\begin{hardeningbox}
\begin{itemize}[nosep,leftmargin=1.4em]
  \item \textbf{One-indexed everything:} \(\tin\in[1,n]\), BIT loop
        \(i\le n\), and \(\tout\) \emph{inclusive}.  If you port a \(0\)-based
        Ch.~5 tour (its \(\tout\) is half-open), \emph{shift by one} before
        touching this BIT -- mixing the two conventions is the classic silent
        WA on subtree ranges.
  \item \textbf{Recursion depth:} a chain tree sends \code{dfs} \(2\cdot10^5\)
        frames deep.  Use an explicit stack, or the compiler flag
        \texttt{-Wl,-stack,}, or compile with \code{int} recursion limit raised;
        a stack overflow shows as RE, not TLE.
  \item \textbf{\(\operatorname{sgn}\) is \(\pm1\), not the depth:} the bug is
        writing \(\mathit{val}\cdot\depth\) instead of \(\mathit{val}\cdot\operatorname{sgn}\).
        Verify on the root: its own contribution must be \(+\mathit{val}\).
  \item \textbf{Fast I/O:} \(m\le2\cdot10^5\) lines of output -- buffer it;
        \code{printf}/\code{scanf} as shown, or untied \code{cin}/\code{cout}.
\end{itemize}
\end{hardeningbox}

% ---------------------------------------------------------------
\subsection{Common Pitfalls}
% ---------------------------------------------------------------

\begin{warningbox}[Trap 1: mixing Ch.~5's \(0\)-based tour with this \(1\)-based BIT]
Ch.~5 (\chref{c24}) builds \(\tin\in[0,n)\) with a \emph{half-open}
\(\tout=\tin+\mathit{size}\); this chapter's BIT uses \(\tin\in[1,n]\) and an
\emph{inclusive} \(\tout=\tin+\mathit{size}-1\).  Feeding the first into
\(\code{query}(\tout)-\code{query}(\tin-1)\) is off by one on \emph{both}
ends: subtree sums come out one node too big, and slot~0 indexes past the
array.  Pick one convention per program and translate in exactly one place.
\end{warningbox}

\begin{warningbox}[Trap 2: treating a root-path read as a subtree read]
\(\propref{c14-path}\): the point read at \(\tin(u)\) sums updates at
\(u\)'s \emph{ancestors}, while \(\propref{c14-eulerrange}\)-B's subtree add is
seen by \(u\)'s \emph{descendants}.  They share the two writes but answer
opposite questions.  Getting ``number of active ancestors of \(u\)'' with a
subtree-add semantics (instead of the enter/leave activation) returns the
descendant count -- right number, wrong nodes.
\end{warningbox}

\begin{warningbox}[Trap 3: simulating the propagation]
CF~383C's statement reads like a process.  Do not BFS the subtree flipping
signs; \(\thmref{c14-parity}\) already collapsed it to one signed range add.
The tell that you fell into this trap: your solution has a ``level'' counter
inside the update loop and passes samples but TLEs on a chain.
\end{warningbox}

\begin{warningbox}[Trap 4: forgetting the read-time \(\operatorname{sgn}\)]
Both ends of \(\thmref{c14-parity}\) matter: multiply the stored delta by
\(\operatorname{sgn}(x)\) \emph{and} the read by \(\operatorname{sgn}(u)\).
Dropping the read-time factor makes same-parity and opposite-parity nodes
return values of the wrong sign -- samples with a shallow tree (depth 1--2)
often still pass, hiding the bug until a hidden test.
\end{warningbox}

% ---------------------------------------------------------------
\subsection{Practice Problems}
% ---------------------------------------------------------------

\begin{stuckbox}
\begin{itemize}[nosep,leftmargin=1.4em]
  \item Did you get \(\tin,\tout\) from Ch.~5 by label and \emph{not} re-derive
        them, and are they in the same (1-based, inclusive) convention as the BIT?
  \item Is your update a single range add, or are you secretly walking the
        subtree?
  \item For a sign-propagating problem, is \(\operatorname{sgn}\) computed once
        in the DFS and applied at \emph{both} the write and the read?
  \item Is \(\tout+1\) handled when \(\tout=n\) (the guard, not a crash)?
\end{itemize}
\end{stuckbox}

\begin{practiceset}
  \practice{Danil and a Part-time Job}{CF 877E}
    {https://codeforces.com/problemset/problem/877/E}
    {CF 2000 L4}
    {``Invert the colour of the whole subtree of \(v\); read the colour of
     \(u\)''.  Colour is a bit, and flipping a subtree is a range add
     \(\bmod 2\): flatten with Ch.~5 and keep a difference BIT over
     \([\tin(v),\tout(v)]\), reading \(\code{node\_value}(u)\bmod 2\) (the
     \(\oplus\) running parity).  Exactly \(\coderef{c14-subtree}\) shape B,
     with the value taken mod~2 -- no segment tree needed.}
  \practice{Little Girl and Problem on Trees}{CF 276E}
    {https://codeforces.com/problemset/problem/276/E}
    {CF 2100 L5}
    {``Add \(x\) to every node at distance exactly \(w\) from \(v\) within
     \(v\)'s subtree''.  Nodes at a fixed depth occupy a contiguous block in
     each subtree's Euler span, so keep \emph{one BIT per depth} over the tour
     positions and range-add \(x\) on depth \(\depth(v)+w\) inside
     \([\tin(v),\tout(v)]\); a node read is one point query in its depth's BIT.
     The lesson: when the query keys on (subtree, depth), index the Euler
     array by depth groups -- the same \(\tin/\tout\) reduction, refined.}
\end{practiceset}

\paragraph{Back-cites and forwards.}
The \(\tin/\tout\) stamps and ``subtree \(=\) interval'' come wholly from
Ch.~5 (\chref{c24}, \(\thmref{c24-range}\)); we supply only the range engine it
deferred.  CSES~1137 ``Subtree Queries'' is Ch.~5's own solved case with a
static prefix array -- replay it here with \(\coderef{c14-subtree}\) shape A and
updates become free.  CF~343D ``Water Tree'' is the same subtree-add /
root-path-read pair with the URL owned by Ch.~6 (\chref{c06}) -- cited by name,
not re-linked.  Forward: ``add on one axis, read on another, offline'' returns
in \secref{c14-offline}; the k-th-by-lifting predicate the Summit CF~1181D
binary-searches over is this section's structure seen from the value side
(\(\secref{c14-order}\)).

% ---------------------------------------------------------------
\begin{recallbox}[Section 18.4 in 10 seconds]
Ch.~5 flattened the tree; a BIT now owns its ranges.  Node\(\,\to\,\)slot
\(\tin(v)\), subtree\(\,\to\,\)interval \([\tin(v),\tout(v)]\)
(\(\defref{c14-eulerbit}\)).  Point-update + subtree-sum is the plain BIT,
subtree-add + point-read is the difference BIT (\(\propref{c14-eulerrange}\),
\(\coderef{c14-subtree}\)); the same two writes read the other way give
root-path sums (\(\propref{c14-path}\), \(\coderef{c14-path}\)).  A depth-parity
sign folds into the stored delta so a propagating update stays one range add
(\(\thmref{c14-parity}\)) -- that solves CF~383C
(\(\probref{c14-prop}\), \(\coderef{c14-383-soln}\)).  Keep one tour
convention per program.
\selftest{
  \item State, in one line each, the reduction for ``subtree sum'' and for
        ``root-path sum'' on a static tree.
  \item Why can \emph{one} difference BIT answer both, and what differs: the
        writes or the read?
  \item For a query ``add to subtree with a sign flipping by level'', what is
        stored and what is the node read?
  \item Give the exact symptom of mixing a \(0\)-based half-open tour with a
        \(1\)-based inclusive BIT.
  \item Which chapter owns \(\tin/\tout\), and which one owns the range engine?
}
\end{recallbox}

% =====================================================================
%  18.5  Offline Sweeps: sort one axis, BIT the other   sec:c14-offline
%  Exports: prop:c14-dom   lst:c14-dom
%           prop:c14-nested lst:c14-nested
%           prop:c14-pbs   lst:c14-pbs   fig:c14-schedule
%  WORKED section (no \solvedproblem). Discharges Ch.17's forward pointer
%  (\chref{c14} in ch17 99_chapter_end): PBS scheduling by reference
%  (\secref{c28-offline}), predicate engine swapped to a BIT.
%  Anti-repetition: Ch.17 (\chref{c28}) owns PBS-with-rollback-DSU -- we reuse
%    the LOCKSTEP SCHEDULE only, and say why a BIT needs no rollback.
%    Ch.1 (\chref{c01}) owns merge count; Ch.5 owns prefix/nested structure.
%  Worked examples (hand-built so no sample is fabricated):
%    * 2D dominance on points (1,1),(2,3),(3,2),(4,4): strict-dominator counts
%      (0,1,1,3).
%    * CSES 2169 Nested Ranges Count on ranges A=[1,6],B=[2,4],C=[8,10],D=[5,7]:
%      (contains, contained-by) = A(1,0) B(0,1) C(0,0) D(0,0), brute-forced.
%    * PBS k-th in static range on a=(4,1,3,2): Q[1,4] k=2 -> 2 ; Q[1,2] k=1 -> 1.
%  Practice URLs (live-verified, urlsweep-free): CF 1284D (2100), CSES 2168.
% =====================================================================
\section{Offline Sweeps: Sort One Axis, BIT the Other}\label{sec:c14-offline}

\begin{previouslybox}[Ch.~17 ran this schedule with a DSU; we swap the engine]
Chapter~17 (\chref{c28}) introduced \term{parallel binary search} (PBS): many
``largest value with property \(P\)'' questions that all binary-search in
\emph{lockstep}, so one shared pass per bit level answers every query at once.
There the predicate \(P\) (``are these two cells connected yet?'') was charged
by a \emph{rollback} DSU, because union operations must be undone between
rounds.  This section keeps the scheduling idea verbatim (\(\secref{c28-offline}\))
and only changes the predicate engine to a BIT -- and notes that a BIT is even
nicer, because \emph{adding to a BIT is monotone along the sweep}, so no
rollback machinery is needed at all.  The wider lesson, back to Ch.~1
(\chref{c01}): an inversion count was a divide-and-conquer there; offline it is
just ``sort one axis, count the other with a running structure.''
\end{previouslybox}

\begin{topicintro}
\textbf{Offline} means all queries are handed to you before any answer is
required, so you may re-order them.  The recurring move is to identify a
\emph{free} ordering axis (one you can sort by without changing answers), sweep
it in monotone order, and let a frequency/sum BIT over the \emph{other} axis
answer each query with the running structure.  Three shapes recur:
\emph{2D dominance counting} (sort by \(x\), count \(y\)-prefixes), \emph{range
nesting} (sort by left end, count right ends -- CSES~2169), and the
\emph{threshold search} where the answer is a value, so you run PBS and each
predicate is one dominance sweep.  In all three the BIT stores counts over a
compressed coordinate, and the schedule -- not the data structure -- is the new
idea.

\textbf{Prerequisites.} \secref{c14-core} (the two walks), \secref{c14-range}
(coordinate compression), \secref{c14-order} (frequency BIT counting by rank),
and Ch.~17's PBS framing (\(\secref{c28-offline}\)) for the third example.

\textbf{After this section you can}
\begin{itemize}[nosep]
  \item count 2D dominated points with a sort and a \(y\)-BIT; \hfill\lv{3}--\lv{4}
  \item solve range-containment \emph{counting} (CSES~2169) by sorting one end
        and BIT-ing the other; \hfill\lv{4}
  \item recognise a query whose answer is a \emph{threshold}, wrap it in PBS,
        and compute the predicate with a fresh sweep each level; \hfill\lv{5}
  \item state why a BIT predicate needs no rollback (unlike Ch.~17's DSU). \hfill\lv{4}
\end{itemize}
\end{topicintro}

\begin{recognitionbox}[Pattern card --``queries can be reordered; each is a prefix count over a second axis'']
\reckey{``for each point, how many points have both coordinates smaller'';
``for each interval, how many it contains / is contained in''; ``k-th smallest
value inside a fixed subarray, many queries''; ``how many pairs \((i,j)\) with
\(i<j\) and a condition on values''; ``the least \(t\) such that a sweep-count
reaches \(k\)''; anything whose brute force is \emph{sort by one key, then a
running tally over another}.}
\recuse{
  \item one comparison axis is free -- sort on it, BIT on the other
        (\(\propref{c14-dom}\), \(\propref{c14-nested}\));
  \item the answer is a value with a monotone predicate -- PBS (\(\secref{c28-offline}\))
        wrapping a dominance sweep (\(\propref{c14-pbs}\));
  \item containment of intervals -- sort left-ends, count right-ends
        (CSES~2169, \(\coderef{c14-nested}\)).
}
\recnotuse{
  \item you must answer \emph{online} (interleaved output, no re-ordering) --
        use the plain/frequency BIT of \secref{c14-core}/\secref{c14-order} or
        a segment tree, Ch.~19 (\chref{c13});
  \item the second axis needs a non-invertible combine (range max with
        arbitrary writes) -- segment tree, not a count BIT;
  \item both axes genuinely need the full 2D structure with rectangle
        updates -- a 2D BIT or range tree, whose \(\bigO{\log^2}\) cost and
        memory are the warning below;
  \item the array is a single static block of queries with no updates -- one
        prefix array already answers in \(O(1)\).
}
One-line test: \emph{can I sort the queries without lying, so that each answer
is ``how much of the other axis has the sweep already seen''?}
\end{recognitionbox}

% ---------------------------------------------------------------
\subsection{2D Dominance: Sort by \(x\), Count \(y\)-Prefixes}
% ---------------------------------------------------------------

\begin{proposition}{Dominated-point counts in \(\bigO{n\log n}\)}{c14-dom}
Given \(n\) points \((x_i,y_i)\) (comparable \(y\) after compression), the
number of points that \(p\) \emph{strictly dominates} -- those \(q\) with
\(q_x\le p_x\) and \(q_y\le p_y\), \(q\ne p\) -- is obtained by sorting points
by \(x\) ascending (ties broken by \(y\) ascending), sweeping that order, and,
just before inserting \(p\)'s \(y\)-rank, reading \(\code{count\_le}(\mathit{rank}_y)\)
on a frequency BIT over \(y\).  The sweep guarantees the BIT holds exactly the
points with smaller \(x\) (or equal \(x\) but not-yet-processed \(y\)), so the
prefix count is the dominance tally.  Total \(\bigO{n\log n}\).
\end{proposition}

\begin{lstlisting}[caption={2D dominance count: sort by x, frequency BIT over compressed y}, label={lst:c14-dom}]
// points[] = (x, yr) with yr = compressed y-rank; sorted by x asc, then y asc
sort(points.begin(), points.end());
vector<int> bit(m + 1, 0);                 // frequency BIT over y (18.3)
for (auto& p : points){
    int dom = 0;
    for (int r = p.yr; r > 0; r -= r & -r) dom += bit[r];  // #seen with y<=p.y
    // dom = points strictly dominated by p (excluding p itself)
    for (int r = p.yr; r <= m; r += r & -r) bit[r]++;      // p joins the seen set
}
\end{lstlisting}

\begin{tracebox}[Dominance on \((1,1),(2,3),(3,2),(4,4)\)]
\(y\)-ranks are already \(1,2,3,4\); sorted by \(x\) (as listed).  Read the
BIT \emph{before} inserting each point:
\[
\renewcommand{\arraystretch}{1.1}
\begin{array}{c|c|c|c}
(x,y) & \text{seen before it} & \code{count\_le}(y) & \text{dominated}\\ \hline
(1,1) & \varnothing & 0 & 0\\
(2,3) & \{(1,1)\} & \#y\!\le\!3=1 & 1\\
(3,2) & \{(1,1),(2,3)\} & \#y\!\le\!2=1 & 1\\
(4,4) & \{(1,1),(2,3),(3,2)\} & \#y\!\le\!4=3 & 3
\end{array}
\]
Dominator counts \((0,1,1,3)\).  Nothing here is a new data structure -- it is
\(\coderef{c14-freq}\) plus a \emph{sort}; the sort is the whole of the new
idea.
\end{tracebox}

\begin{figure}[H]
\centering
\begin{tikzpicture}[diagram, scale=0.8, transform shape]
  \draw[->,gray!70] (0,0) -- (5.2,0) node[below right,font=\tiny]{\(x\)};
  \draw[->,gray!70] (0,0) -- (0,4.2) node[above left,font=\tiny]{\(y\)};
  \foreach \px/\py/\lab in {1/1/{(1,1)},2/3/{(2,3)},3/2/{(3,2)},4/4/{(4,4)}} {
    \filldraw[primaryblue] (\px,\py) circle (2pt) node[above right,font=\tiny] {\lab};}
  \draw[dashed,warningred] (0,0) rectangle (4,4);
  \node[font=\tiny,text=warningred] at (2,-0.5) {sweep \(x\uparrow\): the BIT holds all points left of the dashed edge};
\end{tikzpicture}
\caption{The dominance sweep visualised.  Sorting by \(x\) turns ``count points
below-and-left of \(p\)'' into ``count what the \(y\)-BIT has already swallowed
when the vertical edge reaches \(p\)'s \(x\)''.  Point \((4,4)\) sits with all
three earlier points inside its dominated quadrant, so it reads \(3\).}
\label{fig:c14-schedule}
\end{figure}

% ---------------------------------------------------------------
\subsection{Nested Ranges: sort one end, count the other (CSES 2169)}
% ---------------------------------------------------------------

Range \([l_i,r_i]\) \emph{contains} range \([l_j,r_j]\) iff \(l_i\le l_j\) and
\(r_i\ge r_j\).  That is dominance with one axis flipped -- and after sorting by
one end, the other end is a BIT prefix.

\begin{proposition}{Contains / contained-by counts for every range}{c14-nested}
Sort the ranges by \(l\) ascending, breaking ties by \(r\) \emph{descending}.
\begin{itemize}[nosep,leftmargin=1.4em]
  \item \emph{contains\(_i\)} (how many ranges \(i\) contains): sweep the order
        right-to-left, inserting each range's \(r\) into a frequency BIT; at
        \(i\) read \(\code{count\_le}(r_i)\).  Everything already inserted has
        \(l\ge l_i\), and among those the prefix \(r\le r_i\) are exactly the
        ranges inside \(i\)'s span.
  \item \emph{contained-by\(_i\)} (how many ranges contain \(i\)): sweep the
        same order left-to-right; at \(i\) the number already inserted with
        \(r\ge r_i\) equals \((\text{seen})-\code{count\_le}(r_i-1)\).
\end{itemize}
The tie-break (\(l\) equal \(\Rightarrow\) larger \(r\) first) is what makes
equal-left ranges nest correctly and stops a range counting itself.  Total
\(\bigO{n\log n}\) after compressing the \(r\)-values.
\end{proposition}

\begin{lstlisting}[caption={Nested Ranges Count (CSES 2169): sort left ends, BIT over right ends}, label={lst:c14-nested}]
// rng[] = (l, r, id); r compressed to rank in 1..m
sort(rng, [](A,B){ return A.l!=B.l ? A.l<B.l : A.r>B.r; });   // l asc, r DESC
vector<int> bit(m+1,0);
// contains[i]: sweep from the right
for (int i=n-1;i>=0;i--){
    int c=0; for(int x=rng[i].r;x>0;x-=x&-x) c+=bit[x];       // # with r<=rng[i].r
    contains[rng[i].id]=c;                                    //   (all have l>=l_i)
    for(int x=rng[i].r;x<=m;x+=x&-x) bit[x]++;
}
fill(bit.begin(),bit.end(),0);
// contained_by[i]: sweep from the left, count r >= rng[i].r among seen
for (int i=0;i<n;i++){
    int le=0; for(int x=rng[i].r-1;x>0;x-=x&-x) le+=bit[x];   // # with r< rng[i].r
    contained_by[rng[i].id]= i - le;                          // seen = i ; >= r_i
    for(int x=rng[i].r;x<=m;x+=x&-x) bit[x]++;
}
\end{lstlisting}

\begin{tracebox}[Nested Ranges on four hand-built ranges A,B,C,D]
Sort by \(l\) asc, \(r\) desc: \(A(1,6), B(2,4), D(5,7), C(8,10)\).  Compress
\(r\)-values \(4,6,7,10\to\) ranks \(1,2,3,4\).
\emph{contains} (right-to-left sweep; BIT over \(r\)-rank):
\(C\): query\((4){=}0\); insert~4.
\(D\): query\((3){=}0\); insert~3.
\(B\): query\((1){=}0\); insert~1.
\(A\): query\((2){=}\)count of ranks\(\le2\) seen\(=\{1\}\)=1. So
contains \(=(A\,1,\,B\,0,\,C\,0,\,D\,0)\).
\emph{contained-by} (left-to-right; \(i-\)query\((r_i-1)\)):
\(A\): \(0-\text{query}(1)=0\).  \(B\): \(1-\text{query}(0)=1\).
\(D\): \(2-\text{query}(2)=2-0=0\).  \(C\): \(3-\text{query}(3)=3-2=0\).  So
contained-by \(=(A\,0,\,B\,1,\,C\,0,\,D\,0)\).
Pairing per original range: \(A(1,0),B(0,1),C(0,0),D(0,0)\) -- verified against
a brute scan of the four pairs.  (This is a hand-built set; CSES~2169 has its
own statement and samples.)
\end{tracebox}

\begin{insightbox}[Why the tie-break is load-bearing]
Sorting equal-left ranges by \emph{descending} right end means the widest comes
first, so a narrower same-left range is counted as contained by it, and neither
sees itself (it is inserted only after its own query).  Use \emph{ascending}
right on the tie and equal-left ranges mis-nest -- the classic off-by-one WA of
this problem, and the reason \(\propref{c14-nested}\) states the order twice.
\end{insightbox}

% ---------------------------------------------------------------
\subsection{Parallel Binary Search over a BIT: k-th in a Static Range}
% ---------------------------------------------------------------

The last two answers were \emph{counts}.  Some questions ask for a
\emph{threshold}: the least value \(v\) making a dominance-count reach \(k\) --
for instance the k-th smallest number inside a fixed subarray.  The predicate
``is the k-th smallest of \(a[\ell..r]\) at most \(v\)?'' is ``\(\#\{i\in[\ell,r]:
a_i\le v\}\ge k\)'', which is exactly a dominance count (\(\propref{c14-dom}\))
if we sort by value.  Because the predicate is monotone in \(v\), we binary
search it -- and, with many queries, we run Ch.~17's \term{parallel} version so
one shared sweep serves every query at the same bit level.

\begin{proposition}{k-th smallest in a range by PBS with a position BIT}{c14-pbs}
Each query holds an interval \([\ell,r]\) and a target \(k\); the answer is a
value.  Maintain per-query \((lo,hi)\) over the value range.  Repeat until all
\((lo{=}hi)\): \emph{bucket} every unfinished query at its current \(mid\);
sweep the value domain ascending, inserting each array element's \emph{position}
into a BIT as its value becomes \(\le mid\); for every query bucketed at this
\(mid\), the range count \(\code{query}(r)-\code{query}(\ell-1)\) is the
predicate ``\(\ge k\)?'' -- if true set \(hi=mid\), else \(lo=mid+1\).  Because
\(mid\)-values are processed in ascending order within a round, the BIT only
gains elements (a \(y\)-prefix sweep), so there is \textbf{no rollback}: at the
end of a round we discard the whole BIT and rebuild next round.  This is the
only place Ch.~17's rollback DSU is replaced -- same schedule, easier engine.
\end{proposition}

\begin{lstlisting}[caption={PBS skeleton: lockstep binary search whose predicate is one BIT sweep}, label={lst:c14-pbs}]
// a[] static with compressed values 1..m; queries (l,r,k) -> kth smallest value
int lo[MQ], hi[MQ];                       // per query, over value ranks
for (int it = 0; it < LOG; it++){         // LOG = ceil(log2 m) rounds, all in lockstep
    vector<int> bucket[m + 1];            // queries grouped by their current mid
    for (int q : open) bucket[(lo[q]+hi[q])/2].push_back(q);
    fill(bit.begin(), bit.end(), 0);      // FRESH BIT each round (no rollback needed)
    // sweep value v ascending; positions with a[] <= v sit in the BIT:
    vector<int> pos = positions_sorted_by_value();   // precomputed once
    for (int v = 1, ptr = 0; v <= m; v++){
        while (ptr < n && val_at[pos[ptr]] <= v) add(bit, pos[ptr++], 1);
        for (int q : bucket[v]) {                       // predicate at mid = v
            long long cnt = query(bit, R[q]) - query(bit, L[q] - 1);
            if (cnt >= K[q]) hi[q] = v; else lo[q] = v + 1;
        }
    }
    open = queries with lo < hi;
}
// answer value = lo[q] for each q  (lo == hi at exit)
\end{lstlisting}

\begin{tracebox}[PBS on \(a=(4,1,3,2)\): k-th in a range]
Compressed values \(1<2<3<4\) as ranks \(1,2,3,4\).  Two queries:
\(Q_1=[1,4],k{=}2\); \(Q_2=[1,2],k{=}1\).  Both start \((lo,hi)=(1,4)\).\\
\emph{Round~1, mid\(=2\):} fresh BIT, insert positions with value\(\le2\):
value1@pos2, value2@pos4 \(\Rightarrow\) BIT holds \(\{2,4\}\).
\(Q_1\): count in \([1,4]=2\ge2\) \(\Rightarrow\) \(hi=2\).
\(Q_2\): count in \([1,2]=1\ge1\) \(\Rightarrow\) \(hi=2\).\\
\emph{Round~2, mid\(=1\):} insert value\(\le1\): \(\{2\}\).
\(Q_1\): count in \([1,4]=1\ge2\)? no \(\Rightarrow\) \(lo=2\).  Now \(lo{=}hi{=}2\)
\(\Rightarrow\) answer \(2\).
\(Q_2\): count in \([1,2]=1\ge1\)? yes \(\Rightarrow\) \(hi=1\), \(lo{=}hi{=}1\)
\(\Rightarrow\) answer \(1\).\\
So k-th smallest of \(\{4,1,3,2\}\) with \(k{=}2\) is \(\mathbf{2}\) (sorted
\(1,2,3,4\)), and the 1st of \(\{4,1\}\) is \(\mathbf{1}\) \checkmark.  Two
rounds, one shared sweep each -- the lockstep from Ch.~17, predicate now a BIT.
\end{tracebox}

\begin{comparisonbox}[PBS-with-BIT (here) vs PBS-with-rollback-DSU (Ch.~17)]
\begin{itemize}[nosep,leftmargin=1.4em]
  \item \textbf{Same} \emph{schedule}: all queries advance one bit per round;
        a round is one monotone pass; complexity \(\bigO{(n+q)\log(\text{domain})
        \log n}\).
  \item \textbf{Different} \emph{engine}: Ch.~17 connected cells with DSU, whose
        unions are not invertible, so it \emph{rolls back} to the round's start.
        Here the predicate is a \(y\)-prefix count: the BIT along an ascending
        value sweep only \emph{appends}, so we rebuild it per round and never
        roll back -- the \code{fill(bit,...)} reset in \(\coderef{c14-pbs}\).
  \item \textbf{Different} \emph{shape of the predicate}: DSU answers
        ``connected by threshold''; the BIT answers ``at least \(k\) elements
        below threshold in a span''.  Recognising which one your problem needs
        is the whole design step.
\end{itemize}
This is exactly the obligation Ch.~17 pinned on this chapter: the sweep-a-batch
idea, run with a Fenwick predicate instead of a rollback DSU.
\end{comparisonbox}

% ---------------------------------------------------------------
\subsection{Common Pitfalls}
% ---------------------------------------------------------------

\begin{warningbox}[Trap 1: sorting the wrong way on ties]
In nested ranges (\(\propref{c14-nested}\)) equal left ends must order by
\emph{descending} right end, or two ranges sharing a left end fail to nest and
a range counts itself.  In dominance (\(\propref{c14-dom}\)) equal \(x\) must
order by \(y\) so that a query never sees an equal-\(x\) point it should.  A
tie-break is part of the reduction, not a cosmetic detail.
\end{warningbox}

\begin{warningbox}[Trap 2: rolling a BIT back when you can just rebuild]
Carrying over Ch.~17's rollback habit here costs complexity and bugs.  The BIT
predicate is monotone along the sweep, so a round only appends; discard and
rebuild each round (an \(\bigO{m}\) clear).  If you find yourself writing an
undo stack for a Fenwick, you have mis-identified the predicate's shape.
\end{warningbox}

\begin{warningbox}[Trap 3: forgetting to compress the second axis]
The BIT indexes the \emph{other} axis (\(y\), or \(r\), or position).  If those
coordinates reach \(10^9\), the array is impossible until you rank-compress
them (\(\propref{c14-compress}\)).  Compress \emph{all} coordinates that appear
in any query bound too, or a prefix query reads a rank that was never inserted.
\end{warningbox}

\begin{warningbox}[Trap 4: 32-bit counts and the \(\bigO{\log^2}\) temptation]
Dominance and nesting counts reach \(\binom{n}{2}\) (accumulate in \code{long
long}).  And when a problem really needs \emph{rectangle} updates rather than a
sweep, a naive 2D BIT is \(\bigO{\log^2 n}\) and \(O(n\log n)\) memory -- often
blowing the limit where a single-axis sweep would not.  Check whether a sort
collapses your 2D problem to a 1D BIT before reaching for the heavier tool.
\end{warningbox}

% ---------------------------------------------------------------
\subsection{Practice Problems}
% ---------------------------------------------------------------

\begin{stuckbox}
\begin{itemize}[nosep,leftmargin=1.4em]
  \item Which axis is genuinely \emph{free} to sort on without changing answers?
  \item What does the running BIT store -- a count over compressed \(y\), over
        right ends, or over positions?
  \item Is your answer a \emph{count} (one sweep) or a \emph{threshold} (PBS
        wrapping sweeps)?
  \item Did you compress every coordinate, tie-break correctly, and use 64-bit
        accumulators?
\end{itemize}
\end{stuckbox}

\begin{practiceset}
  \practice{New Year and Conference}{CF 1284D}
    {https://codeforces.com/problemset/problem/1284/D}
    {CF 2100 L5}
    {Lecture \(i\) has interval \([sa_i,ea_i]\) at venue \(a\) and \([sb_i,eb_i]\)
     at venue \(b\); Hyunuk is happy iff \emph{no} pair overlaps in exactly one
     venue.  Reduce to: for every pair overlapping in \(a\), check they also
     overlap in \(b\), then swap venues.  Sort lectures by \(sa\) and sweep,
     keeping the currently \(a\)-active lectures in an ordered structure keyed
     by their \(b\)-endpoints; if any active lecture's \(b\)-interval is
     disjoint from the current one, that pair is the witness \(\Rightarrow\) NO.
     The \emph{same} sort-one-axis/read-the-other sweep as \(\coderef{c14-dom}\),
     but the live structure is a multiset (the combine is an interval test, not
     a count -- the ``non-invertible combine'' case from the pattern card, made
     concrete).}
  \practice{Nested Ranges Check}{CSES 2168}
    {https://cses.fi/problemset/task/2168}
    {CSES L4}
    {The boolean twin of the worked CSES~2169: for each range report only
     ``does it contain some other / is it contained in some other''.  Run the
     very sort of \(\coderef{c14-nested}\), but you no longer need a BIT -- a
     running \(\min/\max\) of right ends suffices, because existence is a
     prefix extremum, not a prefix sum.  The lesson: recognise when the sweep
     wants a count (BIT) versus an extreme (\(\bigO{1}\) running value).}
\end{practiceset}

\paragraph{Back-cites and forwards.}
The lockstep \emph{scheduling} of \(\propref{c14-pbs}\) is Ch.~17's idea, cited
by reference (\(\secref{c28-offline}\)); only the predicate engine is swapped to
\secref{c14-core}'s BIT, discharging the forward pointer Ch.~17 left on this
chapter.  The counting sweeps reuse \secref{c14-order}'s frequency BIT and
\secref{c14-range}'s compression; Ch.~1 (\chref{c01}) owns the merge-count
view of the same inversions.  Forward: the Summit CF~1181D is precisely a PBS
whose predicate is a \(y\)-prefix count on a BIT (\(\thmref{c14-lift}\)'s
lifting answers each bucketed query's ``which slot crosses \(k\)''), composing
this section's schedule with \secref{c14-order}'s search.

% ---------------------------------------------------------------
\begin{recallbox}[Section 18.5 in 10 seconds]
Offline \(=\) you may sort.  Find the free axis, sweep it, and let a BIT over
the \emph{other} (compressed) axis answer a prefix count: dominance
\(\propref{c14-dom}\) (sort \(x\), count \(y\)), nested ranges
\(\propref{c14-nested}\) / CSES~2169 (sort left ends, count right ends).  When
the answer is a \emph{threshold}, wrap that count in Ch.~17's parallel binary
search (\(\propref{c14-pbs}\), \(\coderef{c14-pbs}\)) -- one shared sweep per
bit level, and because a BIT only appends along the sweep, \textbf{no rollback}.
Sort ties deliberately, compress both axes, accumulate in \code{long long}.
\selftest{
  \item For points \((1,1),(2,3),(3,2),(4,4)\), say which axis you sort and
        what the BIT indexes, then give \((4,4)\)'s dominance count.
  \item Why does Nested Ranges sort ties by \emph{descending} right end?
  \item What makes the PBS-with-BIT predicate monotone, and why does that let
        us skip Ch.~17's rollback?
  \item Give the one change from Ch.~17's PBS and the one thing kept identical.
  \item When does a range sweep want a running \(\min/\max\) instead of a BIT?
}
\end{recallbox}

% =====================================================================
%  Chapter 18 review -- the 14-station revision contract.
%  Revision-only: tables, one-liners, pointers back via \secref/\thmref/
%  \propref/\defref/\coderef/\figref/\probref.  No re-derivations.
%  Summit: CF 1181D Irrigation -- water-fill the min-count host city, find
%  the level with a monotone count-space search and name the (offset)-th
%  active city with the frequency-BIT k-th lift (lst:c14-kth).  Samples
%  verified live against the CF statement:
%    (1) n=6 m=4 q=10  a=3 1 1 1 2 2  k=7..16 -> 4 3 4 2 3 4 1 2 3 4
%    (2) n=4 m=5 q=4   a=4 4 5 1      k=15,9,13,6 -> 5 3 3 3
%  Exports prob:c14-summit, lst:c14-kit, lst:c14-summit.
% =====================================================================

\chaptersection{Chapter 18 Summary}

One sentence: \emph{a Fenwick tree is an array that secretly stores a binary
tree over the index axis, so a single \(O(\log n)\) walk answers a prefix sum
and a second walk folds in a point update -- and every question Chapter~16 and
Chapter~17 kept handing back as "just take a prefix" (a range sum, a \(k\)-th
live item, a subtree profile swept offline) is one \(\mathit{lowbit}\) hop
sequence away.} Chapter~5 flattened a tree into an array; this chapter gives
that array a \emph{live prefix engine}. Each row is a pointer, not a proof.

\begin{center}
\renewcommand{\arraystretch}{1.25}
\begin{tabular}{@{}p{0.22\linewidth} p{0.50\linewidth} p{0.20\linewidth}@{}}
\toprule
\textbf{Result} & \textbf{Statement} & \textbf{Where} \\
\midrule
Responsibility interval & Slot \(i\) owns \([\,i-\mathit{lowbit}(i)+1,\ i\,]\);
\(\mathit{lowbit}(i)=i\mathbin\&(-i)\). & \defref{c14-interval} \\
Prefix walk & \(\sum_{1}^{i}A\) is read by jumping \(i\!\to\!i-\mathit{lowbit}(i)\);
\(\bigO{\log n}\). & \thmref{c14-query}, \coderef{c14-query} \\
Point-update walk & Add \(\delta\) to \(A[i]\) by pushing it into every interval
covering \(i\): \(i\!\to\!i+\mathit{lowbit}(i)\). & \coderef{c14-update} \\
Linear build & Hand each partial block to the slot its block extends: \(\bigO{n}\)
build, no updates. & \propref{c14-linear}, \coderef{c14-build} \\
Difference BIT & Two point updates \(+d@l,\ -d@(r{+}1)\) make a range add; a
prefix read is the point value. & \defref{c14-diff}, \coderef{c14-diff-build} \\
Two-BIT range sum & \(\mathit{prefix}(x)=(x{+}1)Q_1(x)-Q_2(x)\) turns range-add
into range-sum. & \thmref{c14-twobit}, \coderef{c14-twobit} \\
Compression & Any value/coordinate set becomes ranks \(1..n\); the BIT only needs
a dense index axis. & \propref{c14-compress}, \coderef{c14-compress} \\
Prefix-min caveat & A min-BIT is \emph{append-only}: valid while values never
decrease back out. & \propref{c14-prefixmin}, \coderef{c14-min} \\
Order statistic & Binary-lift down the tree for the smallest index whose prefix
count reaches \(k\): the \(k\)-th live item. & \thmref{c14-lift}, \coderef{c14-kth} \\
Inversion sweep & A right-to-left frequency-BIT pass counts "smaller after"
online; Ch.~1 owns the merge-sort twin. & \propref{c14-inversion}, \coderef{c14-inv} \\
Euler \(+\) BIT & Ch.~5's \(\tin/\tout\) flat array driven by a BIT: subtree
add/sum by interval, path reads by enter\(+\)1/leave\(-\)1. & \propref{c14-eulerrange}, \propref{c14-path} \\
Dominance & Sort one axis, count the other with a BIT: offline 2D
\(\#\{p_x\le q_x,\ p_y\le q_y\}\). & \propref{c14-dom}, \coderef{c14-dom} \\
PBS over BIT & Parallel binary search whose predicate is a BIT sweep; a BIT only
appends along the sweep, so \emph{no rollback}. & \propref{c14-pbs}, \coderef{c14-pbs} \\
\bottomrule
\end{tabular}
\end{center}

\medskip
{\small\textbf{Algorithms from this chapter.}}
\begin{center}
\renewcommand{\arraystretch}{1.2}
\begin{tabular}{@{}p{0.27\linewidth} l l >{\raggedright\arraybackslash}p{0.26\linewidth}@{}}
\toprule
\textbf{Algorithm} & \textbf{Prep} & \textbf{Per op} & \textbf{Code} \\
\midrule
Prefix sum / point add & \(\bigO{n}\) build & \(\bigO{\log n}\) & \coderef{c14-query}, \coderef{c14-update} \\
CSES 1648 (Dynamic Range Sum) & \(\bigO{n}\) & \(\log n\) & \coderef{c14-1648-soln} \\
Difference BIT (range add / point read) & \(\bigO{n}\) diff & \(\log n\) & \coderef{c14-diff-build} \\
CSES 1651 (Range Update) & \(\bigO{n}\) & \(\log n\) & \coderef{c14-1651-soln} \\
Two-BIT (range add / range sum) & \(\bigO{n}\) & \(\log n\) & \coderef{c14-twobit} \\
Coordinate compression & \(n\log n\) sort & \(\log n\) & \coderef{c14-compress} \\
Frequency \(k\)-th by lifting & \(\bigO{n}\) & \(\log n\) select & \coderef{c14-kth} \\
Inversion / smaller-after & one sweep & \(\log n\)/elem & \coderef{c14-inv} \\
CF 61E (Enemy is weak) & sort ranks & \(\log n\) & \coderef{c14-61-soln} \\
Euler subtree add/sum & \(\bigO{n}\) DFS & \(\log n\) & \coderef{c14-subtree} \\
CF 383C (Propagating tree) & \(\bigO{n}\) Euler & \(\log n\) & \coderef{c14-383-soln} \\
2D dominance count & \(n\log n\) sort & \(\log n\) & \coderef{c14-dom} \\
Nested ranges (CSES 2169) & sort left ends & \(\log n\) & \coderef{c14-nested} \\
PBS over a BIT & \(\bigO{\log R}\) sweeps & \(\log n\) & \coderef{c14-pbs} \\
CF 1181D (Summit, Irrigation) & count + sort & \(\log n\) & \coderef{c14-summit} \\
\bottomrule
\end{tabular}
\end{center}

\chaptersection{Master Decision Flowchart}

Post-study ladder: walk top to bottom, stop at the first branch that matches.
It decides \emph{which BIT flavour} the problem wants and hands off what this
chapter does \emph{not} own: arbitrary-range min with deletions (Ch.~19 segtree),
the merge-sort inversion twin (Ch.~1), lazy propagation (Ch.~20).

\begin{lstlisting}[language={}, basicstyle=\ttfamily\small, numbers=none, frame=single]
  You need a live PREFIX over an index axis.  Which flavour?
  |
  |-- point update + range sum, values combine by + (invertible)?
  |     +-- plain BIT (18.1, lst:c14-query + lst:c14-update). Build O(n)
  |        with lst:c14-build. Range sum = prefix(r)-prefix(l-1).
  |
  |-- add d to A[l..r], but you only ever READ a single point?
  |     +-- difference BIT (18.2, lst:c14-diff-build): +d@l, -d@(r+1),
  |        point read = prefix. Two updates per range (prop:c14-diff).
  |
  |-- add d to A[l..r] AND sum A[l..r]?
  |     +-- two-BIT identity (18.2, lst:c14-twobit) OR lazy segtree.
  |        If you also need min / non-invertible combine -> Ch.19/20 (c13/c15).
  |
  |-- values are huge / sparse, only their ORDER matters?
  |     +-- coordinate compression first (18.2, lst:c14-compress), then any BIT.
  |
  |-- "k-th live item / select / how many below x"?
  |     +-- frequency BIT + binary lifting (18.3, lst:c14-kth). Prefix count
  |        is monotone in the index, so lift down the tree.
  |
  |-- count inversions / smaller-after / greater-before, ONLINE?
  |     +-- one right-to-left frequency sweep (18.3, lst:c14-inv). Merge-sort
  |        (Ch.1, c01) is the offline divide-and-conquer twin -- same bound.
  |
  |-- min over a growing prefix, never pulled back out?
  |     +-- prefix-min BIT (18.2, lst:c14-min), APPEND-ONLY (prop:c14-prefixmin).
  |        Arbitrary decrement -> use a segtree; the caveat is a correctness bug.
  |
  |-- subtree add / subtree sum, or root-path reads, on a TREE?
  |     +-- Euler-flatten (Ch.5, c24) then drive tin/tout with a BIT
  |        (18.4, lst:c14-subtree; path via enter+1/leave-1, lst:c14-path).
  |
  |-- 2D pairs / nested ranges / "how many points dominate q"? offline OK?
  |     +-- sort ONE axis, BIT the other (18.5, lst:c14-dom, lst:c14-nested).
  |
  |-- answer is a THRESHOLD (least x passing a monotone BIT test), many queries?
  |     +-- parallel binary search over one shared BIT (18.5, lst:c14-pbs);
  |        the BIT only appends along the sweep, so NO rollback (Ch.17 did).
  |
  |-- "k-th value in a STATIC range [l..r]"?
        +-- offline value-sweep / PBS over a frequency BIT (18.3 + 18.5); the
           persistent segtree is Ch.19's answer -- this is the offline route.
\end{lstlisting}

\chaptersection{Common Mistakes Checklist}

\begin{itemize}[nosep, leftmargin=1.4em]
  \item \textbf{\(\mathit{lowbit}\) on a signed/zero index.} \(\mathit{lowbit}(i)=i\&(-i)\)
        is only meaningful for \(i\ge 1\); a 0-based BIT loop spins forever
        (\(0+\mathit{lowbit}(0)=0\)). Fenwick is \emph{1-indexed} (\defref{c14-interval}).
  \item \textbf{Query walk vs update walk swapped.} Prefix sum does
        \(i-=\mathit{lowbit}(i)\); point add does \(i+=\mathit{lowbit}(i)\). One
        sign flips the tree from a reader into a writer (\(\coderef{c14-query}\)
        vs \(\coderef{c14-update}\)).
  \item \textbf{Building by \(n\) updates (\(\bigO{n\log n}\)) when \(\bigO{n}\)
        is on offer.} The linear build hands each block forward
        (\(\coderef{c14-build}\)); on \(n=10^6\) the \(\log n\) build is the TLE.
  \item \textbf{Difference BIT missing the \(r+1\) update.} Range add is
        \emph{two} point writes; omitting \(-d@(r{+}1)\) leaks the delta past the
        right end (\(\defref{c14-diff}\)).
  \item \textbf{Treating a min-BIT as deletable.} Prefix-min is append-only; a
        value that later decreases cannot be un-committed
        (\(\propref{c14-prefixmin}\)). That direction is a segtree.
  \item \textbf{Lifting search off-by-one.} The \(k\)-th test is
        \(\mathit{acc}+t[\mathit{next}] < k\) (strict), then return
        \(\mathit{pos}+1\); a \(\le\) returns the wrong neighbour
        (\(\thmref{c14-lift}\), \(\coderef{c14-kth}\)).
  \item \textbf{Forgetting compression on a sparse axis.} A BIT over values up to
        \(10^9\) is not the problem; you must first map values to ranks
        (\(\coderef{c14-compress}\)).
  \item \textbf{\(\tin/\tout\) convention drift.} Ch.~5's \(\tout\) is inclusive;
        subtree sum is \(\mathit{query}(\tout)-\mathit{query}(\tin-1)\) and range
        add writes at \(\tout+1\). Mixing inclusive/exclusive mis-answers every
        subtree (\(\propref{c14-eulerrange}\)).
  \item \textbf{64-bit sums.} Prefix sums and inversion counts reach
        \(\bigO{n^2}\)/\(\bigO{n\cdot\text{val}}\); the tree array is
        \code{long long} the moment values or counts grow (\(\coderef{c14-query}\)).
  \item \textbf{Re-rolling a BIT in PBS.} A BIT sweep is monotone in the sorted
        axis, so PBS needs no undo stack -- unlike Ch.~17's rollback DSU. Copying
        the rollback habit here is pure wasted constant (\(\propref{c14-pbs}\)).
\end{itemize}

\chaptersection{WA Clinic}

Verdicts are symptoms; the causes are the short list this chapter drilled.
Diagnose top-down -- every leaf names a listing.

\begin{figure}[H]
\centering
\begin{tikzpicture}[diagram,
    v/.style={boxnode, fill=primaryblue!12, draw=primaryblue, text width=22mm, align=center},
    b/.style={boxnode, fill=difforange!12, draw=difforange, text width=12mm},
    l/.style={boxnode, fill=successgreen!10, draw=successgreen, text width=56mm,
              font=\scriptsize, align=left}]
  \node[v] (verdict) {\textbf{Verdict on the judge}};
  \node[b, below left=8mm and 4mm of verdict] (wa) {\textbf{WA}};
  \node[b, below right=8mm and 4mm of verdict] (tle) {\textbf{TLE}};
  \node[l, below right=1mm and 2mm of wa.south] (wa1)
    {0-based loop spins or mis-reads: \(\mathit{lowbit}(0)=0\), or the query walk
     uses \(i+=\mathit{lowbit}\) instead of \(i-=\mathit{lowbit}\)
     (\coderef{c14-query}, \coderef{c14-update})};
  \node[l, below left=1mm and 2mm of wa.south] (wa2)
    {difference BIT missing the \(-d@(r{+}1)\) write; inclusive \(\tout\) used as
     exclusive in a subtree sum (\coderef{c14-diff-build}, \propref{c14-eulerrange})};
  \node[l, below=1mm of tle.south] (tle1)
    {\(\bigO{n\log n}\) build where \(\bigO{n}\) was required; un-compressed axis
     sized \(10^9\) (\coderef{c14-build}, \coderef{c14-compress})};
  \node[l, below left=5mm and 30mm of tle1.north] (tle2)
    {int overflow on prefix sums / inversions -- tree array not
     \code{long long} (\(\coderef{c14-inv}\))};
  \draw[flow] (verdict) -- (wa);
  \draw[flow] (verdict) -- (tle);
  \draw[flow] (wa) -- (wa1);
  \draw[flow] (wa) -- (wa2);
  \draw[flow] (tle) -- (tle1);
  \draw[flow] (tle) -- (tle2);
\end{tikzpicture}
\caption*{WA: the walk went the wrong way, or a range write forgot its partner.
TLE: you built with updates, sized the tree by value, or overflowed to
\code{int}. RE: an un-compressed \(10^9\) index into a fixed array.}
\end{figure}

\chaptersection{Quick Reference Card}

\begin{quickref}[Chapter 18 at a glance]
\textbf{Plain BIT.} \(\mathit{lowbit}(i)=i\&(-i)\).
\code{query(i)}: \(s\!+=\!t[i];\,i-=\!\mathit{lowbit}(i)\) sums \(A[1..i]\).
\code{update(i,d)}: \(t[i]\!+=\!d;\,i+=\!\mathit{lowbit}(i)\) adds to \(A[i]\).
Range sum \(=\) \code{query(r)}\(-\,\)code{query(l-1)}
(\(\thmref{c14-query}\), \(\coderef{c14-query}\)).

\medskip
\textbf{O(n) build.} \(t[i]+=\!A[i]\); let \(j=i+\mathit{lowbit}(i)\); if
\(j\le n\) then \(t[j]+=\!t[i]\) (\(\propref{c14-linear}\), \(\coderef{c14-build}\)).

\medskip
\textbf{Range add / point read.} Difference BIT: write \(+d\) at \(l\), \(-d\) at
\(r{+}1\); the value at \(x\) is \code{query(x)}
(\(\defref{c14-diff}\), \(\coderef{c14-diff-build}\)).

\medskip
\textbf{Range add / range sum.} Two BITs \(Q_1,Q_2\); \(\mathit{prefix}(x)=(x{+}1)Q_1(x)-Q_2(x)\);
a range add is two \(\mathit{internalAdd}\) calls
(\(\thmref{c14-twobit}\), \(\coderef{c14-twobit}\)).

\medskip
\textbf{Order statistic.} Frequency BIT; lift down the tree for the least index
with prefix count \(\ge k\): \(k\)-th live item, strictly-\(<x\) count
(\(\thmref{c14-lift}\), \(\coderef{c14-kth}\)).

\medskip
\textbf{On a tree.} Flatten once (Ch.~5 \(\tin/\tout\)); a subtree is an interval
(BIT range op), a root path is enter\(+\)1/leave\(-\)1
(\(\propref{c14-eulerrange}\), \(\coderef{c14-subtree}\), \(\coderef{c14-path}\)).

\medskip
\textbf{Offline.} Sort one axis, BIT the other for dominance/nested counts
(\(\coderef{c14-dom}\), \(\coderef{c14-nested}\)); wrap that count in a parallel
binary search for a threshold -- no rollback, a BIT only appends
(\(\propref{c14-pbs}\), \(\coderef{c14-pbs}\)).

\medskip
\textbf{Translation.} "sum a range, edit a point" \(\to\) plain BIT. "add to a
range, read a point" \(\to\) difference BIT. "add to a range, sum a range" \(\to\)
two-BIT (or segtree). "which live value is the \(k\)-th? how many below \(x\)?"
\(\to\) frequency lifting. "subtree/path on a tree" \(\to\) Euler \(+\) BIT.
"how many points dominate \(q\)" / "least \(x\) passing a prefix test" \(\to\)
offline sweep, PBS over a BIT.
\end{quickref}

\begin{lstlisting}[caption={Chapter 18 kit: the plain BIT, the range-add/range-sum pair, and the k-th lift, complete and production-shaped}, label={lst:c14-kit}]
struct BIT {                                  // 1-indexed, long long payload
    int n; vector<long long> t;
    BIT(int n):n(n),t(n+1,0){}
    static int lowbit(int i){return i&(-i);}
    void add(int i,long long d){for(;i<=n;i+=lowbit(i))t[i]+=d;}   // point write
    long long pre(int i){long long s=0;for(;i>0;i-=lowbit(i))s+=t[i];return s;} // prefix
    long long range(int l,int r){return pre(r)-pre(l-1);}          // invertible combine
    // O(n) build: hand each partial block forward
    void build(const vector<long long>&a){t.assign(n+1,0);
        for(int i=1;i<=n;i++){t[i]+=a[i];int j=i+lowbit(i);if(j<=n)t[j]+=t[i];}}
    // k-th: least idx whose prefix sum reaches k (monotone counts). hb = 2^floor(log2 n)
    int kth(long long k,int hb){int pos=0;long long acc=0;
        for(int pw=hb;pw>0;pw>>=1){int nx=pos+pw;
            if(nx<=n && acc+t[nx]<k){pos=nx;acc+=t[nx];}} return pos+1;}
};
// range add + range sum with two BITs (thm:c14-twobit)
struct RangeAddSum{int n;BIT q1,q2;RangeAddSum(int n):n(n),q1(n),q2(n){}
    void ia(int i,long long v){q1.add(i,v);q2.add(i,v*(i-1));}
    void range_add(int l,int r,long long v){ia(l,v);ia(r+1,-v);}
    long long pref(int x){return (long long)(x+1)*q1.pre(x)-q2.pre(x);}
    long long range_sum(int l,int r){return pref(r)-pref(l-1);}};
// highest power of two <= n, for kth:   int hb=1;while(hb*2<=n)hb*=2;
\end{lstlisting}

\chaptersection{Write-From-Memory Drills}

Five signature skeletons, each lifted from a listing already in the chapter.
Cover the marked line, reproduce from memory, then check the answer one line
below.

\begin{lstlisting}[caption={Drill 1 -- prefix query + point update (from lst:c14-query / lst:c14-update)}]
int lowbit(int i) { return i & (-i); }
long long query(int i) {                 // sum of A[1..i]
    long long s = 0;
    for (; i > 0; i -= lowbit(i)) s += tree[i];   // WALK DOWN
    return s;
}
void update(int i, long long d) {        // add d to A[i]
    for (; i <= n; i += lowbit(i)) tree[i] += d;  // WALK UP
}
\end{lstlisting}
\noindent{\scriptsize Key: the query walks \emph{down} (\(-\)lowbit) collecting
disjoint responsibility blocks; the update walks \emph{up} (\(+\)lowbit) touching
every block that contains \(i\). Range sum \(=\) \code{query(r)-query(l-1)}.}

\begin{lstlisting}[caption={Drill 2 -- O(n) linear build (from lst:c14-build)}]
for (int i = 1; i <= n; i++) {
    tree[i] += a[i];                        // own slot's value in
    int j = i + lowbit(i);                  // the block that EXTENDS i's
    if (j <= n) tree[j] += tree[i];         // hand the partial block forward
}
\end{lstlisting}
\noindent{\scriptsize Key: each \(t[i]\) already sums \([i-\mathit{lowbit}(i)+1,i]\),
so adding it to \(t[j]\) extends \(j\)'s interval in one pass -- \(\bigO{n}\), not
\(n\) updates (\(\propref{c14-linear}\)).}

\begin{lstlisting}[caption={Drill 3 -- difference BIT + two-BIT range sum (from lst:c14-diff-build / lst:c14-twobit)}]
// range add / POINT read
void range_add_diff(int l, int r, long long d) { add(l, d); add(r + 1, -d); }
long long point_read(int x) { return query(x); }         // prefix == value

// range add / RANGE sum  (two BITs)
void internal_add(int i, long long v) { bit1.add(i, v); bit2.add(i, v * (i - 1)); }
void range_add(int l, int r, long long v) { internal_add(l, v); internal_add(r + 1, -v); }
long long prefix(int x) { return (x + 1) * bit1.query(x) - bit2.query(x); }
long long range_sum(int l, int r) { return prefix(r) - prefix(l - 1); }
\end{lstlisting}
\noindent{\scriptsize Key: forgetting the \(-d@(r{+}1)\) write in the difference
BIT leaks the delta past the right end; the two-BIT form is just the difference
idea weighted by position (\(\thmref{c14-twobit}\)).}

\begin{lstlisting}[caption={Drill 4 -- k-th by binary lifting (from lst:c14-kth)}]
int kth(long long k) {           // least idx with prefix count >= k; hb = 2^floor(log2 n)
    int pos = 0; long long acc = 0;
    for (int pw = hb; pw > 0; pw >>= 1) {
        int nxt = pos + pw;
        if (nxt <= n && acc + tree[nxt] < k) {   // STRICT <, else wrong neighbour
            pos = nxt; acc += tree[nxt];
        }
    }
    return pos + 1;              // first index pushing the prefix to k
}
\end{lstlisting}
\noindent{\scriptsize Key: descent is the \(\bigO{\log n}\) search down the hidden
binary tree; the test is \(\mathit{acc}+t[\mathit{nxt}]<k\) so we stop just short
and return \(\mathit{pos}+1\) (\(\thmref{c14-lift}\)).}

\begin{lstlisting}[caption={Drill 5 -- Euler subtree add / sum (from lst:c14-subtree)}]
// tin/tout from Ch.5 DFS (tout INCLUSIVE); BIT over the preorder flat array
long long subtree_sum(int v)  { return query(tout[v]) - query(tin[v] - 1); }
void subtree_add(int v, long long d) {
    add(tin[v], d);              // difference BIT on the Euler axis
    add(tout[v] + 1, -d);        // close the interval at tout (inclusive!)
}
long long node_value(int v) { return query(tin[v]); }   // point read after adds
\end{lstlisting}
\noindent{\scriptsize Key: a subtree is exactly the interval \([\tin,\tout]\); the
\(+1/-1\) pair is the difference-BIT range add applied on the flattened axis
(\(\propref{c14-eulerrange}\)). Off-by-one lives entirely in the \(\tout\)
convention.}

\chaptersection{Complete Problem Checklist}

Every problem this chapter offers -- 4 fully solved in-section, worked offline
entries, practice, and one Summit. Tick \emph{only} after an unaided solve; a
hinted attempt counts as half.

\begin{center}
\footnotesize
\renewcommand{\arraystretch}{1.25}
\begin{tabular}{@{}c p{0.19\linewidth} l p{0.20\linewidth} p{0.24\linewidth} c@{}}
\toprule
\textbf{\#} & \textbf{Problem} & \textbf{Source} & \textbf{Key technique} & \textbf{What to look for} & \textbf{Done?} \\
\midrule
\multicolumn{6}{l}{\textbf{Section 18.1 -- BIT core}} \\
W18.1 & Dynamic Range Sum Queries & CSES 1648 & \lv{3} plain BIT, set-as-delta & point \emph{set} \(\to\) \(\delta=x-A[i]\); range \(=\) two prefixes & \(\square\) \\
\midrule
\multicolumn{6}{l}{\textbf{Section 18.2 -- range variants}} \\
W18.2 & Range Update Queries & CSES 1651 & \lv{4} difference BIT & two writes per range; point read \(=\) prefix & \(\square\) \\
W18.2b & Range add + range sum & worked & \lv{4} two-BIT identity & \((x{+}1)Q_1-Q_2\); mind the \(r{+}1\) write & \(\square\) \\
\midrule
\multicolumn{6}{l}{\textbf{Section 18.3 -- order statistics}} \\
W18.3 & Enemy is weak & CF 61E & \lv{5} frequency BIT, two sweeps & per-middle product of (greater-left)\(\times\)(smaller-right) & \(\square\) \\
W18.3b & \(k\)-th live item / Hotel & worked & \lv{4} binary lifting & strict-\(<k\) descent; return \(\mathit{pos}+1\) & \(\square\) \\
\midrule
\multicolumn{6}{l}{\textbf{Section 18.4 -- Fenwick on trees}} \\
W18.4 & Propagating tree & CF 383C & \lv{5} Euler + parity difference BIT & depth-parity sign folded into the delta; inclusive \(\tout\) & \(\square\) \\
\midrule
\multicolumn{6}{l}{\textbf{Section 18.5 -- offline sweeps (worked)}} \\
W18.5a & 2D dominance count & worked & \lv{4} sort one axis, BIT the other & compress the second; count \(\le\), not \(<\) & \(\square\) \\
W18.5b & Nested ranges & CSES 2169 & \lv{4} sort left ends & tie-break order decides contain/contained & \(\square\) \\
W18.5c & PBS over a BIT & worked & \lv{5} lockstep search, BIT predicate & no rollback -- the sweep only appends & \(\square\) \\
\midrule
\multicolumn{6}{l}{\textbf{Summit (chapter end)}} \\
W18 & Irrigation & CF 1181D & \lv{5} water-fill level + \(k\)-th active city & count-space search, then lift the offset-th index & \(\square\) \\
\bottomrule
\end{tabular}
\end{center}

\chaptersection{Recommended Study Order}

Five focused days. Day~1 internalises the walks and the linear build; day~2 the
difference and two-BIT range variants plus compression; day~3 the frequency BIT
and lifting; day~4 Fenwick on the Euler array; day~5 the offline sweeps and the
Summit.

\begin{center}
\small
\renewcommand{\arraystretch}{1.25}
\begin{tabular}{@{}llp{0.55\linewidth}@{}}
\toprule
\textbf{Day} & \textbf{Focus} & \textbf{Tasks} \\
\midrule
1 & Read + do 18.1 & re-trace the query/update walks on \(A\) (\figref{c14-array}) by hand; write \coderef{c14-query}/\coderef{c14-update} cold; CSES~1648 end to end \\
2 & Read + do 18.2 & difference cascade on \(A\); CF... none -- re-derive \coderef{c14-twobit}'s \((x{+}1)Q_1-Q_2\); CSES~1651; one compression pass with \coderef{c14-compress} \\
3 & Read + do 18.3 & lift search for the \(k\)-th on the frequency array (\figref{c14-lift}); right-to-left inversion sweep (\coderef{c14-inv}); CF~61E \\
4 & Read + do 18.4 & overlay \(\tin/\tout\) of \(T\) (\figref{c14-euler}) on the BIT; subtree add/sum by hand; CF~383C parity fold \\
5 & Do 18.5 + Summit & dominance count and nested-ranges sweeps (\(\coderef{c14-dom}\), \(\coderef{c14-nested}\)); PBS skeleton \(\coderef{c14-pbs}\); then CF~1181D; walk the flowchart; first 3 mixed \\
\bottomrule
\end{tabular}
\end{center}

\chaptersection{Difficulty Progression}

Tiers are competency gates -- each row names the \emph{skill} the tier unlocks.

\begin{center}
\small
\renewcommand{\arraystretch}{1.3}
\begin{tabular}{@{}lp{0.34\linewidth}p{0.38\linewidth}@{}}
\toprule
\textbf{Tier} & \textbf{Problems} & \textbf{What passing this tier proves} \\
\midrule
\lv{3}--\lv{4} \ \(\bigstar\bigstar\bigstar\bigstar\)
  & CSES~1648, CSES~1651, worked two-BIT
  & you can build a BIT cold, drive the query/update walks, and pick difference vs two-BIT for the right range semantics \\
\lv{4}--\lv{5} \ \(\bigstar\bigstar\bigstar\bigstar\bigstar\)
  & CF~61E, CF~383C, worked dominance/nested
  & you turn an index into a \emph{meaning} (rank, Euler position, sorted axis) and lift the tree for an order statistic \\
\lv{5}{+} \ \(\bigstar\bigstar\bigstar\bigstar\bigstar\)
  & CF~1181D (Summit), CF~1045G, CF~220E, CF~1000F
  & you compose two BIT levels (search a threshold \emph{and} select the item) and harden for \(n,q\le 5\cdot10^5\), \(k\le 10^{18}\) \\
\bottomrule
\end{tabular}
\end{center}

\chaptersection{Reduction Gallery}

Every entry is ``one BIT flavour + one twist.'' Each line collapses a problem
family onto a listing already in the chapter.

\begin{itemize}[nosep, leftmargin=1.4em]
  \item ``sum a range after point edits'' (CSES~1648) \(\to\) plain BIT, set as
        \(\delta\), difference of two prefixes (\(\coderef{c14-query}\), \(\probref{c14-rangesum}\)).
  \item ``add to a range, read a point'' (CSES~1651) \(\to\) difference BIT, two
        writes per range (\(\defref{c14-diff}\), \(\probref{c14-rangeupd}\)).
  \item ``add to a range, sum a range'' (worked) \(\to\) two-BIT identity
        \((x{+}1)Q_1-Q_2\) (\(\thmref{c14-twobit}\), \(\coderef{c14-twobit}\)).
  \item ``\(k\)-th live item / how many \(<x\)'' (Hotel, worked) \(\to\) frequency
        BIT + binary lifting (\(\thmref{c14-lift}\), \(\coderef{c14-kth}\)).
  \item ``decreasing triplet / smaller-after count'' (CF~61E, LC~315) \(\to\) one
        or two frequency sweeps over ranks (\(\propref{c14-inversion}\), \(\probref{c14-triplets}\)).
  \item ``subtree add/sum, root path'' (CF~383C, CSES~1137) \(\to\) Ch.~5 flatten,
        then difference BIT on the Euler axis (\(\propref{c14-eulerrange}\), \(\probref{c14-prop}\)).
  \item ``how many points dominate \(q\) / nested ranges'' (worked) \(\to\) sort one
        axis, BIT the other (\(\propref{c14-dom}\), \(\coderef{c14-dom}\)).
  \item ``least threshold passing a prefix test, many queries'' (CF~1181D) \(\to\)
        PBS over one shared BIT sweep (\(\propref{c14-pbs}\), \(\probref{c14-summit}\)).
\end{itemize}

\chaptersection{Mixed Problem Set}

Ten entries from CF~1800 to CF~2400, all verified against the live problem set
and absent from the rest of the book. None says ``use a Fenwick tree'' --
recognise the flavour from what the problem needs summed, selected, or swept.

\begin{mixedset}
  \practice{Pashmak and Parmida's problem}{CF 459D}
    {https://codeforces.com/problemset/problem/459/D}{CF 1800 L3}
    {Build \(L[i]\)=\#occurrences of \(a_i\) in \(a[1..i]\) and \(R[i]\)=\# in
     \(a[i..n]\) with two map passes, then count pairs \(i<j\) with \(L[i]>R[j]\):
     that is a strictly-greater-before sweep -- insert \(L\) values into a
     frequency BIT keyed by rank and query the complement. The very
     right-to-left count of \(\coderef{c14-inv}\) on the \(R\) array
     (\(\secref{c14-order}\)).}
  \practice{Closest Equals}{CF 522D}
    {https://codeforces.com/problemset/problem/522/D}{CF 2000 L4}
    {Over queries \([l,r]\), the smallest \(|x-y|\) with \(a_x=a_y\), else \(-1\).
     Unique candidate keys are the \((p,\text{next}[p])\) equal-value pairs; sweep
     \(l\) downward and commit each pair's gap at coordinate \(\text{next}[p]\) into
     a \emph{prefix-min} BIT (append-only \(\Rightarrow\) legal), answer
     \(\min\) over \([l,r]\) by a range read. \(\coderef{c14-min}\) with the
     append caveat from \(\propref{c14-prefixmin}\) (\(\secref{c14-range}\)).}
  \practice{Yaroslav and Divisors}{CF 301D}
    {https://codeforces.com/problemset/problem/301/D}{CF 2200 L5}
    {For each query \([l,r]\), count pairs \((i,j)\), \(l\le i\le j\le r\), with
     \(a_i\mid a_j\) (or vice versa). Pre-list all valid \((i,j)\) pairs, sort them
     and the queries by right end, sweep with a BIT over left ends -- the exact
     nested-ranges count of \(\coderef{c14-nested}\) applied to pair intervals
     (\(\secref{c14-offline}\)).}
  \practice{Radio stations}{CF 762E}
    {https://codeforces.com/problemset/problem/762/E}{CF 2200 L5}
    {Count pairs of stations with \(\min(r_i,r_j)\ge|x_i-x_j|\) and
     \(|f_i-f_j|\le k\). Sort stations by range endpoint descending and activate
     them into a BIT over the (compressed) frequency axis; for a new station the
     admissible partners are a \(2k{+}1\) window of frequencies, read as frequency
     \(\times\) coordinate-prefix counts. The dominance sweep of
     \(\coderef{c14-dom}\) with a \(\coderef{c14-freq}\) inside
     (\(\secref{c14-offline}\)).}
  \practice{The Number Games}{CF 980E}
    {https://codeforces.com/problemset/problem/980/E}{CF 2200 L5}
    {Keep \(n-k\) tree nodes, connected and containing root \(1\), maximising
     \(\sum 2^i\) over kept labels -- value dominance forces a greedy over
     labels \(v=n..1\). Including \(v\) commits its whole path to the already-kept
     structure; a BIT over the Euler order counts how many still-claimable path
     nodes a subtree holds, bounding the pick. Euler-array driving from
     \(\coderef{c14-subtree}\) / \(\coderef{c14-path}\) (\(\secref{c14-trees}\)).}
  \practice{AI robots}{CF 1045G}
    {https://codeforces.com/problemset/problem/1045/G}{CF 2200 L5}
    {Count pairs of robots that see each other (\(|\Delta\text{position}|\) within
     both sight radii) with \(|q_i-q_j|\le K\). Group by IQ and two-pointer over
     the \(K\)-window; within a merge of two IQ groups, sort by position and count
     mutual-visibility pairs with one BIT over the compressed \(x\pm r\) endpoints
     -- a straight 2D dominance \(\coderef{c14-dom}\) after
     \(\coderef{c14-compress}\) (\(\secref{c14-offline}\)).}
  \practice{Fenwick Tree}{CF 1967C}
    {https://codeforces.com/problemset/problem/1967/C}{CF 2300 L5}
    {A game read directly off the \(\mathit{lowbit}\) structure: who wins depends on
     the parity of the number of set bits reached by alternating
     \(i\!\to\!i-\mathit{lowbit}(i)\) descents. Internalise \(\coderef{c14-query}\)'s
     walk and \(\defref{c14-interval}\)'s responsibility intervals until the
     \(\bigO{\log n}\) hop count is reflex (\(\secref{c14-core}\)).}
  \practice{Little Elephant and Inversions}{CF 220E}
    {https://codeforces.com/problemset/problem/220/E}{CF 2400 L5}
    {Count subarrays whose \emph{removal} leaves \(\le k\) inversions. Two pointers
     on the removed window while maintaining the surviving inversion count -- add
     and remove an element's inversion contribution through a frequency BIT over
     ranks. The online count of \(\coderef{c14-inv}\), made sliding, with
     \(\coderef{c14-freq}\) as the live structure (\(\secref{c14-order}\)).}
  \practice{Little Elephant and Tree}{CF 258E}
    {https://codeforces.com/problemset/problem/258/E}{CF 2400 L5}
    {Subtree colour updates and ``how many distinct colours fully cover subtree
     \(v\)'' queries. Flatten with \(\tin/\tout\) (Ch.~5), then a range-add /
     point-read difference BIT per colour tracks coverage depth, and the offline
     query ordering is the sweep of \(\coderef{c14-dom}\) over the Euler axis
     (\(\secref{c14-trees}\), \(\coderef{c14-diff-build}\)).}
  \practice{One Occurrence}{CF 1000F}
    {https://codeforces.com/problemset/problem/1000/F}{CF 2400 L5}
    {Answer queries \([l,r]\): report any value appearing exactly once, else \(0\).
     Sweep right end \(r\) left-to-right, keep each value's last and second-last
     positions in a BIT over positions (contribute \(+1\) at the newest, remove the
     stale ones); read a point in \([l,r]\) for a singleton. Offline nested-range
     bookkeeping as in \(\coderef{c14-nested}\) (\(\secref{c14-offline}\)).}
\end{mixedset}

\chaptersection{The Summit -- CF 1181D Irrigation}

\begin{notebox}[Why this is the chapter top rope]
Sections 18.3 and 18.5 each own one half: 18.3 turns a prefix \emph{count} into a
\(k\)-th \emph{item} by lifting the tree; 18.5 sorts queries and sweeps one axis
with a BIT. The Summit is both at once, twice: it must first \emph{find a water
level} by searching a monotone count-space prefix function, then \emph{select the}
\(\mathit{offset}\)-th still-short city -- a \(k\)-th lift over a live frequency
BIT -- and it does so for \(q\le 5\cdot10^5\) queries with year numbers up to
\(10^{18}\), so nothing can be simulated year by year. Compress the counts, sort
queries offline, and let the two BIT walks answer each in \(\bigO{\log n}\).
\end{notebox}

\solvedproblem{c14-summit}{Irrigation}{Codeforces 1181D}
  {https://codeforces.com/problemset/problem/1181/D}{CF 2200 L5}
  {water-fill over host counts + offline sorted queries + frequency-BIT k-th lifting over city index}

\begin{problemstatement}[CF 1181D Irrigation]
There are \(m\) cities, numbered \(1..m\). The first \(n\) olympiads are given:
year \(i\) was hosted by city \(a_i\). From year \(n+1\) on, the host is chosen by
a rule: \emph{the city that has hosted the fewest times so far; ties go to the
smallest index}. Each of \(q\) queries gives a year \(k_i\) with
\(n+1\le k_i\le 10^{18}\); report the host city that year.

\textbf{Input.} \(n,m,q\) (\(1\le n,m,q\le 5\cdot10^5\)); then \(n\) integers
\(a_1..a_n\) (\(1\le a_i\le m\)); then \(q\) lines, each \(k_i\)
(\(n+1\le k_i\le 10^{18}\)).

\textbf{Output.} \(q\) integers -- the host city in each queried year.

\textbf{Sample 1.} \(n=6,\ m=4,\ q=10\), \(a=(3,1,1,1,2,2)\),
\(k=7,8,\dots,16\) \(\to\) \texttt{4 3 4 2 3 4 1 2 3 4}.\\
\textbf{Sample 2.} \(n=4,\ m=5,\ q=4\), \(a=(4,4,5,1)\),
\(k=15,9,13,6\) \(\to\) \texttt{5 3 3 3}.
\end{problemstatement}

\solstep{Key Observations}
\begin{enumerate}[nosep,leftmargin=1.6em]
  \item Let \(c[x]\) be how many of the first \(n\) years city \(x\) hosted.
        The rule always lifts a current minimum, so it \emph{equalises} the
        counts: define \(S(L)=\sum_x \max(0,\,L-c[x])\), the number of post-rule
        selections needed to bring \emph{every} city to at least \(L\). \(S\) is
        non-decreasing in \(L\) -- a monotone prefix function.
  \item Write a query as \(j=k-n\), the index of the selection after the rule
        starts (1-based). The chosen city sits in the unique \emph{band}
        \(S(L-1)<j\le S(L)\): precisely the cities with \(c[x]<L\), all standing
        at \(L-1\), picked in increasing index order. So the answer is the
        \(\big(\,j-S(L-1)\,\big)\)-th smallest index among cities with \(c[x]<L\).
  \item Finding \(L\) is a monotone search over a count axis -- \(18.5\)'s PBS /
        binary-search trick, evaluating \(S\) with prefix sums of a
        \emph{frequency-of-count} BIT. Picking the offset-th active city is
        \(18.3\)'s \(k\)-th lift (\(\thmref{c14-lift}\), \(\coderef{c14-kth}\)).
\end{enumerate}

\solstep{Reduction to a Lift over a Sweep}
Two BITs, two axes. \(\mathcal{B}_{\mathit{cnt}}\): a frequency BIT over count
values (how many cities host exactly \(h\) times) with running sums of
\(\#\text{cities}\) and \(\sum\text{counts}\), so \(S(L)=L\cdot C(L)-P(L)\) is a
prefix read. \(\mathcal{B}_{\mathit{idx}}\): a frequency BIT over city indices,
city \(x\) active iff \(c[x]<L\); its \((\,j-S(L-1)\,)\)-th active slot is found
by \(\coderef{c14-kth}\)'s descent. Process queries sorted by \(k\): \(L\) is
non-decreasing across them, so \(\mathcal{B}_{\mathit{idx}}\) only ever
\emph{appends} activations -- exactly the append-only regime that lets
\(\propref{c14-pbs}\)'s PBS skip a rollback. Each query: advance \(L\) while
\(S(L)<j\), activate newly-qualifying cities, then lift the
\((j-S(L-1))\)-th active index. \(\bigO{(n+q)\log n}\) after a \(\bigO{n}\) count.

\solstep{Trace on Sample 1}
\begin{tracebox}[Water-fill on \(n=6,m=4\), \(a=(3,1,1,1,2,2)\)]
Counts \(c=(3,2,1,0)\) for cities \(1..4\). \(S(L)=\sum_x\max(0,L-c[x])\):
\[
{\small
\begin{array}{c|ccccc}
L & 0 & 1 & 2 & 3 & 4\\ \hline
S(L) & 0 & 1 & 3 & 6 & 10
\end{array}}
\qquad
{\small
\begin{array}{c|c|c|c|c|c}
k & j{=}k{-}6 & \text{band }L & \text{batch }(\text{cities }c{<}L) & \text{offset} & \text{host}\\ \hline
7&1&1&\{4\}&1&4\\
8&2&2&\{3,4\}&1&3\\
9&3&2&\{3,4\}&2&4\\
10&4&3&\{2,3,4\}&1&2\\
11&5&3&\{2,3,4\}&2&3\\
12&6&3&\{2,3,4\}&3&4\\
13&7&4&\{1,2,3,4\}&1&1\\
14&8&4&\{1,2,3,4\}&2&2\\
15&9&4&\{1,2,3,4\}&3&3\\
16&10&4&\{1,2,3,4\}&4&4
\end{array}}
\]
Read the batch column as \(\mathcal{B}_{\mathit{idx}}\)'s active set (indices
ascending) and the offset as the lift target; the answers
\(4,3,4,2,3,4,1,2,3,4\) match the official output. The jump from \(j=6\) (offset
\(3\) at level \(3\)) to \(j=7\) (level \(4\), offset resets to \(1\)) is where
city~1 first joins because \(S(3)=6<7\le S(4)\).
\end{tracebox}

\solstep{Implementation}
\begin{lstlisting}[caption={CF 1181D: sort counts, search the water level, lift the offset-th active city}, label={lst:c14-summit}]
#include <bits/stdc++.h>
using namespace std;
int bit_[500005];                        // frequency BIT over city indices
int M, HB;                               // HB = 2^floor(log2 M)
void add(int i,int d){for(;i<=M;i+=i&-i)bit_[i]+=d;}
int kth(long long k){                    // least idx whose prefix >= k
    int pos=0; long long acc=0;
    for(int pw=HB;pw>0;pw>>=1){int nx=pos+pw;
        if(nx<=M && acc+bit_[nx]<k){pos=nx;acc+=bit_[nx];}}
    return pos+1;}
int main(){
    int n,q; scanf("%d%d%d",&n,&M,&q);
    vector<int> cnt(M+1,0);
    for(int i=0,a;i<n;i++){scanf("%d",&a);cnt[a]++;}
    vector<pair<int,int>> ord;           // (count, city), sorted count asc
    for(int x=1;x<=M;x++) ord.push_back({cnt[x],x});
    sort(ord.begin(),ord.end());
    vector<int> cs(M); vector<long long> ps(M+1,0);
    for(int t=0;t<M;t++) cs[t]=ord[t].first;
    for(int t=0;t<M;t++) ps[t+1]=ps[t]+cs[t];      // prefix sum of counts
    auto S=[&](long long L)->__int128{             // selections to reach level L
        long long k=lower_bound(cs.begin(),cs.end(),L)-cs.begin(); // #{cnt<L}
        return (__int128)L*k-ps[k];};
    HB=1; while(HB*2<=M)HB*=2;
    vector<pair<long long,int>> qs(q);
    for(int i=0;i<q;i++){long long k;scanf("%lld",&k);qs[i]={k-n,i};}
    sort(qs.begin(),qs.end());            // by j ascending => L monotone
    vector<int> ans(q); int act=0;        // act: cities with cnt < L activated
    for(auto&qu:qs){
        long long j=qu.first;             // 1-based selection index
        long long lo=1,hi=(long long)2e9;
        while(lo<hi){long long mid=(lo+hi)/2;if(S(mid)>=j)hi=mid;else lo=mid+1;}
        long long L=lo;                   // smallest level with S(L) >= j
        while(act<M && ord[act].first<L){ add(ord[act].second,1); act++; } // append-only
        long long used=(long long)(S(L-1)); // = S(L-1) fits in int64 (<=j)
        long long off=j-used;             // 1-based inside the band
        ans[qu.second]=kth(off);          // off-th smallest active city index
    }
    for(int i=0;i<q;i++) printf("%d ",ans[i]);
}
// Two BIT axes: S(L) is a count-space prefix read (sorted counts + prefix sums);
// kth(off) is 18.3's index-space lift.  Sorting queries by j makes L non-decreasing,
// so activations are appended once -- no rollback, exactly prop:c14-pbs's argument.
// __int128 guards L*k on the level search; j,k <= 1e18 need long long.
\end{lstlisting}

\begin{hardeningbox}[AC \(\longrightarrow\) red: hardening the Summit]
\begin{itemize}[nosep,leftmargin=1.4em]
  \item \textbf{Never simulate years.} \(k_i\le 10^{18}\); the only tractable route
        is the closed-form band \(S(L-1)<j\le S(L)\) with \(j=k-n\). A per-year
        loop is the WA/TLE reflex to kill first.
  \item \textbf{Two axes, one idea.} The level search lives in \emph{count} space
        (\(\mathcal{B}_{\mathit{cnt}}\)); the city pick lives in \emph{index} space
        (\(\mathcal{B}_{\mathit{idx}}\), \(\coderef{c14-kth}\)). Conflating them is
        the classic off-by-one that shifts every answer by one city.
  \item \textbf{Activation is append-only.} Because \(L\) never decreases across
        sorted queries, cities enter \(\mathcal{B}_{\mathit{idx}}\) once, in count
        order. Do \emph{not} rebuild the BIT per query -- hoist it, exactly the
        no-rollback argument of \(\propref{c14-pbs}\).
  \item \textbf{Tie direction.} Within a band, hosts are chosen by \emph{smallest
        index}; the sort of \(\mathit{ord}\) is \((\text{count},\text{city})\)
        ascending, and \(\mathcal{B}_{\mathit{idx}}\) indexes by city so the lift
        already returns the smallest. Sorting by city alone breaks the level.
  \item \textbf{\(\mathit{lowbit}\) and \(k\) are \code{long long}.} \(j\) and the
        prefix target exceed \(2^{31}\); the offset into the band still fits
        \(m\le 5\cdot10^5\), but \(S(L)\) and the running \(\mathit{acc}\) do not.
  \item \textbf{I/O.} \(n,q\le 5\cdot10^5\): \code{scanf}/fast I/O; sorting the
        queries by \(k\) (not reading order) is what makes the sweep monotone.
\end{itemize}
\end{hardeningbox}

\begin{recallbox}[Chapter 18 in ten seconds]
A Fenwick tree is an array hiding a binary tree on the index axis: \(\mathit{lowbit}(i)=i\&(-i)\)
defines slot \(i\)'s responsibility interval, a down-walk reads a prefix, an
up-walk writes a point (\(\defref{c14-interval}\), \(\thmref{c14-query}\)), and a
\(\bigO{n}\) build hands blocks forward (\(\propref{c14-linear}\)). Change the
\emph{stored quantity} and you get a family: differences for range-add/point-read
(\(\defref{c14-diff}\)), two BITs for range-add/range-sum (\(\thmref{c14-twobit}\)),
frequencies for order statistics lifted down the tree (\(\thmref{c14-lift}\)), the
Euler array for subtrees and paths (\(\propref{c14-eulerrange}\)). Change the
\emph{order of operations} and you get offline sweeps: sort one axis, BIT the
other (\(\propref{c14-dom}\)), wrap it in parallel binary search for a threshold
(\(\propref{c14-pbs}\)) -- a BIT only appends, so no rollback.
\selftest{
  \item Give the two walks and the loop-variable sign on each. \hfill\lv{3}
  \item Why is a min-BIT append-only, and what does that force on decrements? \hfill\lv{4}
  \item In the \(k\)-th lift, is the descent test \(<k\) or \(\le k\), and what do you return? \hfill\lv{4}
  \item State the two BIT axes the Summit uses and why one needs no rollback. \hfill\lv{5}
}
\end{recallbox}

\chaptersection{Tool Kit Inventory}

The machine assembled through Chapter~18. Earlier chapters' rows are carried
forward; this chapter's are appended. Pointers only.

\begin{center}
\small
\renewcommand{\arraystretch}{1.25}
\begin{tabular}{@{}p{0.21\linewidth} l c p{0.25\linewidth} p{0.17\linewidth}@{}}
\toprule
\textbf{Tool} & \textbf{Since} & \textbf{Lv} & \textbf{Trigger phrase} & \textbf{Where} \\
\midrule
\multicolumn{5}{l}{\itshape Carried forward} \\
Euler tour \(\tin/\tout\) flatten & 5 & \lv{3} & ``a subtree is a contiguous range'' & \chref{c24} (\(\thmref{c24-range}\)) \\
Merge-sort inversion count & 1 & \lv{4} & ``offline divide-and-conquer twin'' & \chref{c01} \\
Rollback-DSU PBS scheduling & 17 & \lv{5} & ``binary-search all queries in lockstep'' & \secref{c28-offline} \\
\midrule
\multicolumn{5}{l}{\itshape Chapter 18 -- new this chapter} \\
Plain BIT (query/update/build) & 18.1 & \lv{2} & ``live prefix sum, point edit'' & \coderef{c14-query}, \coderef{c14-build} \\
Difference BIT & 18.2 & \lv{3} & ``add to a range, read a point'' & \coderef{c14-diff-build} \\
Two-BIT range sum & 18.2 & \lv{4} & ``add to a range, sum a range'' & \coderef{c14-twobit} \\
Coordinate compression & 18.2 & \lv{3} & ``values sparse / up to 1e9'' & \coderef{c14-compress} \\
Frequency BIT + \(k\)-th lift & 18.3 & \lv{4} & ``select the \(k\)-th live item'' & \coderef{c14-kth} \\
Inversion / smaller-after sweep & 18.3 & \lv{4} & ``count inversions online'' & \coderef{c14-inv} \\
Euler-driven BIT (subtree/path) & 18.4 & \lv{4} & ``subtree add/sum on a tree'' & \coderef{c14-subtree}, \coderef{c14-path} \\
2D dominance / nested sweep & 18.5 & \lv{4} & ``sort one axis, BIT the other'' & \coderef{c14-dom}, \coderef{c14-nested} \\
PBS over a BIT & 18.5 & \lv{5} & ``least threshold, prefix predicate'' & \coderef{c14-pbs} \\
Count-space + index-space lift & Summit & \lv{5} & ``find a level, pick the offset-th'' & \coderef{c14-summit} \\
\bottomrule
\end{tabular}
\end{center}

\chaptersection{Related Topics}

\begin{itemize}[nosep, leftmargin=1.4em]
  \item \chref{c24} (Ch.~5, Euler Tour Technique) --- backward: owns the
        \(\tin/\tout\) flattening and the ``subtree \(=\) interval'' fact
        (\(\thmref{c24-range}\)); 18.4 reuses them by label and never re-derives
        them -- this chapter supplies the \emph{range engine} Ch.~5 promised, and
        its CSES~1137 practice is re-solved here with a BIT (\(\probref{c14-prop}\)).
  \item \chref{c28} (Ch.~17, DSU \& Sack) --- backward/forward: the parallel
        binary-search \emph{scheduling} is theirs (\(\secref{c28-offline}\)); 18.5
        discharges Ch.~17's forward pointer by swapping the rollback-DSU predicate
        for a BIT sweep (\(\propref{c14-pbs}\)). The Summit is the joint
        composition the handoff advertised.
  \item \chref{c01} (Ch.~1, complexity \& counting) --- backward: owns the
        merge-sort inversion count; 18.3 is the online frequency-BIT twin
        (\(\propref{c14-inversion}\)) -- same \(\bigO{n\log n}\), different mechanics.
  \item \chref{c06} (Ch.~6, Constructing Trees) --- back-cite: owns CF~1354D
        Multiset and PBDS order-statistics trees; 18.3 is the array-index rival
        (prose only, URLs not re-linked).
  \item \chref{c13} (Ch.~19, Segment Trees) --- forward: the BIT-vs-segtree
        boundary is drawn in 18.2 -- a BIT wins on prefix sums and simplicity, a
        segtree owns arbitrary-range min/non-invertible combines and lazy; no
        segtree code appears here.
  \item \chref{c15} (Ch.~20, Lazy Propagation) --- forward: range-add/range-sum
        with a \emph{third} operation (or non-commutative combine) belongs there,
        not to the two-BIT trick (\(\thmref{c14-twobit}\)'s remark).
\end{itemize}

\begin{lastword}[Last Word --- Chapter 18]
\emph{``One \(\mathit{lowbit}\), two walks, and an array starts answering prefix
questions in \(\bigO{\log n}\). Change what you store -- sums, differences,
frequencies, Euler slots -- and the same two walks become a whole family; change
the order you touch them -- sort one axis, sweep the other -- and offline
questions fall to the same structure.''}

\medskip
The plain BIT is honest and near-free (\(\coderef{c14-query}\)); the \(\bigO{n}\)
build makes it production-cheap (\(\propref{c14-linear}\)). Store differences and
a range add is two writes (\(\defref{c14-diff}\)); store a position-weighted pair
and range sums follow (\(\thmref{c14-twobit}\)). Index a frequency array and the
down-walk becomes a monotone search, so lifting the tree selects the \(k\)-th live
item (\(\thmref{c14-lift}\)) -- the reflex the Summit spends twice, once on a count
axis and once on a city axis (\(\probref{c14-summit}\)). Drive Ch.~5's flat array
and subtrees/paths are just ranges again (\(\propref{c14-eulerrange}\)); sort one
axis and sweep the other and dominance, nested ranges, and a threshold search are
one shared BIT with no undo (\(\propref{c14-dom}\), \(\propref{c14-pbs}\)).

\medskip
Carry three reflexes forward:
\begin{enumerate}[nosep, leftmargin=1.6em]
  \item \emph{Index \(=\) meaning.} A BIT slot can be an array position, a rank, an
        Euler entry, or a count value -- naming which is half the solution
        (\(\figref{c14-taxonomy}\)).
  \item \emph{Append-only is the BIT's contract.} Prefix sums and monotone sweeps
        fit; arbitrary decrements and non-invertible combines are the segment
        tree's job (\(\propref{c14-prefixmin}\), forward to \chref{c13}).
  \item \emph{Offline beats overkill.} If you may sort, a single sorted sweep plus
        a plain BIT replaces rollback stacks and persistent structures
        (\(\secref{c14-offline}\)).
\end{enumerate}

\medskip
And three habits, drilled until boring: keep the query and update walks'
loop-signs straight and 1-indexed; re-check the \(\tout\) inclusivity and the
\(r{+}1\) write against the smallest sample that could expose them; and before
reaching for a segment tree, ask whether the combine is a prefix sum -- if so this
chapter's eight lines already solve it.

\medskip
\textbf{Next:} \chref{c13} (Segment Trees) takes the range engine you now own and
generalises the combine -- arbitrary intervals, min/max with deletions, lazy
propagation -- paying a constant factor for the extra power. The BIT-vs-segtree
line drawn in \secref{c14-range} is exactly where that chapter picks up.
\end{lastword}


\closechapter
